\documentclass[12pt,a4paper]{article}

\usepackage{mathtools}
\usepackage{mathrsfs}
\usepackage{float}
\usepackage{caption}
\usepackage{tikz}
\usepackage{tikz-cd}
\usetikzlibrary{arrows,positioning,calc,decorations.pathmorphing}

%% Current project v3.070 source-reference bindings; previous macro bodies unchanged.
% Current II-D reference root for PapersMinimal v3.063; mathematical macros otherwise unchanged.
% Macros-v2_188: re-based current paper identities and A--G family; no change to geometry.
% Macros-v2_184: active Paper III-E v1.200/III-F v1.267 binding for framework v3.031; mathematical macros unchanged.
% Macros-v2_183: updated shared cross-reference bindings for Paper1B/III-E/III-F v3.030; mathematical macros unchanged.
% Macros-v2_182: active Paper II-F/V-A cross-references for framework v3.029; definitions unchanged.
% Paper-series v2.381: new I-B torsion/modular results, bibliography and current source roots.
% Paper-series v2.380: new exceptional bibliography and synchronized paper references.
% Paper-series v2.379 current 31-paper macro/aux registry.
% Paper-series v2.378: current 31-paper AUX/bibliography registry updated for Paper I-B v1.2 literature additions.
% Framework v2.371: current-source reference registry refreshed; old theorem labels retained.
% Framework v2.370: current 30-paper source/AUX root rebinding; native mathematical macros unchanged.
% Framework v2.368: synchronized merged 30-root registry; mathematical notation unchanged.
% Framework v2.358: active 30-paper source registry synchronized; existing mathematical macros unchanged.
% Framework v2.357: selective new source-owned certificates; active reference roots updated for five revised papers.
% Framework v2.356: active 30-root registry synchronized to the 30 papers in this release.
% Framework v2.275: exact positive dense-twin word cubature, finite repeatability and eighth-moment classification, nested sampling frames, joint action and POVM-jet exports; preserves v2.274.
% Framework v2.274: filename and root-source synchronization.
% Framework v2.273: restored missing Problem 3 theorems from v2.268; exact exceptional frames, twin geometry, convolution repeatability and projector/occurrence firewalls; merged on v2.272 without overwriting independent branches.
% Framework v2.272: merged v2.270 covariant native-readout mathematics with v2.271 Problem 4 finite-action Hopf hypercharge/holonomy corrections; no universal bulk-gauge uniqueness claim.
% Independent branches: v2.270 readout/guard exports and v2.270 K16 action-consumed charge closure; all prior scoped assumptions retained.
% Framework v2.269: synchronized 30 active AUX roots; mathematics and bibliography preserved from Macros-v2_161.
% Framework v2.266: Paper II-C converse-derived outcome resources; II-E, III-C, IV-A/B, V-A and VI-B downstream mathematical closures; 30-root registry synchronized.
% Framework v2.265: II-C coordinate-stage constitution and exact first-symbol host inheritance; prior mathematical domains retained.
% Framework v2.264: paper II-B reverse/audit and IV-A observable-minimal exchange; synchronized external references to 30 active roots.
% Framework v2.263: Paper I-C adds bridge-controlled root-carry parity and source-relative old-cap cancellation; other root mathematical bodies retained; active 30-root references synchronized.
% Framework v2.262: integrates Paper II-F radius, Casimir and readout results; active 30-root references synchronized.
% Framework v2.261: merged Paper II-B GP2-level/geometric repairs; preserves v2.260 II-D repairs; active references synchronized.
% Framework v2.260: Paper II-D signed RS, finite-symbol typing, and conditional curvature/Jordan repairs; references synchronized.
% Framework v2.259: conditional native SU(2) atlas and explicit later discharge; v2.258 mathematical bodies retained elsewhere.
% Framework v2.258: preserves parallel II-A/dynamics update and adds complete readout/frequency closures.
% Framework v2.257: exact minimal complete position readout and arbitrary-preparation frequency law.
% Framework v2.256: exact native inertia and sharp bounded ordered-source closure; preserves concurrent v2.255 work.
% Framework v2.255: preserves parallel detector integration and adds exact finite readout closures.
% Framework v2.254: Pasterski bibliography and synchronized 30-paper external-reference registry.
% Framework v2.253: typed raw coordinate defect in I-A and synchronized VI-A import; active references synchronized.
% Framework v2.252: synchronized source registry; optional III-G shell-frame calculation.
% Framework v2.251: forward common-action execution, joint source jets and scoped coupled control.
% Framework v2.251: exact full-heat variations, fermion source jets and explicit physical completion frontier.
% Framework v2.251: certified interacting source jets, pole stability and scoped completion proofs.
% Framework v2.251: synchronized active source registry; mathematical notation unchanged.
% Framework v2.251: synchronized active source registry; mathematical notation unchanged.
% Framework v2.251: synchronized active source registry; mathematical notation unchanged.
% Framework v2.251: synchronized source registry; no mathematical notation changes.
% Framework v2.244: synchronized 30-paper AUX registry and rooted phase-source action-order tests.
% Framework v2.242: signed rooted phase lift, complete transverse return and exact family cubic.
% Framework v2.241: proved rooted conormal geometry, conditional cap realization and scoped action tests.
% Framework v2.239: all 30 active source roots synchronized; no mathematical macro changes.
% Framework v2.238: Paper IV A/B/C split and 30-root reference registry.
% Framework v2.237: Paper VII-D technical split and 28-root reference registry.
% Framework v2.236: Paper VI-C technical-volume reference added; all current root bindings advanced.
% Framework v2.235: 26-root AUX and bibliography bindings synchronized; mathematical macros unchanged.
% BEGIN INLINE SOURCE: Paper00-TeXBuildCompatibility-v1_0.tex
% Rendering-only repairs for inherited long appendices and bitmap text fonts.
% Integer counters through 26 retain the original labels; later values use AA,
% AB, ... . This file changes no action coefficient, state or theorem premise.
\RequirePackage{alphalph}
\makeatletter
\let\BHPrimitiveAlph\@Alph
\newalphalph{\BHExtendedAlph}{\BHPrimitiveAlph}{26}
\let\@Alph\BHExtendedAlph
\AtBeginDocument{%
  \@ifpackageloaded{microtype}{\microtypesetup{expansion=false}}{}%
}
\makeatother

% END INLINE SOURCE: Paper00-TeXBuildCompatibility-v1_0.tex

% Framework v2.164: current reference/bibliography profile; mathematical definitions retained.
% Framework v2.162: finite metric, functional Hilbert and scoped Ward integration.
% Framework v2.161: inherits Macros-v2_57.tex exactly except the current external-reference binding.
% Macros v2.57 / framework v2.156: current reference registry only; mathematical definitions unchanged.
% Framework v2.143 / Macros v2.55: current external-reference bindings; mathematical notation unchanged.
% Framework v2.110 / Macros v2.54: current external-reference registry and bibliography; mathematical notation inherited.
% Framework v2.67 / Macros-v2_53: adds Born-action-cell and joint-EH readout symbols; absolute unit realization is metrological; MATCH_{J/scal}^{EH} remains the sole G24 normalization bridge.
% Framework v2.35 / Macros-v2_52: grouped paper numbering, active AUX registry, and Born-v6_20 bibliography binding synchronized; notation unchanged.
% Framework v2.26: active-source AUX and bibliography bindings synchronized; notation unchanged.
% Framework v2.09: reviewed early-reciprocity integration; scoped physical imports.
% Framework v2.08: Born--Euler mathematics and conditional complete-action potential criterion.
% Framework v2.07: active AUX registry includes Paper I-B v1.2; shared bibliography Born-v6_18.
% Macros-v2_46 / framework v2.04: mathematical atlas and downstream realization correction.
% Macros-v2_46 / framework v2.03: scoped structure integration and current dependency bindings.
% Macros-v2_46 / framework v2.01: shared source bindings; prior revision comments are historical.
% v2.42 / framework v1.99: Synchronized RG/AS owner and bibliography bindings.
% v2.41 / framework v1.99: active manuscript AUX bindings and geometric-closure bibliography synchronized.
% Release-v1.80 centralized source synchronization for scalar-circle/comparison-band integration.
% Release v1.79: Hurwitz hierarchy and kinetic-normalized action response integration; physical mass matching remains conditional.
% Release v1.78: approved M8 closure-batch integration; physical and analytic domains retained.
% Current v1.78 registry: all active TeX roots are synchronized below.
% Earlier release comments are historical scope records, not current status.
% Release-v1.77 synchronization: native same-owner Noether closure, generator--flow/history typing, analytic-certification firewalls, scalaron/K16, LF/QCD, quantum-Ward, and RAR boundary-owner transport are registered; shared mathematical notation is unchanged.
% Release-v1.76 synchronization: registry points to the v1.76 active paper roots.
% Release-v1.76 synchronization: registers the local A+Phi+g Noether/constraint closure and dynamics-frontier propagation; shared mathematical notation is unchanged.
% Release-v1.75 synchronization: registry points to the v1.75 active paper roots.
% Release-v1.75 synchronization: integrates the joint-Dirac Einstein bulk real form and the local odd-scalar heat-coefficient action on explicit analytic first-jet/interior domains; v1.74 Yang--Mills closure is retained. Fermionic reality, pure-metric curvature-squared/exact-nonlocal terms, noncanonical boundaries, physical vacuum selection, and actual solution-history tokens remain open.
% Release-v1.74 synchronization: registry points to the v1.74 active paper roots and retains the v1.73 theorem labels.
% Release-v1.74 synchronization: canonical Macros-v2_34.tex; generated local nonlinear Yang--Mills action, constrained connection-field phase space, source-free Gauss momentum map, and Legendre closure are integrated on their proved domains. The remaining non-gauge ACT_cl sectors and actual solution-history selection remain open.
% Release-v1.71 synchronization: closes the former JOINT_x/JOINT_5
% connection-occurrence gap by canonical structure-group extension of the
% nonflat Born connection, faithful associated representation, and the local
% cotangent--Clifford operator constructor.  Actual fluctuation members,
% actions, vacua, solutions, calibrations, and analytic QFT domains remain
% separately typed.  The active-AUX registry is advanced once.
% Release-v1.71 canonical macro and active-AUX registry.
% ======================================================================
% Macros-v2_39.tex -- canonical publication macro file
% Release v1.71 retains the v1.70 null-ray/QCD boundaries and additionally
% removes the duplicate structural connection actuality premise.  It does not
% claim global bundle topology, select a gauge/Higgs member, or promote the
% finite construction to an interacting continuum theory.
% Release v1.68 constructs the finite light-front total-fraction operators and
% parton number/momentum operator measures, proves their simplex and Hausdorff-limit
% theorems, and separates positive cut measures from scheme-renormalized
% light-ray distributions.  It also separates finite real Schur poles, stable
% norm-resolvent/Riesz pole-residue limits, and continued resonance
% pole/residue limits.  The Paper-VII-C scattering package is claim-indexed:
% local QFT and spectral claims import no S-matrix or infrared premise, while
% asymptotic amplitudes, LSZ, S-matrix poles, completeness, and cross sections
% import only their named massive or massless/long-range clauses.
% Release v1.67 factorizes BND48_H into a same-action Hessian/Schur
% descendant and still-open typed boundary-trace, critical-lift, and
% refinement-limit conditions.  It proves the finite-collar Dirichlet-to-
% Neumann law sqrt(y)tanh(L sqrt(y)), its nonuniform critical limit, and the
% finite-return ultra-infrared turnover.  The Paper-I-A Hardy constant mode is
% an exact geometric near-hit, not yet the physical metric boundary mode.
% The current local f(R) Hessian remains massive and therefore still refutes
% a RAR prediction on the minimal action branch.
% Release v1.66 derives the transverse depth-six normalizer bridge from the
% already generated C6 Reynolds image and Paper-I-A binary-dihedral normalizer.
% The bridge is the unique least 2D3-stable scalar-containing realization of
% the faithful quotient 2D3/C6=C2, has trace defect six and hence K=48, but is
% transverse to the nested active carrier.  Exact K=44 and K=52 contraction,
% semigroup, and energy-midpoint identities are registered as structural
% readouts rather than projector events.  The active E48 carrier gap remains
% unchanged.  A same-action boundary-symbol/Calderon intertwiner is isolated
% as the necessary and still unproved owner bridge to the static DtN sort;
% representation equivalence alone neither derives RAR nor creates a K=48
% action junction, pole, occupation, or cosmological transition.
% Release v1.65 derives the native E_48-centered carrier mass gap from the
% GP1--GP3-generated active spectrum, distinguishes its carrier-free interval
% from a particle located at E_48, and withdraws the unsupported promotion to
% a positive continuum Yang--Mills mass gap.  The latter remains a valid
% conditional transfer theorem on an enlarged continuum branch, but it is not
% an active framework prediction or an unfinished native-gap proof.
% Release v1.64 synchronizes Derivation Batches 14--22.  It registers the
% GP1--GP3 heat-profile orbit, same-Hessian electroweak response and pole
% ownership, K16/K48 action-response corrections, perturbation state and
% non-Gaussianity typing, finite-vacuum/gap/QCD limit repairs, factorized
% position-area ownership, the actual static metric response, and the ordered
% Wilson/string-tension limits.  Strict continuum convergence, positive
% physical gap, surviving higher cumulants, and positive tension remain open.
% Release v1.61 closes TIMEWIDTH16 as an independent direct-production datum.
% JUNC_{2|7} continues to own the finite intrinsic microscopic profile/width,
% while the actual same-action background/production history determines its
% physical-time pullback, support, and duration.  The current Paper-I-A--VII-C AUX
% registry is advanced once; shared mathematical notation is unchanged.
% Release v1.60 derives the final-octet primitive-remnant depletion, canonical
% Hopf-current representative, exact Hopf descent to dOmega/(4pi), and
% closed-FLRW a^{-3} neutral-current shape.  It retypes K=48 as the onset
% of the final one-unit geometric transfer and records cDM/RAR as distinct
% physical projections of that common neutral parent while preserving
% occurrence, Hilbert-owner, stress-matching, and static-kernel premises.
% Release v1.59 transports the v1.58 exceptional-shell/K16 closure through
% Papers III-F and VI-B--VII-C.  It replaces Paper VI-B's obsolete mixed-Hessian
% output typing by the third-variation pair vertex and Bogoliubov occupation
% support, derives the post-seed Higgs-pair feedback vertex and its exact
% alpha--delta kernel factors, propagates the unique K16/K48 action-junction
% split, identifies the QCD delta count with the upstream 2O shell atom only
% as a common-scalar numerical compatibility, and narrows THERM16 to dynamical
% approach/beta selection because the finite Gibbs/KMS target is already
% unique on a supplied simple observable algebra once its Hamiltonian and
% inverse temperature are supplied; direct-sum central weights remain typed.
% Release v1.58 registers the exceptional nearest-shell closure, the derived
% transverse K=16 action event, the unique interior action junction, the
% type-correct ordered pair vertex, the thin-layer squeeze law, and the
% radiation-like/matter-like kinematic split.  REG_0, transfer magnitude,
% physical layer width, and thermalization remain explicitly conditional.
% Release v1.57 registers the active-depth-six obstruction, the native
% K=48 carrier desert, the typed 7|9 owner reduction, and the dynamical
% reheating handoff across Papers III-B--VI-B.
% Release v1.56 registers the least-support n_2 inflation dressing of the
% persistent n_7 Laurent response, the analytic global-convexity and positive
% endpoint-measure proofs, the derived inflation core and K=16 coefficient
% freeze, and the no-particle dark-sector split into static RAR plus the
% conditional homogeneous effective-curvature response at reflected K=48.
% There is no n_6 action, scalaron-DM mass, or second Einstein--Hilbert term.
% Release v1.55 registers the protected-octet resolved/residual reflection
% q_sp+Z_W=8, its j<->8-j and K<->64-K scale covariance, and the derived
% rank-one-exit quadratic grade C_1^2/U=C_2 in Papers III-C, III-E, and VIII.
% The reflection is a typed trace/scale theorem, not a time reversal or a
% physical event-occurrence theorem; reheating transfer remains conditional.
% Release v1.54 registers the exact positive-part Laurent indices, the
% active-domain/terminal-endpoint action split, the no-second-EH firewall, and
% the q_DM=6/K_DM=48 scalaron readout typing across Papers III-D/E/F, VI,
% VIII, and IX.
% Release v1.53 registers the common nonnegative Laurent obstacle projection
% for n_1^+ and n_7^+, discharges the three former exit/freeze assumptions,
% and advances the Paper-III-D/E/F, VIII, and IX consumers.  Shared notation
% and bibliography ownership are unchanged.
% Release v1.51 registers the PQTI derivation batch and its updated
% Paper-III-A/C/D/E/F, VIII, and IX consumers.  Shared notation and
% bibliography ownership are unchanged.
% Release v1.50 advances the Paper-III-B AUX target to v1.40 after the
% rigorous audit repair; shared notation and bibliography ownership are unchanged.
% Release v1.49 advances the Paper-II-E/F AUX target to v3.14 after the
% minimal audit repair; shared notation and bibliography ownership are unchanged.
% Release v1.46 synchronizes Paper II-A's minimal frozen constructor-origin
% theorem and its five dependent AUX targets.  Shared notation is unchanged.
% Release v1.45 updates the Paper-II-B AUX target to v2.23 after the
% GP3-rigidity and appendix-placement repair; shared notation is unchanged.
% Release v1.44 advances the centralized Paper-I-A--VII-C AUX registry after the
% Paper-I-A v5.16 protected-completion scope repair.  Every current
% paper imports this one file; no paper-level opt-in or \externaldocument list
% is permitted.  The bibliography remains centralized here through Born-v6_13.
% Shared bibliography, theorem infrastructure, colors, notation alignment,
% and current Paper-I-A--VII-C external-reference registration.
% This file centralizes common imports, bibliography sorting, colors, theorem
% infrastructure, and the notation shared by the complete series.
% v2.9 (release v1.43) made the Paper-I-A--VII-C AUX registry automatic and
% centralized the bibliography and all external-reference ownership.
% v2.8 release-registry refresh (v1.33) registers the rooted generated-parent
% cascade and the POL/REG/free-law/11+4/Born-group closure revisions.  No
% shared mathematical notation or Paper-II-A fundamental principle is changed.
% v2.8 release-registry refresh (v1.32) registers the universal generated-
% continuation revisions in Papers II-A, III-A, III-C, IV, and VIII.  No
% shared mathematical notation or Paper-II-A fundamental principle is changed.
% v2.8 release-registry refresh (v1.30) registers the CRP4/projective-state
% source revisions and their direct registry-only consumers.  No shared
% mathematical notation or Paper-II-A axiom is changed.
% v2.7 release-registry refresh (v1.27) updates only the current external
% Paper-I-A--VII-C basenames after the placement/assumption synchronization; no
% shared mathematical notation or Paper-II-A axiom is changed.
% v2.7 synchronizes the geometric defect-repair revision: Paper-II-A v3.42,
% Paper-III-C v1.26, Paper-III-D v1.27, Paper-III-E v1.32, and Paper-III-F
% v1.34.  It registers the common trace-defect/spectral-filtration theorem,
% its owner-typed response pullback, and the reduced GRED and TERM interfaces.
% No shared mathematical notation is changed.
% v2.6 synchronizes Paper-III-E v1.31's cone-type repair, minimal-response
% classification, and explicit classical action-role boundary, together with
% Paper-III-F v1.32's downstream response-atlas synchronization.  No shared
% mathematical notation is changed.
% v2.5 synchronizes the Paper-III-D v1.26 faithful-trace-duality proof,
% exact ordered-owner joint locus, and typed alternative-presentation
% bookkeeping repair, together with its Paper-III-E and closure-frontier
% references.  No shared mathematical notation is changed.
% v2.4 synchronizes the Paper-III-C v1.25 orbit/process/owner repair, its
% downstream Papers III-B--F, V, VI, VIII, and IX revisions, and Born-v6_13.
% v2.3 synchronizes the Paper-II-A v3.41 witness/support/principal inheritance
% layer, the Paper-III-A v1.24 joint-realizability/readout firewall, all dependent
% host-rate revisions, and the current Paper-I-A--VII-B source registry.
% v2.2 synchronizes the generated active-realization, determinant-depth,
% dimension-weighted hierarchy, and typed cosmological-event interfaces.
% v2.1 synchronizes the exceptional-support interface revisions: Paper I-A
% v5.9, Paper II-A v3.30, Paper II-E/F v3.7, and Paper IV v4.16.
% v2.0 introduces the stage-safe G0 occurrence notation, synchronizes the
% active Paper-I-A--IV source registry (current rigorous Paper-III closure:
% III-A v1.17, III-B v1.22, III-C--E v1.17, III-F v1.14, IV v4.14), and makes
% external-reference self-skip depend on logical paper role rather than a
% versioned job name.
% ======================================================================

% ======================================================================
% macros.tex — Shared macros for the Born Geometry series
% ======================================================================

% ======================================================================
% 1. PACKAGES, BIBLIOGRAPHY, COLORS, AND UNIT DEFINITIONS
% ======================================================================
% Core packages shared across the complete Paper-I-A--VII-C series.  Loading them here makes
% the macro file a genuine series preamble while remaining compatible with existing
% paper preambles that load the same packages before \input.
\usepackage{amsmath}
\usepackage{amssymb}
\usepackage{amsthm}
\usepackage{mathrsfs}
\usepackage{array}
\usepackage{booktabs}
\usepackage{xcolor}
\makeatletter
\@ifpackageloaded{geometry}{}{\usepackage{geometry}\geometry{margin=1in}}
\makeatother
\usepackage{mdframed}

% Shared support packages required by the macros below and by most papers.
\usepackage{ragged2e}
\usepackage{longtable}
\usepackage{tabularx}
\usepackage{enumitem}
\usepackage[most]{tcolorbox}

% Series-wide numeric bibliography convention using the portable BibTeX path.
% All Papers I-A--VII-C use the ordinary \cite interface.
\usepackage[numbers,sort&compress]{natbib}
\providecommand{\printbibliography}{%
  \bibliographystyle{plainnat}%
  \bibliography{Born-v6_59}%
}

% Shared colors.  \providecolor avoids clobbering a paper-specific override.
\providecolor{postulatecolor}{RGB}{230,242,255}
\providecolor{compare}{RGB}{230,245,230}
\providecolor{definitioncolor}{RGB}{230,242,255}

\providecolor{openriskcolor}{RGB}{255,235,235}
\providecolor{bridgecolor}{RGB}{240,240,255}
\providecolor{gold}{RGB}{255,215,0}

% ======================================================================
% 2. CUSTOM ENVIRONMENTS
% ======================================================================
% Common table column types used in all papers.  Guarded so that a paper may
% define its own versions before loading this file.
\makeatletter
\@ifundefined{NC@find@L}{%
  \newcolumntype{L}[1]{>{\RaggedRight\arraybackslash}p{#1}}%
}{}
\@ifundefined{NC@find@C}{%
  \newcolumntype{C}[1]{>{\Centering\arraybackslash}p{#1}}%
}{}
\@ifundefined{NC@find@Y}{%
  \newcolumntype{Y}{>{\RaggedRight\arraybackslash}X}%
}{}
\makeatother

% ======================================================================
% 3. THEOREM STYLES AND ENVIRONMENTS
% ======================================================================
% Custom underlined theorem style (useful for remarks or notes)
\newtheoremstyle{underlined}
  {3pt}   % Space above
  {3pt}   % Space below
  {}      % Body font (empty means normal/roman text, typical for remarks)
  {}      % Indent amount
  {}      % Theorem head font (can be \itshape, \bfseries, etc.)
  {.}     % Punctuation after theorem head
  { }     % Space after theorem head
  {\underline{\thmname{#1}\thmnumber{ #2}}\thmnote{ (#3)}} % Theorem head formatting

% Standard plain theorem style environments.  Each environment is guarded so
% the shared macro file can be loaded after a paper-specific theorem setup.
\theoremstyle{plain}
\makeatletter
\@ifundefined{theorem}{\newtheorem{theorem}{Theorem}[section]}{}
\@ifundefined{lemma}{\newtheorem{lemma}{Lemma}[section]}{}
\@ifundefined{proposition}{\newtheorem{proposition}{Proposition}[section]}{}
\@ifundefined{corollary}{\newtheorem{corollary}{Corollary}[section]}{}
\makeatother

% Definition style environments (typically non-italicized body)
\theoremstyle{definition}
\makeatletter
\@ifundefined{definition}{\newtheorem{definition}{Definition}[section]}{}
\@ifundefined{postulate}{\newtheorem{postulate}{Postulate}[section]}{}
\@ifundefined{closure}{\newtheorem{closure}{Closure}[section]}{}
\makeatother

% Applying the custom underlined style
\theoremstyle{underlined}
\makeatletter
\@ifundefined{remark}{\newtheorem{remark}{Remark}[section]}{}
\@ifundefined{convention}{\newtheorem{convention}{Convention}[section]}{}
\@ifundefined{principle}{\newtheorem{principle}{Principle}[section]}{}
\makeatother

% Framed proof environment.
% Every \begin{proof} ... \end{proof} is automatically framed/shaded.
% This avoids mdframed/trivlist nesting, which can drop the top border after
% theorem-like environments.
\makeatletter
\renewenvironment{proof}[1][\proofname]{%
  \par\addvspace{6pt}%
  \noindent
  \begin{tcolorbox}[%
    enhanced,
    breakable,
    colback=black!4,
    colframe=black!35,
    boxrule=0.35pt,
    arc=1pt,
    left=6pt,
    right=6pt,
    top=4pt,
    bottom=4pt,
    boxsep=0pt,
    before skip=0pt,
    after skip=6pt,
    parbox=false
  ]%
  \pushQED{\qed}%
  \normalfont\itshape #1\@addpunct{.}\enspace
  \normalfont\ignorespaces
}{%
  \unskip\nobreak\hfill\popQED
  \end{tcolorbox}%
  \par\addvspace{2pt}%
}
\makeatother

\mdfdefinestyle{postulatebox}{
  backgroundcolor=yellow!10,
  linecolor=black,
  linewidth=0.4pt,
  skipabove=\baselineskip,
  skipbelow=\baselineskip,
  nobreak=true
}

% ======================================================================
% PROOF-READABILITY HELPERS
% ======================================================================

\newcommand{\ProofPlan}[1]{%
  \par\smallskip
  \noindent\textbf{Proof plan. }#1%
  \par\smallskip
}

\newlist{proofsteps}{enumerate}{2}
\setlist[proofsteps,1]{
  label=\textbf{(\alph*)},
  ref=\alph*,
  leftmargin=2.2em,
  itemsep=0.45\baselineskip,
  topsep=0.35\baselineskip
}
\setlist[proofsteps,2]{
  label=\textbf{(\roman*)},
  ref=\roman*,
  leftmargin=2.0em,
  itemsep=0.25\baselineskip,
  topsep=0.2\baselineskip
}


% ======================================================================
% 4. NUMBER SYSTEMS AND ALPHABETS
% ======================================================================
\providecommand{\RR}{\mathbb{R}} % Real numbers
\providecommand{\PP}{\mathbb{P}} % Projective space / Probability
\providecommand{\CC}{\mathbb{C}} % Complex numbers
\providecommand{\ZZ}{\mathbb{Z}} % Integers
\providecommand{\NN}{\mathbb{N}} % Natural numbers
\providecommand{\HH}{\mathbb{H}} % Quaternions / Hyperbolic space
\providecommand{\LFRayGeom}{\mathsf{LF\!\mbox{-}RAY\!\mbox{-}GEOM}}
\providecommand{\LFRep}{\mathsf{LF\!\mbox{-}REP}}
\providecommand{\LFState}{\mathsf{LF\!\mbox{-}STATE}}
\providecommand{\LFLim}{\mathsf{LF\!\mbox{-}LIM}}
\providecommand{\LFScat}{\mathsf{LF\!\mbox{-}SCAT}}
% Number systems / projective spaces
\providecommand{\CP}{\mathbb{CP}}
\providecommand{\RP}{\mathbb{RP}}

% ======================================================================
% 5. GROUPS AND LIE THEORY
% ======================================================================
\providecommand{\SO}{\ensuremath{\mathrm{SO}}}     % Special Orthogonal group
\providecommand{\SU}{\ensuremath{\mathrm{SU}}}     % Special Unitary group
\providecommand{\Spin}{\ensuremath{\mathrm{Spin}}} % Spin group
\providecommand{\CS}{\ensuremath{\mathrm{CS}}}     % Chern-Simons (capitalized)
\providecommand{\cs}{\ensuremath{\mathrm{cs}}}     % Chern-Simons (lowercase)

% ======================================================================
% 6. COMMON MANIFOLDS AND TOPOLOGY
% ======================================================================
\providecommand{\Sone}{S^1}                 % Standard circle

% ======================================================================
% 7. MATH OPERATORS (Declared via \DeclareMathOperator)
% ======================================================================
\DeclareMathOperator{\End}{End}   % Endomorphisms
\DeclareMathOperator{\Aut}{Aut}   % Automorphisms
\DeclareMathOperator{\spec}{spec} % Lowercase spectrum
\DeclareMathOperator{\tr}{tr}     % Lowercase trace (if you already have \Tr)
\DeclareMathOperator{\supp}{supp} % Support
\DeclareMathOperator{\TrA}{Tr_{A}} % Partial trace over subsystem A
\DeclareMathOperator{\TrB}{Tr_{B}} % Partial trace over subsystem B
% Lie/algebra operators
\DeclareMathOperator{\Ad}{Ad}
\DeclareMathOperator{\Sym}{Sym}
\DeclareMathOperator{\su}{su}
% ======================================================================
% 8. STANDARD OPERATOR-LIKE SYMBOLS (Safe definitions via \providecommand)
% ======================================================================
\providecommand{\vol}{\operatorname{vol}}   % Volume form
\providecommand{\Vol}{\operatorname{Vol}}   % Total volume
\providecommand{\Ric}{\operatorname{Ric}}   % Ricci tensor
\providecommand{\Rm}{\operatorname{Rm}}     % Riemann tensor
\providecommand{\Tr}{\operatorname{Tr}}     % Uppercase trace
\providecommand{\Spec}{\operatorname{Spec}} % Uppercase spectrum
\providecommand{\diag}{\operatorname{diag}} % Diagonal matrix
\providecommand{\rank}{\operatorname{rank}} % Rank
\providecommand{\scal}{\mathcal R} % Scalar curvature
\providecommand{\Cl}{\operatorname{Cl}}     % Clifford algebra
\providecommand{\diam}{\operatorname{diam}} % Diameter

% ======================================================================
% 9. FUNCTIONS, TRIGONOMETRY, AND CALCULUS
% ======================================================================

% Codifferential naming conventions
% (Avoid \ddag: commonly defined by physics.sty and symbol packages)
\providecommand{\codiff}{\mathrm{d}^{\dagger}} % Codifferential operator

% Quantum mechanics bra-ket notation (needed for basis tables)
\providecommand{\ket}[1]{\lvert #1 \rangle}
\providecommand{\bra}[1]{\langle #1 \rvert}
\providecommand{\braket}[2]{\langle #1 \mid #2 \rangle}

% ======================================================================
% 10. PHYSICAL SCALES AND CONSTANTS (Paper 1c / Born Geometry)
% ======================================================================
\providecommand{\Rp}{\ensuremath{R_{\mathrm P}}}     % matched common EH curvature scale (not defined from observed G)
\providecommand{\hBorn}{\ensuremath{h_{\mathrm B}}} % physical representative 2\pi\hbar of one primitive Born scalar action cell
\providecommand{\UEHjoint}{\ensuremath{U_{J}}} % joint-Dirac EH curvature-coordinate scale 8 kappa_J
\providecommand{\Gread}{\ensuremath{G_{\rm scal}}} % scalar-action coordinate c_op^3/(h_B U); physical matching explicit
\providecommand{\GreadJoint}{\ensuremath{G_{J,\mathrm{raw}}}} % raw ordinary-trace coordinate = G_scal/96
\providecommand{\Rds}{R_{\mathrm{dS}}}    % de Sitter Curvature
\providecommand{\Mpl}{M_{\mathrm{Pl}}}    % Planck mass
\providecommand{\lp}{\ell_{\mathrm{Pl}}}  % Planck length
\providecommand{\Pmax}{P_{\mathrm{max}}}  % Maximum momentum
\providecommand{\Mplbar}{\bar{M}_{\mathrm{Pl}}} % Reduced Planck mass
\providecommand{\eff}{\mathrm{eff}}       % Subscript for effective quantities
\providecommand{\Cosmo}{\ensuremath{\Lambda}}                % Cosmological constant (use \Cosmo in source)
\providecommand{\LCDM}{\ensuremath{\Cosmo\mathrm{CDM}}}  % \LambdaCDM (typesets with \Cosmo)
\providecommand{\Lcut}{\ensuremath{\Lambda_{\mathrm{cut}}}}  % Cutoff scale (brace when used as script)
\providecommand{\Rcs}{\ensuremath{R_{\mathrm{cs}}}}      % Configuration-space radius
\providecommand{\CosmoEff}{\ensuremath{\Cosmo_{\eff}}}   % Effective cosmological constant

% Disambiguation: \Rmom is the momentum-space (S^3) radius. 
% Reserve plain R for the 4D curvature scalar in f(R).
\providecommand{\Rmom}{\ensuremath{R_{\mathrm{mom}}}} 
\providecommand{\Punit}{P_{0}} % Momentum unit set by \Rmom (= \hbar/\Rmom)
\providecommand{\Rvec}{\ensuremath{R_{\mathrm{vec}}}} % invariant vector-space radius in the reduced round sector
% Specific numeric values and expressions
\providecommand{\Vlogcap}{L_{\mathrm{Hopf}}^{(1)}}
\providecommand{\Hopf}{\operatorname{Hopf}}
% Scalar phase (S^1) sector + universal cover notation (Paper I-A)
\providecommand{\tper}{T_\star}      % Time period/scale
\providecommand{\Eper}{E_\star}      % Energy period/scale
\providecommand{\ttilde}{\tilde t}   % Universal cover time
\providecommand{\Etilde}{\tilde E}   % Universal cover energy

% ======================================================================
% 10A. SERIES-WIDE LOGICAL / NOTATIONAL DISAMBIGUATIONS
% ======================================================================
% Use these symbols to avoid collisions among: (i) the Paper--III
% no-insertion/geometric-closure class, (ii) Paper--I-A trimer diagonals,
% (iii) Paper--V-A/V-C family-space diagonals, (iv) scalar-index running,
% and (v) protected shell-clock labels.
\providecommand{\Rprim}{\ensuremath{\mathsf R_{\mathrm{prim}}^{\Sigma}}} % primitive finite structural readouts
\providecommand{\Rloop}{\ensuremath{\mathsf R_{\circ}^{\Sigma}}} % closed-loop readout role
\providecommand{\Rbdry}{\ensuremath{\mathsf R_{\partial}^{\Sigma}}} % boundary-conditioned readout role
\providecommand{\Rgen}{\ensuremath{\mathsf R_{\mathrm{fin}}^{\Sigma}}} % finite generated closure of primitive roles
\providecommand{\PhysOcc}{\operatorname{PhysOcc}}
\providecommand{\Actual}{\operatorname{Actual}}
\providecommand{\GeomLic}{\operatorname{GeomLic}}
\providecommand{\InvDist}{\operatorname{InvDist}}
\providecommand{\Dcl}{\ensuremath{\mathsf{D}_1}}      % unmarked positive Dirac-even closed spectral-readout form class
\providecommand{\Dheat}{\ensuremath{\mathsf D_{\mathrm{heat}}}} % one-parameter heat-semigroup subclass
\providecommand{\DtriI}{\ensuremath{D^{\mathrm{tri}}_1}} % Paper--I-A cyclic-trimer diagonal
\providecommand{\DtriII}{\ensuremath{D^{\mathrm{tri}}_2}}
\providecommand{\DfamI}{\ensuremath{D^{\mathrm{fam}}_1}} % Family-space diagonal selector
\providecommand{\DfamII}{\ensuremath{D^{\mathrm{fam}}_2}}
\providecommand{\alpharun}{\ensuremath{\alpha_{\mathrm{run}}}} % scalar spectral-index running
\providecommand{\Ksh}{\ensuremath{K_{\mathrm{sh}}}} % protected shell-clock coordinate
\providecommand{\ddef}{\ensuremath{\delta_{\mathrm{def}}}} % universal defect scalar after Paper--III realization

% ======================================================================
% 10B. SERIES-WIDE LOGICAL STATUS TAGS
% ======================================================================
% Upright small-caps wrapper for audit/status labels.  It prevents LaTeX from
% trying to use unavailable small-caps italic font shapes inside theorem bodies
% such as TU/lmr/m/scit.
\providecommand{\AuditLabel}[1]{}
\renewcommand{\AuditLabel}[1]{{\normalfont\scshape #1}}

% These tags are bookkeeping labels for consistency audits. They do not
% create new mathematical hypotheses; they record how a datum is being used
% in the proof chain.
\providecommand{\StatusPrimitive}{}
\renewcommand{\StatusPrimitive}{\AuditLabel{Primitive}}
\providecommand{\StatusDerived}{}
\renewcommand{\StatusDerived}{\AuditLabel{Derived}}
\providecommand{\StatusImported}{}
\renewcommand{\StatusImported}{\AuditLabel{Imported}}
\providecommand{\StatusConditional}{}
\renewcommand{\StatusConditional}{\AuditLabel{Conditional}}
\providecommand{\StatusBridge}{}
\renewcommand{\StatusBridge}{\AuditLabel{Bridge-derived}}
\providecommand{\StatusComparison}{}
\renewcommand{\StatusComparison}{\AuditLabel{Comparison-only}}

% ======================================================================
% 10C. SERIES-WIDE CLAIM-WORDING AUDIT LABELS
% ======================================================================
% These are textual audit labels used to keep modal language consistent
% across the series.  They introduce no new mathematical assumptions.
\providecommand{\ClaimForced}{}
\renewcommand{\ClaimForced}{\AuditLabel{forced}}
\providecommand{\ClaimUnique}{}
\renewcommand{\ClaimUnique}{\AuditLabel{unique}}
\providecommand{\ClaimCanonical}{}
\renewcommand{\ClaimCanonical}{\AuditLabel{canonical}}
\providecommand{\ClaimMinimal}{}
\renewcommand{\ClaimMinimal}{\AuditLabel{minimal}}
\providecommand{\ClaimClosed}{}
\renewcommand{\ClaimClosed}{\AuditLabel{closed}}
\providecommand{\ClaimPrediction}{}
\renewcommand{\ClaimPrediction}{\AuditLabel{prediction}}

% ======================================================================
% 10D. SERIES-WIDE THEOREM-CHAIN AUDIT LABELS
% ======================================================================
% These labels mark the direction of load-bearing theorem chains. They are
% audit annotations only; they introduce no new mathematical assumptions.
\providecommand{\ChainInput}{}
\renewcommand{\ChainInput}{\AuditLabel{Chain input}}
\providecommand{\ChainOutput}{}
\renewcommand{\ChainOutput}{\AuditLabel{Chain output}}
\providecommand{\ChainGuard}{}
\renewcommand{\ChainGuard}{\AuditLabel{Chain guard}}
\providecommand{\ChainNoBackflow}{}
\renewcommand{\ChainNoBackflow}{\AuditLabel{No backflow}}

% ======================================================================
% 10E. SERIES-WIDE LOCAL THEOREM-AUDIT LABELS
% ======================================================================
% These labels mark a statement-level dependency audit around individual
% theorem/definition nodes. They are bookkeeping labels only.
\providecommand{\LocalAudit}{}
\renewcommand{\LocalAudit}{\AuditLabel{Local audit}}
\providecommand{\UsesOnly}{}
\renewcommand{\UsesOnly}{\AuditLabel{Uses only}}
\providecommand{\Exports}{}
\renewcommand{\Exports}{\AuditLabel{Exports}}
\providecommand{\DoesNotUse}{}
\renewcommand{\DoesNotUse}{\AuditLabel{Does not use}}

% ======================================================================
% 10F. SERIES-WIDE NUMERICAL-ROW AUDIT LABELS
% ======================================================================
% These labels classify individual rows in numerical, benchmark, status,
% and comparison tables.  They are audit annotations only and do not change
% any theorem, value, formula, or comparison number.
\providecommand{\RowAudit}{}
\renewcommand{\RowAudit}{\AuditLabel{Numerical-row audit}}
\providecommand{\RowTheorem}{}
\renewcommand{\RowTheorem}{\AuditLabel{Theorem-fixed row}}
\providecommand{\RowBridge}{}
\renewcommand{\RowBridge}{\AuditLabel{Bridge-fixed row}}
\providecommand{\RowRepresentative}{}
\renewcommand{\RowRepresentative}{\AuditLabel{Representative row}}
\providecommand{\RowComparison}{}
\renewcommand{\RowComparison}{\AuditLabel{Comparison row}}
\providecommand{\RowCalibrated}{}
\renewcommand{\RowCalibrated}{\AuditLabel{Calibrated row}}
\providecommand{\RowEstimate}{}
\renewcommand{\RowEstimate}{\AuditLabel{Estimate row}}
\providecommand{\RowOpen}{}
\renewcommand{\RowOpen}{\AuditLabel{Open row}}
\providecommand{\RowNoSelector}{}
\renewcommand{\RowNoSelector}{\AuditLabel{No selector}}

% ======================================================================
% 10G. TABLE HELPERS
% ======================================================================
% Table helpers used by Paper V-C.  The unused readout-label preset has been
% removed; only active table helpers remain.
\providecommand{\gpsection}[1]{%
  \addlinespace[0.55em]
  \multicolumn{3}{@{}l}{\textsc{#1}}\\[-0.25em]
  \midrule
}
\providecommand{\mathcell}[1]{%
  \(\begin{gathered}#1\end{gathered}\)%
}

\newcommand{\DerBy}[1]{%
  \begin{tabular}[t]{@{}l@{}}#1\end{tabular}%
}

\newcommand{\TableNotes}[1]{%
  \par\medskip
  \begin{minipage}{0.94\linewidth}
  \footnotesize\emph{Notes.} #1
  \end{minipage}\par
}
% ======================================================================
% 10H. PAPER VI-A / QCD NOTATION
% ======================================================================
% Shared symbols used by Paper VI-A and downstream QCD/Yang--Mills handoffs.
\providecommand{\Pqcd}{\mathfrak P_{\rm VIII}^{\rm QCD}}
\providecommand{\Iqcd}{\mathfrak I_{\rm QCD}^{(7)}}
\providecommand{\muQCD}{\mu_{\rm QCD}^{\rm int}}
\providecommand{\LambdaQCD}{\Lambda_{\rm QCD}}
\providecommand{\LambdaBH}{\Lambda_{\rm BH}^{\rm QCD}}
\providecommand{\HQCD}{\mathcal H_{\rm QCD}^{\rm BH}}
\providecommand{\HQCDfin}{H_{\rm QCD}^{\rm fin}}
\providecommand{\Hraw}{\mathcal H_{\rm raw}}
\providecommand{\HphysQCD}{\mathcal H_{\rm phys}^{\rm QCD}}
\providecommand{\Vcol}{V_C}
\providecommand{\Ps}{P_{\rm sing}}
\providecommand{\PiOct}{\Pi_8}
\providecommand{\MSbar}{\overline{\rm MS}}
\providecommand{\YM}{\mathrm{YM}}
\providecommand{\Ktri}{K_{\triangle}}
\providecommand{\KYM}{K_{\YM}^{\rm BH}}
\providecommand{\aBH}{a_{\rm BH}}
\providecommand{\epsBH}{\varepsilon_{\rm BH}}
\newcommand{\RowConstructive}{\textnormal{\textsc{Constructive}}}
\newcommand{\RowObligation}{\textnormal{\textsc{Bridge obligation}}}
\newcommand{\RowTarget}{\textnormal{\textsc{Target row}}}
\newcommand{\RowValidation}{\textnormal{\textsc{Validation row}}}
\providecommand{\slashed}[1]{\not{#1}}

% ======================================================================
% 10I. PAPER VI-B / CONTINUOUS-QFT NOTATION
% ======================================================================
% Shared symbols used by Paper VI-B and any downstream continuous-QFT or
% Yang--Mills handoff paper.  All definitions are guarded so individual
% papers may override them locally before loading this file if needed.
\providecommand{\A}{\mathfrak A}
\providecommand{\Hphys}{\mathcal H_{\rm phys}}
\providecommand{\QCD}{\mathrm{QCD}}
\providecommand{\BH}{\mathrm{BH}}
\providecommand{\loc}{\mathrm{loc}}

\providecommand{\RBH}{\mathcal R_{\rm BH}}
\providecommand{\DG}{\mathfrak D_G}
\providecommand{\PW}{\mathcal{P\!W}} % Peter--Weyl spectral-symbol subspace
\providecommand{\Toe}{\operatorname{Toe}} % Toeplitz compression/representative
\providecommand{\Hol}{\operatorname{Hol}} % Holonomy

% ======================================================================
% 11. LATE-LOADED PACKAGES (Must be at the end of the preamble)
% ======================================================================
% Load late packages only if a paper has not already loaded them, avoiding
% option clashes in paper-specific preambles.
\makeatletter
\@ifpackageloaded{hyperref}{}{\usepackage{hyperref}}
\makeatother
% --- Series-wide disambiguations (Pass 3) ---
\providecommand{\diracdelta}{\delta_{\rm D}} % Dirac delta distribution (delta reserved for hierarchy)
\providecommand{\Hg}{H_{\mathrm g}} % Hodge harmonic projector (keep H for Hubble)
\providecommand{\Gr}{G_{\mathrm r}} % Green operator (keep G for Newton)
\providecommand{\FD}{\mathrm{D}} % functional/variational derivative symbol (delta reserved)
\providecommand{\Id}{\mathbf{1}} % identity/Kronecker tensor (avoid \delta_{ij})


% BEGIN INLINE SOURCE: Paper00-ActiveExternalReferences-v2_130.tex
% Framework v2.191: current reference binding; existing definitions retained.
% Framework v2.190: current22-root registry with native clock-evolution proofs.
% Framework v2.189: exact current22-root external registry, including retained v2.188 hardening owners.
% Framework v2.187: current registry synchronization only; source release v2.186 preserved.
% Framework v2.186: authoritative current 22-paper registry.
% Framework v2.182: all current 22 logical paper AUX bindings.
% Framework v2.170: current I-B v1.43 external AUX binding.
% Framework v2.167: exact current 22-root registry after CP--BL2 integration.
% Framework v2.166: exact current 22-root compact-release AUX registry.
% Framework v2.165: authoritative current root bindings.
% Framework v2.164: exact current 22 manuscript AUX bindings.
% Framework v2.162: finite metric, functional Hilbert and scoped Ward integration.
% Framework v2.161: current 22 manuscript AUX bindings; source only.
% Framework v2.156 / registry v2.113: synchronized current 22 manuscript roots; source-only registry.
% Framework v2.143 / registry v2.111: active manuscript AUX basenames; no generated AUX or PDF included.
% ======================================================================
% Central external-reference imports, numbers only, no PDF jumping
% ======================================================================

% xr-hyper is loaded centrally so that papers may retain their existing
% \usepackage{xr-hyper} lines without redefining \externaldocument after this
% registry is installed.  If plain xr was loaded earlier, retain it.
\makeatletter
\@ifpackageloaded{xr-hyper}{}{%
  \@ifpackageloaded{xr}{}{\usepackage{xr-hyper}}%
}
\makeatother

% Make external-document registration idempotent.  The registry identifies a
% logical paper by its Paper<number><optional letter> prefix.  Prefix-qualified
% imports remain independent.  Current paper roots contain no local
% \externaldocument calls; this guard also makes accidental duplicates harmless.
\ExplSyntaxOn
\prop_new:N \g_bh_external_refs_prop
\cs_new_eq:NN \BHOriginalExternalDocument \externaldocument

\cs_new_protected:Npn \bh_external_document:nnnn #1#2#3#4
 {
  \tl_set:Nn \l_tmpa_tl {#3}
  \regex_replace_once:nnN
   { \A (?: xref/ )? (Paper[0-9]+[A-G]?) .* \Z }
   { \1 }
   \l_tmpa_tl
  \tl_set:Nx \l_tmpb_tl { #1 | \tl_use:N \l_tmpa_tl }
  \prop_if_in:NVTF \g_bh_external_refs_prop \l_tmpb_tl
   {
    \PackageInfo{BornHopf}
     {Skipping~duplicate~external~refs~for~#3}
   }
   {
    \prop_gput:NVn \g_bh_external_refs_prop \l_tmpb_tl {#3}
    \IfNoValueTF{#4}
     { \BHOriginalExternalDocument[#1][#2]{#3} }
     { \BHOriginalExternalDocument[#1][#2]{#3}[#4] }
   }
 }

\RenewDocumentCommand{\externaldocument}{O{} O{} m o}
 { \bh_external_document:nnnn {#1}{#2}{#3}{#4} }

\NewDocumentCommand{\BHExternalRefs}{m}
 {
  % Do not import the aux file of the logical paper currently being compiled,
  % even when the job and registry use different version suffixes.
  \tl_set:Nx \l_tmpa_tl { \c_sys_jobname_str }
  \regex_replace_once:nnN
   { \A (Paper[0-9]+[A-G]?) .* \Z }
   { \1 }
   \l_tmpa_tl
  \tl_set:Nn \l_tmpb_tl { #1 }
  \regex_replace_once:nnN
   { \A (Paper[0-9]+[A-G]?) .* \Z }
   { \1 }
   \l_tmpb_tl
  \str_if_eq:eeTF { \tl_use:N \l_tmpa_tl } { \tl_use:N \l_tmpb_tl }
   {
    \PackageInfo{BornHopf}
     {Skipping~external~refs~for~#1~because~this~is~the~current~document}
   }
   {
    \IfFileExists{xref/#1.aux}
     { \externaldocument[][nocite]{xref/#1}[#1.pdf] }
     {
      \IfFileExists{#1.aux}
       { \externaldocument[][nocite]{#1} }
       {
        \PackageInfo{BornHopf}
         {External~aux~file~#1.aux~not~found;~no~labels~imported}
       }
     }
   }
 }
\ExplSyntaxOff

% External references are automatic.  This file is the sole owner of both the
% bibliography and the current Paper-I-A--VII-D AUX registry (30 roots including I-C).  Missing auxiliaries
% are informational on the first pass and are populated by the source-only
% two-pass series build.  Basenames deliberately omit temporary download/copy
% suffixes such as ``(1)'' or ``-1''.
  % Paper I-A: round S3 compendium.
  \BHExternalRefs{Paper1A-Geometry-v5_181}

  % Paper I-B: complete exceptional geometry and finite spectral structures.
  \BHExternalRefs{Paper1B-ExceptionalGeometry-v1_25}
  % Paper I-D: technical companion of I-A through I-C.
  \BHExternalRefs{Paper1D-TechnicalMaterial-v1_69}
  % Paper I-C: reciprocal M8 compendium.
  \BHExternalRefs{Paper1C-M8Geometry-v1_146}

\ifdefined\BHMathOnly
  % Papers I-A and I-B import only the mathematical atlas labels.
\else
  % Paper II: five physical-mathematical main papers plus one technical proof volume.
  \BHExternalRefs{Paper2A-GeometricConstitution-v3_227}
  \BHExternalRefs{Paper2B-MomentumCarrier-v2_147}
  \BHExternalRefs{Paper2C-ReciprocalLocalization-v2_176}
  \BHExternalRefs{Paper2D-SymmetrySpectra-v2_145}
  % The former II-F export labels remain in II-E; current II-F is the technical volume.
  \BHExternalRefs{Paper2E-KinematicExports-v3_140}
  % Paper II-F: complete technical proof and calculation owner for II-A--II-E.
  \BHExternalRefs{Paper2F-TechnicalMaterial-v1_71}

  % Paper III: Dynamics and Intrinsic Geometric Constitution (six components).
  \BHExternalRefs{Paper3A-GeneratedDynamics-v1_181}
  \BHExternalRefs{Paper3B-SpectralAction-v1_192}
  \BHExternalRefs{Paper3C-ProtectedHierarchy-v1_195}
  \BHExternalRefs{Paper3D-RunningResponse-v1_167}
  \BHExternalRefs{Paper3E-GravityAction-v1_203}
  \BHExternalRefs{Paper3F-CompleteDynamics-v1_279}
  % Paper III-G: complete preserved technical proof owner.
  \BHExternalRefs{Paper3G-TechnicalMaterial-v1_57}

  % Paper IV-A: generated finite quantum mechanics.
  \BHExternalRefs{Paper4A-FiniteQuantumMechanics-v4_217}
  % Paper IV-B: complete finite quantum processes and physical response.
  \BHExternalRefs{Paper4B-QuantumResponse-v1_55}
  % Paper IV-C: extended proofs and finite technical calculations.
  \BHExternalRefs{Paper4C-TechnicalMaterial-v1_50}

  % Paper V-A: The Standard Model.
  \BHExternalRefs{Paper5A-FiniteMatterGeometry-v3_218}

  % Paper V-B: Finite Spectrum and Flavor Closure.
  \BHExternalRefs{Paper5B-TerminalSpectrum-v5_202}

  % Paper V-C: Finite Interaction Boundaries of the Standard Model.
  \BHExternalRefs{Paper5C-InteractionResponse-v2_184}
  % Paper V-D: preserved technical proof and calculation owner.
  \BHExternalRefs{Paper5D-TechnicalMaterial-v1_58}

  % Paper VI-A: Gravity and Closed Cosmology.
  \BHExternalRefs{Paper6A-GravityCosmology-v2_268}

  % Paper VI-B: Geometric Dark Sector and Cosmological Readouts.
  \BHExternalRefs{Paper6B-CosmologyReadouts-v1_215}

  % Paper VI-C: complete technical proof and comparison owner.
  \BHExternalRefs{Paper6C-TechnicalMaterial-v1_53}

  % Paper VII-A: Finite QCD Kernel and Nonperturbative Closure.
\BHExternalRefs{Paper7A-QCDKernel-v1_183}

  % Paper VII-B: Hadronic, Thermal, and Nuclear QCD Readouts.
  \BHExternalRefs{Paper7B-QCDReadouts-v1_167}

  % Paper VII-C: Continuous Quantum Field Theory.
  \BHExternalRefs{Paper7C-ContinuousQFT-v2_208}

  % Paper VII-D: complete technical material for VII-A through VII-C.
  \BHExternalRefs{Paper7D-TechnicalMaterial-v1_60}

\fi

% END INLINE SOURCE: Paper00-ActiveExternalReferences-v2_130.tex


\hypersetup{hidelinks}

\usepackage{xr-hyper}

\title{\textbf{Born--Hopf Geometry III-A: Generated Dynamics and Clocked Reciprocal Constitution}\\[0.5em]
\large Dynamical input, co-realized process laws, and visible response}

\author{Peter Kugelmann\thanks{\texttt{pkugelmann@proton.me}}}
\date{Version 1.181 --- 20 November 2026 \\ Framework v3.069}
\setlength{\emergencystretch}{3em}
\begin{document}
\maketitle
\begin{abstract}
Paper III-A derives the dynamical reciprocal role from the constituted
momentum and position geometry of Papers I and II. The clocked Hopf operation,
exact-continuation statements, law--flow--history distinction, and common-action
co-consumption are given their precise geometric domains. A generated
continuation construction does not imply the occurrence of arbitrary physical
trajectories: the retained process and source must be co-realized.
The technical descent and alternative-construction proofs are collected in
Paper~III-G without re-opening a momentum selector.
\end{abstract}
\tableofcontents
\clearpage

\section{Geometric inputs and dynamical constitution}
\label{paper3a:sec:organized-scope}
\subsection{Constitutional inputs and paper ownership}
\label{p3a:sec:series-allocation}
Paper~III-A is the reciprocal-geometry entry of the Paper~III sequence.  It
imports the Paper-I-A reference objects and the current Paper-II operational
constitution, including G0 and Paper~II-C's premise-minimal \(E_1\)-evaluation
reciprocal closure.  The Born/GP2-generated reciprocal role, relation-free direct
product, and Born involution are therefore imported without
\(\mathsf P_{\rm mc}\).  Paper~III-A separately derives the
variance-opposite full metric-profile test action and certifies the derived
\((\mathrm{BR}_{\rm ind})\) and \((\mathrm{IND}_{xp})\) statuses on the
optional full-profile G0 test-occurrence domain.  The
momentum seed is an upstream result of Paper~II-B and is not a Paper~III
selection problem.  Later Paper~III-B action/spectral interfaces are used only
in explicitly marked forward synthesis statements; they do not enter the
native reciprocal-geometry proofs.

\subparagraph{Momentum-seed handoff.}
Paper~II-B proves
\[
   \boxed{\mathcal H_p\cong\mathcal H_1,\qquad m_p=1.}
\]
Generic word stages and lower-ideal families remain auxiliary mathematics and
later-role interfaces.  Accordingly no \(\mathsf P_{\rm sli}\), momentum-side
TWR, or level-one premise is used below to constitute the structural
momentum-role seed.  Its physical occurrence remains conditional on
the upstream transfer-operation G0 root and is co-realized by the
Paper-II-A/II-B theorem; no independent momentum witness remains.  The auxiliary tower and fixed-action selector countermodels are recorded
in Paper~III-G, Appendix~\ref{p3a:app:legacy-selector}.

\label{p3:part:dyn-core}

\subparagraph{Scope.} This section fixes the framework
prefix used by Paper~III, records the imported Paper-I-A/Paper-II data, and
constructs the Born-reciprocal completion from the already constituted
momentum parent and seed using the already discharged GP2 reciprocal-role
objectification and relation-free-composition antecedents.  Position-side
finite levels, local/global position--momentum comparison, reciprocal physical
scale, and all later action/state choices remain separate G0 interfaces.

\subsection{Constituted inputs and theorem scope}
\label{p3:sec:introduction}

\subsubsection{Framework prefix and theorem scope}
\label{p3:subsec:framework-prefix-status}

\subparagraph{Notation.} Throughout this paper we denote the (4D) scalar
curvature by $\scal$ and reserve $\Rmom$ for the already calibrated momentum-space
radius.  The reciprocal profile spectrum initially carries only the transported
profile metric.  When a physical configuration-space unit is required we use the
explicit reciprocal scale calibration
\[
 (\mathrm{CAL}_{xp}):\qquad
 \alpha_{xp}>0,\qquad d_x:=\alpha_{xp}d_{B_A},\qquad
 \Rcs:=\alpha_{xp}\Rmom.
\]
Thus $\Rcs$ is not derived by Born reciprocity alone; without
$(\mathrm{CAL}_{xp})$ the reciprocal theorem yields the isometric abstract role
$S^3_B$ with metric $d_{B_A}$.  We also reserve $\Cosmo$ for the cosmological
constant.  The structural momentum-role seed has $m_p=1$ by Paper~II-B and is
co-realized with the occurring transfer operations.  A
generic position-side label $m_x$ is used only as the derived least-resolution
statistic of a separately occurring finite frontier history; $\Lambda_D(m)$ remains an auxiliary Dirac-shell ceiling
and $\Lcut$ the spectral-action regulator.  We write $\Sigma_{\rm II}^{\rm role}$ for the
typed vector/scalar role signatures and exact/faithful operation grammar of
Paper~II-A, and write
\[
 \mathsf P_{\rm mc}
 :=
 \text{the metric-curry closure clause of
 Definition~\ref{p2:def:metric-partial-evaluation-closure}}.
\]
Thus the Paper-II A-role constitution uses the displayed role signature with
G0 and GP1--GP3, whereas every claim that the canonical mathematical distance
profiles form the \emph{full} physical reciprocal test doctrine additionally
uses $\mathsf P_{\rm mc}$ as a G0 test-occurrence certificate.  Reciprocal-role
existence itself instead follows from the Paper-II-C intrinsic
\(E_1\)-evaluation theorem without this clause.  We retain
$({\rm BR})\equiv(\mathrm{BR}_{\rm ind})$ and
$({\rm IND})\equiv(\mathrm{IND}_{xp})$ only as derived-status aliases.  In Paper~II-A they are not independent constitutive
interfaces: $({\rm BR}_{\rm ind})$ certifies the generated-role
objectification antecedent and $({\rm IND}_{xp})$ certifies absence of a
prior generated joint-support constraint.  Paper~II-C proves both statuses on
its intrinsic evaluation spectrum and invokes
Theorem~\ref{p2:thm:gp2-generated-reciprocal-role-closure}; Paper~III-A
re-proves them for its stronger metric-profile construction under
$\mathsf P_{\rm mc}$.  Later,
after the primitive scalar Poincar\'e class $c_{\rm sc}$ is derived, we use
the separate physical phase-unit calibration
\[
 (\mathrm{CAL}_{\hbar}):\qquad
 \hBorn:=2\pi\hbar,\qquad
 \left[\frac{\omega_{\rm sc}}{\hBorn}\right]=c_{\rm sc}.
\]
Thus \(\hBorn\) is the physical action carried by one primitive scalar Born
phase cell.  This fixes the dimensionful representative of the already derived
integral class; it does not choose its topology, sign, or integer multiplicity.  For
the clocked Hopf doctrine we likewise register the primitive-turn calibration
\[
 \begin{gathered}
 (\mathrm{CAL}_{t}):\qquad T_t>0\\
 =\text{the physical duration assigned}\\
 \text{to one positive primitive scalar turn}.
 \end{gathered}
\]
The structural positive primitive is already derived before this calibration;
$(\mathrm{CAL}_{t})$ only supplies its physical clock representative.  The
three displayed calibrations are G0 calibration certificates: their comparison
to a realized standard is not derived by reciprocity and creates no new
structural content.

Paper~III is the dynamical layer over the framework prefix
\[
        \mathfrak F_{\le2}
        :=
        \mathfrak F_{\le1}
        +
        \substack{\text{Paper-II vector-transfer constitution and scalar classification,}\\
        \text{plus the Paper-II-C intrinsic evaluation reciprocal closure}},
\]
where \(\mathfrak F_{\le1}\) is the Paper-I-A reference layer.  Paper~I-A fixes
the mathematical \(S^3\) theorem bank.  Inside
\(\Sigma_{\rm II}^{\rm role}+\mathrm{G0}+\mathrm{GP1--GP3}\), Paper~II derives the
nonabelian A-role \(S^3_A\), the scalar \(S^1_A\) role, and the faithful
transformation--test pairing.  Paper~II-C's intrinsic evaluation construction
supplies a Born/GP2-generated tagged B-sort and proves that no generated
joint-support relation is present; Paper~II-A's GP2 theorem then objectifies
the B-role and supplies the universal relation-free products
\(S^3_A\times S^3_B\) and \(S^1_A\times S^1_B\), with Born exchange generated
by the canonical comparison isomorphisms.  Under \(\mathsf P_{\rm mc}\),
Paper~III-A additionally realizes the full generated metric-profile family as
physical localization tests and proves its faithful opposite-variance
pullback law.  This strengthens the test doctrine but does not constitute the
reciprocal role anew.

The standing convention is therefore dependency-labelled:
\[
\boxed{
\begin{aligned}
  \Sigma_{\rm II}^{\rm role}+\mathfrak F_{\le2}
  &\;\vdash\;\text{GP2-derived reciprocal-product closure},\\
  \Sigma_{\rm II}^{\rm role}+\mathfrak F_{\le2}+\mathsf P_{\rm mc}
  &\;\vdash\;\text{full-profile G0 test realization}.
\end{aligned}
}
\]
Here GP1 finite/bounded admissibility, GP2 exact/host-faithful realization,
and GP3 irreducible/least reduction have their imported Paper-II meanings.
The words ``theorem'', ``unique'', and ``excluded'' always mean ``under the
hypotheses displayed in that statement''.  Without \(\mathsf P_{\rm mc}\),
the reciprocal role, product, and Born exchange still follow from the imported
Paper-II-C evaluation theorem; only the claim that every distance profile is a
physical test is absent.  With \(\mathsf P_{\rm mc}\), no additional BR/IND
postulate is required.  A
dimensionful reciprocal scale, a stronger observable-to-operator
position--momentum comparison, action, state, regularity, reconstruction, and
polarization still retain their own displayed theorem domains.

\subsubsection{Native dynamical problem}
\label{p3:subsec:native-dynamical-problem}

The Paper-II vector and scalar A-roles together with the Paper-II-C
evaluation-spectrum closure give the continuous vector product
$S^3_x\times S^3_p$.  On the optional full-profile G0 domain,
$\mathsf P_{\rm mc}$ supplies its contravariant action on all distance-profile
tests.  This smooth product
has $H^2=0$ and is not a global symplectic phase space.  Independently, the
finite full-product-invariant endomorphism algebra is abelian, so it contains
no exact invariant Heisenberg pair.  The separately prior-generated scalar
role is then Born-completed to a torus carrying the primitive Poincar\'e class.
Dimensionful calibration, action choice, state, regularity, position-side
finite level, Hamiltonian selection, and local dynamics retain their own
displayed theorem domains.

\subsubsection{Imported \texorpdfstring{$\mathfrak F_{\le2}$}{F<=2} data}
\label{p3:subsec:axioms}

Table~\ref{tab:p3-paperii-import-ledger} is the complete Paper-II theorem package used below.  The following definitions translate that package into the notation needed for Paper-III ledgers.


\begin{table}[H]
\centering
\scriptsize
\begingroup
\setlength{\tabcolsep}{3pt}
\renewcommand{\arraystretch}{1.16}
\begin{tabular}{@{}L{3.1cm}L{3.6cm}L{4.0cm}L{4.0cm}@{}}
\toprule
\textbf{Package} & \textbf{Source} & \textbf{Use in Paper~III-A} & \textbf{Status here} \\\\
\midrule
Typed vector/scalar role signatures, G0, and GP triptych
& Paper~II-A
& G0 actualization boundary; host types; GP1 finite/bounded, GP2 exact/faithful, GP3 irreducible/least
& \textsc{Import}; full structured-lift uniqueness remains subject to the source lift-stability scope \\\\
Derived momentum parent and structural seed
& Paper~II-A/II-B
& round $S^3_p(\Rmom)$ parent and $\mathcal H_p\cong\mathcal H_1$, $m_p=1$
& \textsc{Upstream theorem output relative to the II-A/II-B G0 host occurrence}; no independent occurrence witness is added here \\\\
Scalar succession role and orientation provenance
& Paper~I-A/Paper~II-A
& prior-generated scalar role, positive Hopf primitive, and typed provenance arrow used after the vector scalar-defect theorem
& \textsc{Import}; the sign is derived on the \(\mathsf{HS}_t\) domain, while units and Hamiltonian selection remain later \\\\
Clocked propagation content
& Paper~II-A/Paper~III-A
& \(\kappa_t^{-1}\), the typed Hopf inclusion in \(A_{\rm v}^{\rm op}\), and the faithful reciprocal isometry action
& \textsc{Derived}; one nonstationary seed and its least constructor doctrine replace an independent full-family occurrence premise \\\\
Canonical mathematical profile dual
& Paper~II-A, Theorem~\ref{p2:thm:canonical-metric-profile-separation}
& separating mathematical profile spectrum
& \textsc{Theorem import}; not yet an independent physical role \\\\
Intrinsic \(E_1\)-evaluation reciprocal closure
& Paper~II-C, Theorem~\ref{p2c:thm:coordinate-generated-reciprocal-role}
& Born/GP2-generated reciprocal role, relation-free product, and Born exchange without metric-curry
& \textsc{Upstream theorem}; discharges BR/IND statuses \\\\
Metric-curry closure $\mathsf P_{\rm mc}$
& Paper~II-A
& admits the distance profiles as physical reciprocal tests
& \textsc{Optional G0 test-occurrence certificate}; not a reciprocal-role premise \\\\
$(\mathrm{BR}_{\rm ind})+(\mathrm{IND}_{xp})$ theorem statuses
& Paper~II-A GP2 closure / Paper~II-C discharge / Paper~III-A full-profile replay
& certify generated reciprocal variance and absence of a generated joint-support constraint
& \textsc{Derived statuses}; not independent constitutive interfaces \\\\
$(\mathrm{CAL}_{xp})$
& current Paper~III calibration interface
& converts the transported reciprocal profile metric into physical configuration units, $\Rcs=\alpha_{xp}\Rmom$
& \textsc{G0 calibration}; witness not generated by Born reciprocity \\\\
$(\mathrm{CAL}_{\hbar})$
& current Paper~III calibration interface
& chooses the physical $2\pi\hbar$ representative of the derived primitive scalar phase class
& \textsc{G0 calibration}; not a topological selector \\\\
$(\mathrm{CAL}_{t})$
& current Paper~III calibration interface
& assigns a physical duration $T_t>0$ to one already derived positive primitive scalar turn
& \textsc{G0 calibration}; not an orientation or propagation selector \\\\
Position-side finite family $\mathsf K_{\rm res}(m_x,1)$
& Paper~II-C
& derived scalar comparison tower with momentum entry already fixed
& \textsc{Theorem-level family}; the occurring frontier is the residual G0 datum, derives $m_x=m_x^{\rm res}(F_x)$, and the primitive frontier gives $m_x=1$ \\\\
Hopf/Reeb abelian line
& Paper~I-A/II-D
& contact/phase direction used as a geometric export
& \textsc{Import}; no Standard-Model identification here \\\\
\bottomrule
\end{tabular}
\caption[Paper-II imports used by Paper~III-A]{Current Paper-II imports and interfaces used by Paper~III-A.}
\label{tab:p3-paperii-import-ledger}
\endgroup
\end{table}

\subparagraph{Selector-status classes.}
For the present paper the imported/generated role structure, contingent G0
occurrences, physical calibrations, and open downstream bridges are kept
separate.  The first class includes the momentum seed, reciprocal product,
scalar torus, primitive Poincar\'e class, first-symbol pairing, and generated
propagation doctrine.  The second includes actual frontier histories and
particular tests/processes/trajectories.  The third contains
$(\mathrm{CAL}_{xp})$, $(\mathrm{CAL}_{\hbar})$, and
$(\mathrm{CAL}_{t})$. The scalar-harmonic and finite unit-only reciprocal
maps are already generated upstream. The fourth contains stronger
observable-to-operator and continuum comparisons, Hamiltonian/state choices,
and physical field/light
speed identifications.  Mathematical availability in any class is not by
itself an occurrence or a selector for another class.

\subparagraph{Finite-role dependency boundary.}
Paper~II distinguishes the smooth round parent-role space from its finite
carrier.  The momentum carrier is already fixed upstream as
$\mathcal H_p\cong\mathcal H_1$ with derived label $m_p=1$.  Generic modules
$\mathcal H_m$ and paired notation $S^3_m$ remain auxiliary/later-role
mathematics unless a separate role signature constitutes them.  On the
reciprocal side, Paper~II-C derives the mathematical family
$\mathsf K_{\rm res}(m,1)$; an occurring frontier history supplies the
residual G0 datum, its least-resolution statistic $m_x$ is derived from that
frontier, and the primitive frontier already has $m_x=1$.  An occurring coordinate operation supplies its own finite
coordinate-class point through Paper~II-C's invariant-intensity theorem; no
second completion witness is required for that operation.  Paper~I-B classifies
the class space by
$\mathfrak M_{\rm coord}/\!\sim\cong\mathbb R_{\ge0}^3$, with
$\mathsf C_{\rm PI}$ one special representative; choosing that numerical
representative rather than the class carried by an occurring operation is a
calibration/specialization.  Paper~II-A derives finite metric support once
point-resolution readout occurs.  The exceptional $\mathsf{Supp}_{2T}$ output
is then co-realized only after a generated bridge to the Paper-I-B
exceptional-core constructor, or an actual application of that constructor;
point resolution alone does not supply this domain bridge.

Paper~II supplies the parent-role foundation, the structural momentum seed, the
scalar A-role, and the generated transformation--test doctrine summarized in
Table~\ref{tab:p3-paperii-import-ledger}.  In the current terminology GP1 is
Finite/Bounded, GP2 is Exact/Faithful, and GP3 is Irreducible/Least.

\begin{principle}[Imported stage-safe G0]
\label{p3:g0-actualization}
All physical interfaces below are role-specific specializations of
Paper~II-A's G0 geometric-actualization certificate.  Each asserted occurrence
must exhibit a contingent token of a geometrically licensed, claim-relative
invariant distinction with result-independent provenance.  Its role tag alone
asserts no occurrence and supplies no independent actualization principle.
No instance reopens frozen upstream geometry; a new role, standard, or
comparison interface enters only through a declared downstream dependency
frontier.  Analytic and reconstruction hypotheses remain separately displayed.
No independent momentum witness exists: the imported carrier is a generated
descendant of the occurring transfer-operation root.  A particular momentum
test, process, trajectory, bridge, or calibration remains a G0 token on its
own stated domain.  An unresolved realization family is retained until a
licensed discriminator exists; an independently constituted later frontier may
then discharge that family, but only prospectively and without retroactively
changing the earlier selector status.
\end{principle}

\begin{lemma}[Typed generated-role occurrence inheritance]
\label{p3a:lem:generated-role-occurrence-inheritance}
Let \(D=(d_1,\ldots,d_n)\) be a jointly occurring/generated root tuple in a
frozen presentation \(P\) with licensed grammar \(\Sigma\).  Let
\(\mathsf Q\) be a target sort and let \(w\) be a finite typed word satisfying
\[
 \operatorname{Orig}_{\Sigma}(w;D,\mathsf Q),
 \qquad q:=\operatorname{ReqOut}_{\Sigma}(w,D).
\]
Then
\[
 \operatorname{RootReal}_{P}(D)
 \quad\Longrightarrow\quad
 \operatorname{Obs}_{\rm FRC}(w,D)
 \ \vee\
 \operatorname{CoReal}_{\mathsf Q}(q\mid D;w).
\]
If one GP1-finite-and-bounded, GP2-exact-and-faithful realization of the
complete word is exhibited, the obstruction alternative is false and the
co-realization conclusion follows.
If \(q\) is a generated role space or operation doctrine, the conclusion is
only its co-realization with the complete root tuple; it does not assert
\(\operatorname{PhysOcc}(x)\) for every \(x\in q\).
\end{lemma}

\begin{proof}
This is the role-level specialization of
Paper~II-A,
Theorem~\ref{p2a:thm:rooted-generated-image-corealization}, together with
its multi-root/role-token
Corollary~\ref{p2a:cor:g0-witness-transport-multiroot-boundary}.  In
particular, intrinsicness or mathematical canonicality alone is not enough:
the licensed frozen word and its predeclared required output are indispensable.
Naturality, conservativity, and zero realization defect may verify the word's
faithful realization, but none of them licenses a missing constructor.
\end{proof}

\begin{definition}[Constitutive datum]
\label{p3:def:constitutive-datum}
A constitutive datum is any datum required to make a finite carrier into a physical sector.  Such data include preferential data---points, directions, scales, subgroups, quotients, projectors, weights, channel labels, profiles, and finite branch labels---and operational data---composition laws, covering structures, reciprocal sectors, symplectic/contact structures, scalar conjugate pairs, Hamiltonian generators, and spectral-action profiles.
\end{definition}

\begin{principle}[Imported GP1: finite and bounded]
\label{p3:gp1-finiteness}
Every realized structured stage has finite realized capacity, while every
constituted comparison/rate family is bounded or totally bounded in its typed
comparison structure.  GP1 does not select a spectral level or a dynamical
action.
\end{principle}

\begin{principle}[Imported GP2: exact and host-faithful]
\label{p3:gp2-geometric-primacy}
Every retained physical datum belongs to the constructor-relative generated
host signature: nothing ungenerated is inserted, nothing generated is omitted,
and generated object/operation/relation/test sorts are represented without
collapse or retyping.  A later joint closure may therefore neither erase a
generated factor distinction nor add an ungenerated cross-role relation.
\end{principle}

\begin{principle}[Imported GP3: irreducible and least]
\label{p3:gp3-faithful-closure}
Among GP1-admissible and GP2-exact/faithful complete structured realizations,
GP3 removes only host-faithful realization redundancy.  An endpoint is
irreducible when no proper admissible structured reduction remains and least
when every other admissible realization reduces to it.  If no unique least
class exists, the unresolved irreducible family is retained.
\end{principle}

\begin{theorem}[Imported reciprocal role and optional profile presentation]
\label{p3a:thm:v166-reciprocal-interface-summary}
For the constituted round vector parent \(A_{\rm v}=S_A^3\), Paper~II-C,
Theorem~\ref{p2c:thm:coordinate-generated-reciprocal-role}, supplies
\[
 B_{A_{\rm v}}^{(1)}=\iota_1(A_{\rm v}),\qquad
 A_{\rm v}\times B_{A_{\rm v}}^{(1)}\cong S_A^3\times S_B^3,
\]
with the transported round comparison metric and Born exchange. No
metric-curry occurrence is required for these outputs.

For a compact metric operation parent \(A\) on the frozen relation-free
domain of Definition~\ref{p3:def:reciprocal-interface-hypotheses}, assume
the optional \(\mathsf P_{\rm mc}\) premise and the retained isometric
operation domain stated there. The full profiles
\(b_z(x)=d(z,x)\) are additionally retained as physical tests in their
declared sort. They give the presentation
\[
 B_A^{\rm prof}=\{b_z:z\in A\},\qquad
 \kappa_A^{\rm prof}(z)=b_z,\qquad
 d_{B_A^{\rm prof}}(b_z,b_w)=\|b_z-b_w\|_\infty=d(z,w).
\]
Its total evaluation \(E(x,b_z)=d(x,z)\) and comparison isomorphism do
not impose a graph support relation. On this domain the full-profile
BR/IND statuses and Born involution follow as in
Theorems~\ref{p3:thm:reciprocal-status-discharge} and
\ref{p3:def:intrinsic-born-reciprocal-constitution}.
The associated function algebras, completions and spectra are mathematical
envelopes. Paper~III-G, Appendix~\ref{p3a:app:v166-full-profile} gives the optional
construction and proof; it adds a profile-test presentation, not a premise
for the native reciprocal role.
\end{theorem}
\begin{proof}
The intrinsic evaluation theorem supplies the first assertion. The distance
inequalities give \(\|b_z-b_w\|_\infty\le d(z,w)\), with equality by
evaluation at \(z\), proving the optional isometry. Its physical test
status uses exactly \(\mathsf P_{\rm mc}\). The full-profile statements
and their unchanged relation-free domain are proved in the cited appendix.
\end{proof}



\begin{definition}[Constitutive residue]
\label{p3:def:constitutive-residue-main}
A \emph{constitutive residue} is a datum or interface required by a later
physical layer but not yet intrinsic to, generated by, or fixed by the current
constituted signature.  It may not be silently inserted; a later theorem must
either derive it or enlarge the displayed domain explicitly.  When the residue
is an occurrence, test, preparation, process, trajectory, bridge, readout, or
calibration, G0 supplies its common admissibility type but not its witness.
The detailed
bookkeeping lemmas are deferred to Paper~III-G, Appendix~\ref{p3a:app:framework-bookkeeping}.
\end{definition}

\begin{definition}[Imported vector-momentum constitution]
\label{p3:def:imported-vector-momentum-constitution}
Paper~II supplies the derived round momentum parent-role space
$S^3_p(\Rmom)\cong SU(2)$ and, separately, the exact finite GP2-generated momentum \emph{type-support}, represented
by the canonical one-copy model
\[
   \boxed{\mathcal G^{\rm car}_{\Sigma_p}\cong\mathcal H_1,
          \qquad m_p=1.}
\]
The descriptor $m_p=1$ is the frozen word/level label, not the
multiplicity of an arbitrary physical structured endpoint.  The full
carrier identification $\mathcal H_p\cong\mathcal H_1$ additionally
requires the whole-structured comparison hypotheses of
Theorem~\ref{p2b:thm:momentum-carrier-rigidity}.  Operator formulas on
$\mathcal H_1$ concern its canonical finite model until such a full
endpoint lift is certified.
The generic modules
$\mathcal H_m=\bigoplus_{j_L+j_R\le m}V_{j_L,j_R}$ and the notation $S^3_m$
remain auxiliary representation/tower data unless a later role signature
constitutes them.  This paper does not re-select the momentum carrier.
\end{definition}

\subparagraph{Auxiliary momentum-selector countermodels.}
The word-stage/lower-ideal and fixed-coordinate closure constructions do not
derive the momentum seed in this constitution.  Their sharp
countermodels and conditional selector theorem are retained only in
Paper~III-G, Appendix~\ref{p3a:app:legacy-selector}.

\begin{lemma}[Generated rank-one lattice-dual spectrum]
\label{p3:lem:rank-one-lattice-dual-spectrum}
Let \(A_{\rm s}=V/\Lambda\) with \(\dim V=1\) and
\(\Lambda\cong\mathbb Z\).  The deck/loop lattice \(\Lambda\) is intrinsic to
the universal-cover presentation of \(A_{\rm s}\).  Its commutative group
\(C^*\)-algebra has canonical Gelfand spectrum
\[
        C^*(\Lambda)
        \cong
        C\!\left(\widehat\Lambda\right),
        \qquad
        \widehat\Lambda
        :=
        \operatorname{Hom}(\Lambda,U(1))
        \cong V^*/\Lambda^*,
\]
where
\[
        \Lambda^*
        :=
        \{\ell\in V^*:\ell(\Lambda)\subseteq\mathbb Z\}.
\]
The Gelfand transforms
\[
        u_\lambda(\chi):=\chi(\lambda),
        \qquad \lambda\in\Lambda,
\]
separate \(\widehat\Lambda\) and generate a uniformly dense unital
self-adjoint algebra in \(C(\widehat\Lambda)\).

Moreover, in rank one there is a canonical lattice-dual group isomorphism
\[
        \delta_\Lambda:V/\Lambda\longrightarrow V^*/\Lambda^*.
\]
Consequently the phase-character spectrum is a canonical mathematical model
of the reciprocal B-sort, and the opposite pullback law transports through
\(\delta_\Lambda\).  Its physical target typing and retention as a comparison
morphism are supplied by the loop-character evaluation constructor below, not
by \(C^*\)-algebra or Gelfand-spectrum formation alone.
\end{lemma}

\begin{proof}
Pontryagin--Gelfand duality for the discrete abelian group \(\Lambda\)
gives
\[
        C^*(\Lambda)\cong C(\widehat\Lambda).
\]
The canonical group elements of \(C^*(\Lambda)\) have Gelfand transforms
\(u_\lambda\).  Distinct characters differ on some
\(\lambda\in\Lambda\), so these functions separate points; their finite
linear span is a unital self-adjoint algebra and is uniformly dense by
Stone--Weierstrass.

Choose either primitive generator \(e\) of \(\Lambda\), and let
\(e^*\in\Lambda^*\) satisfy \(e^*(e)=1\).  Define
\[
        \delta_\Lambda([te]):=[te^*].
\]
Changing \(t\) by an integer changes \(te^*\) by an element of
\(\Lambda^*\), so the map is well defined.  Replacing \(e\) by \(-e\)
also replaces \(e^*\) by \(-e^*\) and \(t\) by \(-t\), leaving
\([te^*]\) unchanged.  Hence the map is independent of the only primitive
generator ambiguity and is a canonical continuous group isomorphism.
The map
\[
        \delta_\Lambda:
        A_{\rm s}\longrightarrow\widehat\Lambda
\]
identifies the two compact rank-one group presentations mathematically and
transports the opposite law.
\end{proof}

\begin{definition}[Typed loop-character evaluation constructor]
\label{p3a:def:scalar-loop-character-constructor}
For the constituted scalar role \(A_{\rm s}=V/\Lambda\), let
\(\mathsf{Loop}(A_{\rm s})=\Lambda\) be its generated deck/iteration relation
sort.  Before any character output is evaluated, freeze in the scalar-loop
grammar the character-evaluation schema, declare the reciprocal target sort
\(\mathsf{Test}_{B_{\rm s}}\) and define
\[
 \mathsf{CHAR}_{\rm sc}(\Lambda)
 :=
 \left(
 B_{\rm s}:=\operatorname{Hom}(\Lambda,U(1)),
 \ \delta_\Lambda,\ \delta_\Lambda^{-1},
 \ \{u_\lambda\}_{\lambda\in\Lambda},
 \ \operatorname{ev}_{\rm sc}
 \right),
\]
where
\[
 u_\lambda(\chi):=\chi(\lambda),
 \qquad
 \operatorname{ev}_{\rm sc}(\lambda,\chi):=\chi(\lambda).
\]
Using the rank-one lattice-dual construction of
Lemma~\ref{p3:lem:rank-one-lattice-dual-spectrum}, the constructor sends the
full loop comparison profile
\(\lambda\mapsto(\chi\mapsto\chi(\lambda))\) to its declared separating
B-test output.  Its physical output is this evaluation-test role and its
faithful opposite operation law, together with the canonical typed comparison
isomorphism \(\delta_\Lambda:A_{\rm s}\to B_{\rm s}\) and its inverse.  The
algebra \(C^*(\Lambda)\), its completion, and its Gelfand spectrum are
mathematical envelopes used to prove compactness and completeness; they are
not physical output sorts of \(\mathsf{CHAR}_{\rm sc}\).
Writing \(w_{\rm CHAR}\) for the displayed finite word from the retained loop
relation through \(\operatorname{Hom}(-,U(1))\), the rank-one comparison,
evaluation, and the declared regular opposite-operation closure, every
intermediate is retained in its native sort and
\[
 \operatorname{Orig}_{\Sigma_{\rm sc}^{\rm loop}}
 (w_{\rm CHAR};(A_{\rm s},\Lambda),\mathsf{Test}_{B_{\rm s}})
\]
holds by the frozen schema.  This records the exact model-specific constructor
input; it is not inferred from naturality or added by GP1--GP3.
\end{definition}

\begin{theorem}[G0-derived physical scalar reciprocal role]
\label{p3a:thm:g0-derived-scalar-character-role}
On the constituted scalar-loop domain, the certified word
\(w_{\rm CHAR}\) is nonobstructed and Paper~II-A,
Theorem~\ref{p2a:thm:rooted-generated-image-corealization}, co-realizes the
physical evaluation-test role
\[
 B_{\rm s}=\widehat\Lambda\cong V^*/\Lambda^*\cong U(1)
\]
with the scalar A-role without a second occurrence premise.  Its regular
opposite action is faithful, \(\delta_\Lambda:A_{\rm s}\to B_{\rm s}\) is
the generated comparison isomorphism. Evaluation is total on
\(\Lambda\times B_{\rm s}\); it is not a pairing of two torus points.
The frozen scalar-role relation slot contains no joint-support equation on
\(A_{\rm s}\times B_{\rm s}\).
Consequently Paper~II-A,
Theorem~\ref{p2:thm:gp2-generated-reciprocal-role-closure}, gives the physical
relation-free scalar product
\[
 \boxed{A_{\rm s}\times B_{\rm s}\cong S^1_t\times S^1_E.}
\]
The same root co-realizes the scalar Born exchange
\[
 \mathcal B_{\rm s}(a,b)
 :=\bigl(\delta_\Lambda^{-1}b,\delta_\Lambda a\bigr),
 \qquad
 \mathcal B_{\rm s}^2=\operatorname{id}.
\]
\end{theorem}

\begin{proof}
The map
\[
 \Lambda\longrightarrow U(1)^{B_{\rm s}},
 \qquad \lambda\longmapsto u_\lambda,
\]
is natural under lattice isomorphisms, contains no choice of primitive
generator or scale, and is injective because characters separate points of the
discrete abelian group \(\Lambda\).  It therefore transports every actual
loop/iteration distinction to a distinction in the declared B-test sort and
reflects equality.  The output is exactly the essential image
\(\{u_\lambda\}\), its compact character parameter role, the canonical
rank-one comparison isomorphism, and the explicitly declared regular
operation closure; no additional test, calibration, or support relation is
inserted.  The standard product model
\((V/\Lambda)\times\operatorname{Hom}(\Lambda,U(1))\), equipped with the
retained loop evaluation on \(\Lambda\times B_{\rm s}\), is
GP1-admissible and has zero joint insertion, omission, relation, type, and
closure defect. The A-role point is independent of the loop-class argument;
no torus-point evaluation is introduced.  It therefore rules out
\(\operatorname{Obs}_{\rm FRC}(w_{\rm CHAR},(A_{\rm s},\Lambda))\), and the
rooted generated-image theorem co-realizes the typed required output.  The
preceding naturality and separation checks establish faithful realization;
they are not constructor-license premises.

Pointwise multiplication makes \(B_{\rm s}\) a compact rank-one group.
Through the canonical rank-one isomorphism \(\delta_\Lambda\), translation
induces a faithful opposite action.  The two distinct generated role sorts and their comparison maps are retained,
and no generated relation sets \(b=\delta_\Lambda(a)\) or imposes another
joint-support condition. A comparison isomorphism is not such an equation.
The relation-free-composition antecedents of the GP2 reciprocal-role theorem
are therefore satisfied,
which proves objectification and relation-free composition.  The
Born exchange and its involution now follow from
Corollary~\ref{p2a:cor:born-invertible-corealization}.  The
Pontryagin--Gelfand and Stone--Weierstrass statements in
Lemma~\ref{p3:lem:rank-one-lattice-dual-spectrum} certify mathematical
compactness and completeness only and are not used to cross the physical
output boundary.
\end{proof}

\begin{corollary}[The scalar \(S^1\) pair is one invertibly bridged co-realization]
\label{p3a:cor:s1-universal-corealization-bridge}
The roles \(S^1_t\) and \(S^1_E\), their loop-character evaluation,
opposite operations, relation-free product, and scalar Born exchange form one
rooted incidence component.  Their ultimate origin is the occurring scalar
role together with its intrinsic loop lattice \(\Lambda\); their immediate
constructors are respectively loop extraction, character dualization,
evaluation dualization, product formation, and invertible exchange.
Accordingly no second occurrence premise is required for the reciprocal
circle, its dual operations or dynamics, or the Born involution.

The conclusion is uncalibrated.  It supplies neither
\(\tper\Eper=2\pi\hbar\), a time orientation, a Hamiltonian, an actual
evolution, nor an energy/history token.
\end{corollary}

\begin{proof}
Apply Theorem~\ref{p2a:thm:universal-corealization-bridge-normal-form} to
\(w_{\rm CHAR}\).  Its evaluation-dual operation is covered by
Theorem~\ref{p2a:thm:evaluation-dual-corealization}; its generated comparison
isomorphism and relation-free product place it on the Born specialization.
The excluded outputs occur only in later calibrated or dynamical sorts.
\end{proof}

\subparagraph{Optional full-profile enhancement.}
The complete distance-profile doctrine supplies a faithful presentation of
the already constituted reciprocal role on the additional
\(\mathsf P_{\rm mc}\) domain. Its operation dictionary and mathematical
envelope reconstruction are recorded in
Paper~III-G, Appendix~\ref{p3a:app:v166-full-profile}; the native vector and scalar
products use the intrinsic evaluation and loop-character constructors.

\begin{remark}[Relation-free support versus conservative derived readouts]
\label{p3a:rem:relation-free-readout-firewall}
The exclusion of a graph condition or nontrivial quotient in
Theorem~\ref{p3:def:intrinsic-born-reciprocal-constitution} forbids replacing the
physical product, deleting factor tests, or imposing an ungenerated equality
on its support.  It does not forbid a total licensed readout
\(F:A\times B_A\to Q\) whose graph lift
\[
 c\longmapsto(c,F(c))
\]
is conservative on the complete source and whose target sort is declared
before evaluation.  Such an extension is derived only after the full
multi-root tuple satisfies Paper~II-A's witness-transport clauses and has zero
joint GP2 defect.  Equality \(F(c)=F(c')\) then belongs to the output
comparison sort and does not assert \(c=c'\) in the source.

This is also the simultaneous-realizability firewall.  Memberwise admissibility
of two families does not prove that an arbitrary pair of members is jointly
realizable: only tuples in their complete joint zero-defect locus survive.
Conversely, when a total graph-lift constructor has zero defect on the whole
product, all source pairs remain and the new readout is co-realized without a
selector or duplicate occurrence root.
\end{remark}

\begin{remark}[Local Lorentzian form and global gravitational scope]
\label{p3:rem:born-reciprocity-lorentzian-recovery}
The local operation cone derived below determines a Lorentzian conformal
form on its rank-four domain, with representative
\(-c_{\rm op}^2d\widetilde t^{\,2}+h\). This does not by itself
supply a global spacetime realization, a gravitational metric action or the
dynamics of a particular field. Those conclusions retain their generated
coframe, owner and action interfaces. The conditional coframe descent in
Subsection~\ref{p3a:sec:v189-local-symmetric-quadratic-descent} makes this
distinction explicit.
\end{remark}





\begin{theorem}[GP2-derived Born completion of the Paper-II scalar branch]
\label{p3:thm:scalar_classification_import}
Paper~II classifies the least nontrivial compact connected abelian scalar
operation parent as
\[
        A_{\rm s}=V/\Lambda\cong U(1)\cong S^1_A,
        \qquad \dim V=1,
\]
with primitive lattice \(\Lambda\cong\ZZ\).  Under the scalar instance of
the premise-minimal evaluation/lattice-dual closure, the GP2-objectified
reciprocal role and its pullback action make the primitive phase characters
constitute the
full phase-character B-torus of that lattice,
\[
    B_{\rm s}=\operatorname{Hom}(\Lambda,U(1))
    \cong V^*/\Lambda^*,
        \qquad
        \Lambda^*:=\{\ell\in V^*: \ell(\Lambda)\subseteq\ZZ\},
\]
and hence
\[
        B_{\rm s}\cong U(1)\cong S^1_B.
\]
The discrete reciprocal mode lattice and the compact character role are
therefore different typed objects:
\[
   \Lambda_E:=\Lambda^*\cong\mathbb Z,
   \qquad
   S^1_E:=B_{\rm s}=V^*/\Lambda_E.
\]
Elements of \(\Lambda_E\) label character generators; points of \(S^1_E\)
are characters.  No identification of the discrete labels with the compact
role space is used.
Its pullback law is contravariant, and chronological persistence gives, by
Paper~II-A's GP2 reciprocal-role theorem, the joint scalar role space
\[
        C_{\rm s}=A_{\rm s}\times B_{\rm s}
        \cong S^1_A\times S^1_B=T^2.
\]
Thus the scalar torus is the GP2 relation-free product of the classified
scalar transformation role and its generated reciprocal phase
role.  These are distinct typed roles, not two new geometric principles.
\end{theorem}

\begin{proof}
The Paper-II scalar succession role has a one-dimensional compact connected
abelian parent, hence the displayed quotient \(V/\Lambda\), and its primitive
deck/iteration relation sort is \(\Lambda\).  The typed loop-character
constructor of Definition~\ref{p3a:def:scalar-loop-character-constructor}
transports that complete relation profile into the declared reciprocal-test
sort.  Theorem~\ref{p3a:thm:g0-derived-scalar-character-role} proves, without
a second occurrence premise, that the resulting character-test role is
co-realized, carries a faithful opposite action, and has total joint
evaluation.  Paper~II-A's GP2 closure theorem therefore supplies the
Born/GP2-generated physical B-role and the relation-free joint product.

The mathematical compact model
\(\operatorname{Hom}(\Lambda,U(1))\cong V^*/\Lambda^*\)
and the completeness of its phase functions are independently verified
by Lemma~\ref{p3:lem:rank-one-lattice-dual-spectrum}.  Its \(C^*\)-algebra and spectrum are not used as a
physical-output constructor.  Thus no metric-curry premise or physical
algebra/spectrum premise is used.
\end{proof}




\begin{definition}[Reciprocal-vector bookkeeping package]
\label{p3:def:reciprocal-vector-bookkeeping-package}
We retain the shorthand
\[
\begin{aligned}
\mathcal A_{\mathrm{book}}:={}&
\text{Paper-II vector/scalar classification}\\
&+\text{the Paper-III GP2-derived master reciprocal closure}.
\end{aligned}
\]
only for compact theorem statements.  It is not an additional GP or a
foundational constitution package.  It abbreviates the opposite test action
on the intrinsic evaluation spectrum together with the GP2-derived reciprocal
role and relation-free product; \(\mathsf P_{\rm mc}\) adds only the optional
full distance-profile test doctrine.  The primitive scalar phase class is derived
below inside that domain;
later action, state, stationarity, and polarization claims retain the theorem
domains stated in Papers~III-B--III-F.
\end{definition}


\noindent\textbf{Proof convention.}
Every native closure theorem below states its imported role datum, its
Born/contact or spectral admissibility class, the alternatives excluded by
no-insertion, and the surviving intrinsic datum.  The words ``unique'', ``least'', and ``forced'' are always relative to that explicitly stated constitution domain.  The word ``minimal'' is reserved for the technically weaker preorder notion.


\begin{remark}[Dimensionful calibration and metrological boundary]
\label{p3:rem:dimensionful-calibration-boundary}
The closure theorems below determine role type, reciprocal pairing,
primitive integer classes, dimensionless normalization, and relative
coefficients.  Their physical presentation uses the reciprocal scale
certificate \(\alpha_{xp}\), the primitive-turn duration \(T_t\), and the
physical representative \(\hBorn=2\pi\hbar\) of the primitive Born scalar
action cell.  The operational rate
\[
 c_{\rm op}=\frac{2\pi R_x}{T_t}
\]
is then derived from the represented spatial radius and primitive-turn clock.
The factor \(\alpha_{xp}\) converts the transported reciprocal profile metric
to configuration-space units; it is not fixed by Born reciprocity.  Likewise
\((\mathrm{CAL}_{\hbar})\) represents an already derived primitive integral
class and \((\mathrm{CAL}_t)\) represents an already derived primitive turn.

Newton's constant is not an additional upstream calibration input in this
paper.  Paper~III-F later derives a joint Einstein--Hilbert coefficient and
may attach to it a Newton coupling coordinate.  Comparing that derived
coordinate with a realized numerical standard is metrology/validation; it does
not create the geometric coefficient.  A global simultaneous rescaling of the
chosen dimensional standards therefore does not reopen any dimensionless
structural theorem.

No theorem here proves that a particular later field or photon channel has the
undeformed clocked-host principal owner.  Paper~II-A proves the conditional
inheritance theorem: after that ancestry is established, its characteristic
front speed equals \(c_{\rm op}\).  None of the displayed metrological
representatives selects \(S^3\), \(S^1\), Born doubling, the scalar torus,
the primitive integral class, \(m_p=1\), or a dimensionless cutoff profile.

The full spectral action retains one common regulator profile unless it is
specialized to the primitive heat semigroup; the semigroup profile itself is
then exponential by path-independent composition.  The cosmological
realization of the position-side unfolding \(m_x(\widetilde t)\) is a solution
problem assigned to Paper~VI-A and is not a Paper-III constitutive premise.
\end{remark}

\subsection{Exhaustive readout-role handoff to Paper~III}
\label{p3a:sec:readout-role-handoff}

\begin{definition}[Dynamical occurrence tags]
\label{p3a:def:dynamical-readout-occurrences}
For a finite process diagram \(W\) in the frozen Paper-II/III constructor
grammar, write
\[
\begin{aligned}
 (\mathrm{LOOP}_{\rm occ})[W:X\to X]
 &:\Longleftrightarrow
 \operatorname{FSR}_{\Sigma}(\operatorname{Tr}_X W)\\
 &\quad\text{and the word is the canonical unmarked}\\
 &\quad\text{duality closure of }W,\\[0.35em]
 (\mathrm{BOUND}_{\rm occ})[\epsilon,T,\psi]
 &:\Longleftrightarrow
 \operatorname{FSR}_{\Sigma}(\epsilon T\psi)\\
 &\quad\text{and the displayed boundary data}\\
 &\quad\text{have generated provenance}.
\end{aligned}
\]
These are G0 occurrence predicates.  They do not select \(W,T,\psi\), or
\(\epsilon\), and they do not identify two occurrences merely because their
numerical values agree.
\end{definition}

\begin{theorem}[Two-branch dynamical handoff]
\label{p3a:thm:two-branch-dynamical-handoff}
Every primitive finite structural scalar readout used in the later dynamical
papers has exactly one of the forms
\[
\boxed{
\begin{aligned}
 (\mathrm{LOOP}_{\rm occ})[W:X\to X]
 &\Longrightarrow S_{\circ}(W)=\operatorname{Tr}_XW,
 &&S_{\circ}\in\Rloop,\\
 (\mathrm{BOUND}_{\rm occ})[\epsilon,T,\psi]
 &\Longrightarrow S_{\partial}(\epsilon,T,\psi)=\epsilon T\psi,
 &&S_{\partial}\in\Rbdry.
\end{aligned}}
\]
Disconnected products, sums, determinants, Pfaffians, entropies, and nonlinear
response functions lie in the generated closure \(\Rgen\); they do not define
a third primitive outer role.  In particular, an unmarked closed spectral
action is a loop occurrence, while prepared/tested quantum processes and
external matter fields are boundary-conditioned occurrences.
\end{theorem}

\begin{proof}
The frozen Paper-III diagrams satisfy the finite structural-readout predicate
precisely when their constructors and retained marks have the displayed
provenance.  Apply Paper~II-A
Theorem~\ref{p2:thm:loop-boundary-readout-exhaustion}.  Empty external
boundary gives the loop form. For the specified canonical closing word,
Lemma~\ref{p2:lem:canonical-finite-duality-trace} and
Theorem~\ref{p2:thm:unmarked-natural-closure} give the ordinary trace;
nonempty boundary gives the
matrix-coefficient normal form.  Paper~II-A
Corollary~\ref{p2:cor:generated-readout-closure} classifies every finite
decoration or disconnected composite.  The alternatives are exhaustive and
provenance-disjoint.
\end{proof}

\begin{corollary}[No readout profile from GP3]
\label{p3a:cor:no-gp3-readout-profile}
GP3 compares role-faithful carrier realizations only after the role has been
fixed.  It neither converts \(\Rbdry\) into \(\Rloop\), nor selects a spectral
process law inside \(\Rloop\), nor deletes a generated state, effect, source,
field, grading, or projector.  The later closed spectral class \(\Dcl\) must
therefore be derived from an actual loop occurrence together with the
positive Dirac-even operator input; its particular process member remains G0
occurrence data.
\end{corollary}


\subsubsection{Scalar-circle role dictionary and calibrated spatial clock transport}
\label{p3a:sec:v180-scalar-circle-clock-transport}

\begin{remark}[Scalar evaluation, carrier, and window dictionary]
\label{p3a:rem:scalar-evaluation-carrier-window-dictionary}
The scalar-loop constructor retains the exact characters
$u_n(b)=b(ne_t)$ with $u_{n+m}=u_nu_m$ and $u_{-n}=\overline{u_n}$.
A dual-lattice label, a point of the compact reciprocal circle, a Fourier
character, and a physical Hamiltonian eigenvalue remain distinct sorts.
Likewise the fixed compact momentum level $m_p=1$, the growing position
resolution $m_x$, the position comparison degree $m_{\rm cmp}=2$, and the
independent finite scalar Fourier window of Paper~III-B are not interchangeable
cutoffs.  Born exchange transports a supplied finite window; it does not
select its size.
\end{remark}

\begin{corollary}[Clocked position character support]
\label{p3a:cor:clocked-position-character-support}
The already generated clocked-Hopf co-realization pulls the primitive Hopf
action back to every complete occurring position band.  On degree $k$,
\[
 E_k^x|_{S_t^1}=\bigoplus_{q=-k,-k+2,\ldots,k}(k+1)\chi_q,
\]
so on $\mathcal P_x(m_x)$ the represented clock support is
$\{-m_x,\ldots,m_x\}$.  One primitive half-turn acts by $(-1)^k$ on
$E_k^x$; enlarging the position harmonic support does not multiply the
primitive winding of the clock circle.
\end{corollary}

\begin{corollary}[Calibrated clock family equals the scalar kinetic operator]
\label{p3a:cor:clock-family-scalar-kinetic-identity}
On a compatible calibrated datum with
$E_\star a=\hbar c_{\rm op}$, define $G_i=E_\star Q_i$.  Then the geometric
Casimir identity of Paper~II-D gives
\[
 \boxed{\sum_iG_i^2=\hbar^2c_{\rm op}^2\Delta_a.}
\]
The equality uses the same represented chiral family and the existing
calibration.  A directional $G_i$ is not the total Hamiltonian, and a family
of compatible calibrated data is not an additional physical time history.
\end{corollary}


\paragraph{Technical proof ownership and constituted scope.}
Every extended auxiliary derivation, explicit matrix verification,
and optional comparison moved during this reorganization is
preserved in Paper~III-G, Section~\ref{p3g:sec:constitution}.
The physical momentum carrier is inherited from Paper~II-B:
$\mathcal H_p\cong\mathcal H_1$ and $m_p=1$, without later
scalar, action, regulator, or source data reselecting it.

\section{Generated reciprocal phase-space geometry}
\label{paper3a:sec:organized-reciprocal}
\subsection{Derivation of Compact Reciprocal Geometry}
\label{p3:sec:phasespace}

The compact reciprocal-geometry derivation is organized by Table~\ref{tab:p3-born-completion-ledger}.  The imported momentum row is recalled first; all subsequent rows are native Paper-III closure results.

\begin{table}[H]
\centering
\small
\begingroup
\setlength{\tabcolsep}{3pt}
\renewcommand{\arraystretch}{1.15}
\begin{tabular}{@{}
  L{3.5cm}
  L{8.0cm}
  L{\dimexpr\linewidth-11.5cm-4\tabcolsep\relax}
@{}}
\toprule
\textbf{Ledger row}
& \textbf{Claim}
& \textbf{Discharged by} \\
\midrule
Imported momentum sector
& The compact momentum entry is the Paper-II physical seed \(\mathcal H_p\cong\mathcal H_1\); generic \(S^3_m\) notation is auxiliary.
& Theorem~\ref{p3:thm:momentum} \\

Born reciprocal vector sector
& Metric-curry closure supplies physical profile tests; Paper~III-A proves the
BR/IND certification statuses, and Paper~II-A GP2 objectification/relation-free
composition supplies the reciprocal role, joint product, and Born involution.
& Lemma~\ref{p3:lem:born}, Theorem~\ref{p3:thm:position}, Corollary~\ref{p3:cor:doubled-vector} \\

Vector-sector symplectic defect
& The doubled vector parent-role space has no intrinsic compact symplectic phase form.
& Theorem~\ref{p3:thm:vector-sector-not-symplectic} \\

No hidden scalar pair
& The scalar pair is not already contained in the compact vector product.
& Lemma~\ref{p3:lem:no-intrinsic-vector-scalar}, Corollary~\ref{p3:cor:no-intrinsic-scalar-pair} \\

Scalar succession orientation
& Paper~I-A supplies the positive Reeb primitive, Paper~II-A transports it
through \(\mathsf{HS}_t\), and the full dual fixes the reciprocal energy sign.
& Theorem~\ref{p3:thm:transported-scalar-orientation} \\

Scalar Born-completion input
& The scalar transformation torus and its full test-dual carry the canonical primitive Poincar\'e phase class.
& Theorem~\ref{p3:thm:scalar-born-residue} \\

Primitive scalar pair
& The Poincar\'e evaluation form on the GP2-derived scalar product gives one canonical plane.
& Theorem~\ref{p3:thm:minimal-symplectic-completion}, Corollary~\ref{p3:cor:missing-sector-one-pair} \\

Compact scalar phase torus
& The full dual lattice reconstructs the reciprocal test spectrum; its
opposite variance and relation-free status trigger GP2 closure to the scalar
\(T^2\), whose primitive phase class is then derived.
& Lemma~\ref{p3:lem:minimal-phase-torus-quantization}, Theorem~\ref{p3:thm:scalars} \\

No scalar modulus
& The scalar pair carries no independent running modulus.
& Theorem~\ref{p3:thm:scalar-pair-no-modulus} \\
\bottomrule
\end{tabular}
\caption[Born reciprocal completion ledger]{Born reciprocal completion ledger.}
\label{tab:p3-born-completion-ledger}
\TableNotes{The following subsections prove the native rows in order.  Paper-II supplies the first-row vector parent-role space, the scalar A-branch classification, the intrinsic evaluation reciprocal closure, the mathematical profile doctrine, and the GP2 reciprocal-role theorem listed in Table~\ref{tab:p3-paperii-import-ledger}.  Paper~III-A imports the premise-minimal reciprocal product and, under \(\mathsf P_{\rm mc}\), separately realizes the full-profile physical test action and replays the BR/IND certification.}
\endgroup
\end{table}


\subsubsection{Imported momentum-sector result from Paper~II}

\begin{remark}[Reciprocal \(S^3\) notation and physical scale]
The symbol \(S^3_m\) denotes an auxiliary dimensionless \(\Spin(4)\)-equivariant
finite-stage constitution with label \(m\).  Born reciprocity first gives an
abstract round B-role \(S^3_B\) whose transported profile metric is isometric to
the A-role metric.  It does not by itself choose a new physical unit.  Only under
\((\mathrm{CAL}_{xp})\) do we write
\[
        S^3_p(\Rmom)\qquad\text{and}\qquad S^3_x(\Rcs),
        \qquad \Rcs=\alpha_{xp}\Rmom,
\]
with \(d_x=\alpha_{xp}d_{B_A}\).  The later finite position label \(m_x\) remains
a separate position-side interface and need not equal the already derived
momentum label \(m_p=1\).
\end{remark}

\begin{theorem}[Imported momentum parent and generated finite seed]
\label{p3:thm:momentum}
The current Paper-II chain supplies
\[
  \mathcal P_{\rm vec}^{\rm parent}\cong S^3_p(\Rmom),
  \qquad \mathcal G^{\rm car}_{\Sigma_p}\cong\mathcal H_1,
  \qquad m_p=1.
\]
The first is the round parent; the second is its generated GP2 type-support,
with $\mathcal H_1$ the canonical one-copy representative.  The level
label $m_p=1$ does not assert full-endpoint multiplicity-one.  Identifying
the entire GP3-irreducible physical carrier with $\mathcal H_1$ requires
the additional whole-structured comparison certificate of
Theorem~\ref{p2b:thm:momentum-carrier-rigidity}.
Generic $\mathcal H_m$ stages are auxiliary unless separately constituted.
\end{theorem}
\begin{proof}
The parent is imported from Paper~II-A and the exact six-type support from
Paper~II-B, Theorem~\ref{p2:thm:momentum-host-H1}.  The label $m_p=1$
follows from Corollary~\ref{p2:thm:m1-component} rather than a multiplicity
selector.  Corollary~\ref{p2b:cor:physical-momentum-carrier-corealization}
separates generated support co-realization from conditional endpoint rigidity.
No Paper~III premise enters these upstream statements.
\end{proof}
\begin{remark}[The imported branch is already the constituted simply connected one]
The content imported from Paper~II is stronger than the bare statement
``compact momentum parent is a spherical space form''.  Inside the typed
Paper-II role domain, the simply connected spin-compatible transfer parent
$S^3_A$ is already constituted, and Paper~II-B separately fixes the finite
generated momentum support $\mathcal G^{\rm car}_{\Sigma_p}\cong\mathcal H_1$.  Paper~III therefore does
not choose the momentum-side $S^3$ or its finite seed anew.  Under
$\mathsf P_{\rm mc}$, Paper~II's canonical metric-profile spectrum becomes a
full physical localization-test doctrine.  Independently, Paper~II-C's
intrinsic evaluation spectrum has already discharged the BR/IND statuses and
instantiated the reciprocal role and joint product through GP2.  The
free-quotient spectral
normalization guardrail is supplementary and is deferred to
Paper~III-G, Appendix~\ref{p3a:app:quotient-spectral-guardrail}.
\end{remark}

\begin{theorem}[GP2-derived vector Born completion]
\label{p3:thm:born-reciprocity-exhaustion}
Let \(A=S^3_A\cong SU(2)\) be the constituted Paper-II vector parent-role space.
The Paper-II-C intrinsic evaluation spectrum, the GP2-objectified reciprocal
role, and Paper~II-A's reciprocal closure give
\[
        B_{\rm v}\cong A^{\mathrm{op}}\cong A\cong S^3_B
\]
and chronological persistence gives, up to Born-equivariant role-space isomorphism,
\[
        C_{\rm v}
        =A\times B_{\rm v}
        \cong S^3_A\times S^3_B,
        \qquad
        \mathcal B_{\rm v}(a_1,a_2)=(a_2,a_1).
\]
By Paper~II, Theorem~\ref{p2:thm:intrinsic-coordinate-sector}, the
independently tagged A- and B-role doctrines each retain the intrinsic lowest
nontrivial coordinate type
\[
        V_{\frac12,\frac12}.
\]
Thus the reciprocal completion transports not only the abstract group type
but also the requirement that this vector sector remain admissible on each
side.  It does not impose a common finite level.  The momentum entry is already fixed to \(m_p=1\), while any finite reciprocal-position resolution entry remains a separate position-side interface.
This is a
\(\Sigma_{\rm II}^{\rm role}+\mathrm{G0}+\mathrm{GP1--GP3}\)
reciprocal-closure conclusion, not a choice of reciprocal
physical units, orientation, action, state, or infrared unfolding.
Under the optional \(\mathsf P_{\rm mc}\) certificate, the full
distance-profile family is additionally realized as the physical test
doctrine.
\end{theorem}

\begin{proof}
Paper~II-C's intrinsic \(E_1\)-evaluation spectrum generates the tagged
localization role, proves its faithful opposite action and the absence of a
joint-support relation, and invokes Paper~II-A's GP2 theorem to objectify it as
the autonomous physical role \(SU(2)^{\mathrm{op}}\).  Inversion
\[
        g\longmapsto g^{-1}
\]
is a canonical group isomorphism
\(SU(2)\to SU(2)^{\mathrm{op}}\), so the B-side has the same constituted round
three-sphere type and preserves the exact/faithful source form.  The product
and involution then follow from the GP2-derived reciprocal closure.
Because the intrinsic coordinate eigenspace, its evaluation tests, and their
pullbacks are part of the frozen factor doctrines, GP2 exact/host faithfulness retains their
\((\frac12,\frac12)\) coordinate type on both factors; deleting it on either
side would destroy a generated factor test rather than remove residual realization redundancy.
\end{proof}

At an admitted enriched finite comparison in the auxiliary all-level exchange family
\(\mathcal H_{m_x,x}\otimes\mathcal H_{1,p}\), Born exchange changes the
retained level labels whenever \(m_x\neq1\).  This family is comparison
infrastructure from Paper~I-C, not the physical reciprocal-position carrier:
the latter finite datum is the Born-derived spectral window
\(\mathcal P_x(m_x)\).  The exchange is therefore not an
unmarked automorphism of that selected auxiliary presentation.  Paper~II-A,
Proposition~\ref{p2a:prop:typed-born-branch-transport}, instead transports
the complete licensed comparison to
\(\mathcal H_{1,x}\otimes\mathcal H_{m_x,p}\); Paper~I-C,
Theorem~\ref{p1b:thm:unequal-exchange}, supplies the concrete finite-level
projection/operator intertwiner.  The two presentations and their invariant
comparison retain distinct labels.  Exchange therefore does not require
\(m_x=m_p\), and the comparison completion does not constitute a second
momentum occurrence.





\subparagraph{Operator conventions (as in Paper~II).}
Throughout we adopt the \(\Spin(4)\cong SU(2)_L\times SU(2)_R\) normalization of Paper~II.
Let \(L_i^\pm\) denote the dimensionless left/right \(\mathfrak{su}(2)\) generators on the truncated Hilbert space, with
\[
[L_i^\pm,L_j^\pm]=i\,\epsilon_{ijk}L_k^\pm,\qquad
[L_i^+,L_j^-]=0.
\]
We define the physical rotation and momentum operators by
\begin{equation}
\hat J_i:=\hbar(L_i^+ + L_i^-),
\qquad
\hat P_i:=\frac{\hbar}{\Rmom}\,(L_i^+ - L_i^-).
\end{equation}
With these conventions the compact momentum algebra and all later operator identities agree dimensionally with Paper~II.

\subsubsection{Born reciprocity and the doubled vector sector}

\begin{lemma}[GP2-derived Born-reciprocal vector-sector closure]
\label{p3:lem:born}
The reciprocal vector role derived from the Paper-II-C intrinsic evaluation
spectrum is an autonomously constituted round reciprocal role
\[
   S^3_B\times S^3_p(\Rmom),
\]
where the B-role carries the transported profile metric and is isometric to the
A-role before physical unit conversion.  If \((\mathrm{CAL}_{xp})\) is also
declared, this may be written
\[
   S^3_x(\Rcs)\times S^3_p(\Rmom),\qquad
   \Rcs=\alpha_{xp}\Rmom.
\]
No independent physical position radius follows from Born reciprocity alone.
Under an occurring position frontier the generic finite reciprocal family is
$\mathsf K_{\rm res}(m_x,1)$, with the momentum entry fixed upstream.
\end{lemma}

\begin{proof}
The master closure theorem gives the B-side role as the Born/GP2-generated
variance-opposite role from the A-side round parent.  The momentum
finite carrier is already fixed by Theorem~\ref{p3:thm:momentum}.  The
position-side finite label therefore requires its own role interface, exactly
as in Paper~II-C; no common microscopic truncation label follows from Born
reciprocity alone.
\end{proof}


\begin{theorem}[Born-reciprocal position vector-sector realization]
\label{p3:thm:position}
Over the imported Paper-II momentum parent, assume the premise-minimal
Paper-II-C/GP2 reciprocal closure.  Then the position parent-role space is the Born/GP2-generated reciprocal round role
\[
        \mathcal X_{\mathrm{vec}}^{\rm parent}\cong S^3_B
\]
with transported profile metric.  Under \((\mathrm{CAL}_{xp})\) this is written
\[
        \mathcal X_{\mathrm{vec}}^{\rm parent}\cong S^3_x(\Rcs),
        \qquad \Rcs=\alpha_{xp}\Rmom.
\]
It is abstractly isomorphic to the momentum parent; its physical metric unit is
a separate calibration.  Paper~II-C already derives the auxiliary family
$\mathsf K_{\rm res}(m,1)$.  If an actual
$(\mathrm{HIST}_x^{\rm fin})$ frontier history occurs, its least-resolution
statistic $m_x(\widetilde t)=m_x^{\rm res}(F_x(\widetilde t))$ is derived
and therefore determines the stagewise member.  The primitive frontier
gives $\mathsf K_{\rm res}(1,1)$ without a separate initial-level premise.
These are position-side specializations and do not re-prove the momentum seed.
\end{theorem}

\begin{proof}
Theorem~\ref{p3:thm:momentum} supplies $S^3_p(\Rmom)$ and the already fixed
momentum entry $m_p=1$.  Lemma~\ref{p3:lem:born} supplies the Born/GP2-generated
reciprocal round role.  The finite statements are exactly the
position-side interfaces imported from the current Paper-II reciprocal
kinematics.
\end{proof}

\begin{remark}[Standing convention for the reciprocal physical scale]
\label{p3:rem:reciprocal-scale-calibration}
Whenever a later formula in this paper writes $S^3_x(\Rcs)$, uses the
configuration metric $h$, or compares dimensionful configuration distances, it
is understood to be in the explicitly enlarged $(\mathrm{CAL}_{xp})$ domain.
Without that calibration the corresponding statement is read on the abstract
B-role $S^3_B$ with transported profile metric.  Topological results are
independent of $\alpha_{xp}$.
\end{remark}

\begin{remark}[Reciprocal position geometry versus projected coordinates]
\label{p3:rem:p2-to-p3-position}
Paper~II-C already classifies finite coordinate operators on the derived
momentum seed and, through intrinsic \(E_1\)-evaluation together with the GP2
closure theorem, instantiates the Born/GP2-generated reciprocal round parent-role
space and assigns it the reciprocal position role.  The optional
\(\mathsf P_{\rm mc}\) certificate strengthens that result only by realizing
the entire transported distance-profile family as physical G0 tests.  Any
finite physical position-side presentation is the spectral window
\(\mathcal P_x(m_x)\) selected by the separately displayed position-side
frontier interfaces.  Auxiliary \(\mathcal H_{x,m}\) modules used for
all-level exchange comparison remain mathematical infrastructure; neither is
obtained by reinterpreting a projected momentum coordinate operator as an
independent microscopic position carrier.
\end{remark}

\begin{corollary}[Imported constituted Hopf/Reeb abelian line]
\label{p3:cor:hopf-abelian-line}
On the constituted simply connected Hopf/round parent-role space, use the
upstream provenance-fixed oriented principal-circle representative
\((\eta,\xi)\) of the forced \(\lvert c_1\rvert=1\) Hopf class.  At the
principal-connection level only its homogeneous Reeb-normalized gauge class is
constitutive; the displayed \(\eta\) is a convenient representative.  The
resulting geometry determines a distinguished one-dimensional abelian line
generated by the dimensionless
Hermitian Hopf/Reeb operator
\[
\hat Q_{\rm Hopf}:=-\,i\,\Rmom\,\mathcal L_\xi,
\]
where \(\xi\) is the Reeb/Hopf vector field satisfying \(\eta(\xi)=1\) and
\(\iota_\xi d\eta=0\).  On the unit-radius contact sphere the dimensionless
Reeb generator has \(-i\mathcal L_{\xi_0}\)-eigenvalues \(q\in\mathbb Z\).
On the physical sphere \(S^3(\Rmom)\), the unit-speed vector is
\(\xi=\Rmom^{-1}\xi_0\), so \(-i\mathcal L_\xi\) has eigenvalues
\(q/\Rmom\), and \(\hat Q_{\rm Hopf}\) has integer eigenvalues \(q\).
Reversing this formula representative reverses both \(\xi\) and the displayed
generator.  Thus \(\lvert c_1\rvert=1\) and the abelian line are generated, whereas
the sign is frozen by the upstream orientation provenance.  The spectral
labels \(q\in\mathbb Z\) form the complete Hopf/Reeb character atlas; they
do not by themselves select a physically occurring matter charge.
\end{corollary}

\begin{proof}
Paper~I-A establishes the geometric Hopf/Reeb abelian generator and records the
representative Hopf normalization package.  On the relation-exact source form
\(SU(2)\simeq S^3\), Paper~II,
Theorem~\ref{p2:thm:forced-hopf-chern-unit}, proves that a principal-circle
representative has \(|c_1|=1\); the Paper-I-A orientation convention writes
\(c_1=+1\).  The simply connected global form is already derived in
Paper~II from the full GP1--GP3 chain.

No finite-level identification with the protected commutant of Paper~II-E is
used here.  The Hopf/Reeb operator generates a
\(U(1)\subset SU(2)\) action and is generally non-scalar within a nontrivial
doublet, whereas the Paper~II-E line
\(\mathbb R\Gamma_{\rm sp}\) is the traceless part of the full
\(\Spin(4)\)-commutant on \(\mathcal S_{\rm prot}\) and requires no reality hypothesis.
These are distinct typed abelian lines.

Paper~III therefore imports the abelian line and its mathematical normalization rather
than deriving them anew.
\end{proof}

\begin{remark}[Status of the abelian slot]
The existence and normalization of the abelian Hopf/Reeb direction are fixed
by the imported geometric and forced-Chern-unit results.  It is not the
protected commutant line of Paper~II-E.  Paper~III also does not identify the Hopf/Reeb direction with physical
hypercharge or electromagnetism.  In particular,
\[
 \operatorname{Spec}(\hat Q_{\rm Hopf})=\mathbb Z
 \;\not\Rightarrow\;
 \text{a physically selected charge }q,
\]
and the principal Chern unit \(c_1(P)=+1\) does not by itself select a matter
charge \(q=+1\).  Any such internal Standard-Model interpretation or charge
normalization belongs to a later identification layer, not to Paper~III-A.
\end{remark}

\begin{definition}[Dimensionless unit contact-rate convention]
\label{p3:lem:primitive-contact-speed-normalization}
Let \(c_{\rm nat}\) denote the coefficient of the oriented Reeb parameter after
the Reeb generator itself has been normalized to unit length.  We choose the
positive unit convention
\[
c_{\rm nat}=+1.
\]
This is a dimensionless parameter convention.  It fixes the positive
primitive angular parametrization but not a duration unit.  After one positive
turn is assigned duration \(T_t\),
Theorem~\ref{p3:thm:gp1-speed-calibration} derives the clocked Hopf rate and
the single generated operational ceiling
\(c_{\rm op}=2\pi R_x/T_t\).  That kinematic theorem does not by itself prove
principal ancestry for a light or field channel, choose a Hamiltonian, or
assert that every physical system follows a boundary history.  A sector law
may realize a smaller effective speed but may not insert a second independent
unlabelled ceiling.  Once undeformed same-owner ancestry is proved, the
Paper-II-A inheritance theorem identifies its front speed with
\(c_{\rm op}\).
\end{definition}


\begin{corollary}[Doubled vector parent and finite reciprocal family]
\label{p3:cor:doubled-vector}
Under Theorem~\ref{p3:thm:born-reciprocity-exhaustion}, the continuous vector
role is the abstract reciprocal product
\[
   S^3_B\times S^3_p(\Rmom).
\]
Under \((\mathrm{CAL}_{xp})\) it is written
\[
   S^3_x(\Rcs)\times S^3_p(\Rmom),\qquad
   \Rcs=\alpha_{xp}\Rmom.
\]
Paper~II-C derives the corresponding generic finite comparison family
$\mathsf K_{\rm res}(m,1)$.  On an occurring frontier history,
$m=m_x(\widetilde t)$ is derived from the actual frontier.  The primitive
frontier gives $\mathsf K_{\rm res}(1,1)$, while later higher-harmonic
frontiers may increase the position entry and leave the momentum entry at its
already derived value $1$.  On any dynamical branch retaining the frozen
momentum carrier, Paper~II-A's chronological-persistence theorem supplies
the corresponding $(\mathrm{FRZ}_p)$ conclusion; it is not an additional
premise.
\end{corollary}

\begin{proof}
Immediate from Theorems~\ref{p3:thm:momentum} and~\ref{p3:thm:position} and
the current Paper-II reciprocal finite-stage interface.
\end{proof}

\subsubsection{Why the doubled vector sector is not yet a complete phase space}
\label{p3:subsec:scalar-closure}

\begin{theorem}[Smooth doubled vector parent has a scalar symplectic defect]
\label{p3:thm:vector-sector-not-symplectic}
Under the GP2-derived reciprocal closure and the displayed reciprocal metric calibration, let
\[
   M_{\rm vec}=S^3_x(\Rcs)\times S^3_p(\Rmom).
\]
Then
\[
   H^2(M_{\rm vec};\mathbb R)=0,
\]
so the smooth compact vector parent admits no global symplectic form.  This is a
statement about the smooth parent only; it does not, by de Rham cohomology
alone, exclude every possible finite or noncommutative Hamiltonian structure.
The exact finite obstruction relevant to this paper is established separately
by Lemma~\ref{p3:lem:no-intrinsic-vector-scalar}.
\end{theorem}

\begin{proof}
Each $S^3$ has $H^1(S^3;\mathbb R)=H^2(S^3;\mathbb R)=0$.  The K\"unneth
formula therefore gives $H^2(S^3\times S^3;\mathbb R)=0$.  On a compact
six-manifold a symplectic form $\omega$ would satisfy
$\int \omega^3>0$, while exactness of every closed two-form would force
$\int \omega^3=0$ by Stokes' theorem.  Hence no global symplectic form
exists.
\end{proof}

\begin{lemma}[No full-product-intrinsic vector-scalar canonical variables: smooth and finite forms]
\label{p3:lem:no-intrinsic-vector-scalar}
Let
\[
M_{\rm vec}=S^3_x(\Rcs)\times S^3_p(\Rmom)
\]
be the smooth doubled reciprocal vector parent-role space before dynamical asymmetry.  Its pre-asymmetry vector
symmetry is the full product symmetry
\[
        G_{\rm vec}:=\Spin(4)_x\times\Spin(4)_p
\]
together with Born exchange of the two reciprocal factors.  By an intrinsic vector-scalar observable
we mean a smooth observable on \(M_{\rm vec}\) that is defined without adding an origin, preferred
frame, sector intertwiner, or relative-coordinate marker, and is invariant under this full
pre-asymmetry vector symmetry.

Every such intrinsic vector-scalar observable is constant.  Hence two such intrinsic vector-scalars
\[
\tau,\epsilon:M_{\rm vec}\to\mathbb R
\]
cannot form a local canonical pair:
\[
\{\tau,\epsilon\}=0
\]
for every local Poisson bracket on the vector parent-role space.

To test whether an all-level compact exchange comparison could manufacture
such a scalar canonical pair, consider the auxiliary finite representation family at label
$m_x$ (with the physical momentum entry already fixed to $1$), and define its
full product-invariant endomorphism algebra
\[
 \mathcal A^{\rm inv}_{m_x,1}
 :=\operatorname{End}_{G_{\rm vec}}
 \bigl(\mathcal H_{x,m_x}\otimes\mathcal H_{p,1}\bigr).
\]
Here \(\mathcal H_{x,m_x}\) is an auxiliary exchange-comparison module, not
the physical reciprocal-position stage \(\mathcal P_x(m_x)\).  This
comparison algebra can contain non-scalar sector projectors, but it is
abelian.
Born exchange at a generic asymmetric stage is a typed algebra isomorphism
\[
 \sigma_{m_x,1}:\mathcal A^{\rm inv}_{m_x,1}
 \stackrel{\sim}{\longrightarrow}\mathcal A^{\rm inv}_{1,m_x},
\]
not an endomorphism of \(\mathcal A^{\rm inv}_{m_x,1}\).  Thus a
Born-compatible observable is a pair
\((O,\sigma_{m_x,1}O)\) across the exchanged stages.  Only at the symmetric
seed \(m_x=1\) is an exchange-fixed subalgebra of the same typed algebra
defined; that subalgebra is still abelian.
Consequently no two invariant observables \(T,E\) inside this auxiliary
finite comparison algebra can obey an exact Heisenberg relation \([T,E]=i\hbar\mathbf1\).  At the symmetric seed $m_x=1$, the full product-group commutant has complex dimension
\(36\).
\end{lemma}

\begin{proof}
Full product transitivity makes every smooth invariant function constant, so
its differential vanishes and every local Poisson bracket of two such
functions is zero.  On the auxiliary finite comparison family the product representation is
multiplicity free; Schur's lemma therefore identifies its invariant
endomorphism algebra with a finite direct sum of copies of \(\mathbb C\), so
it is abelian.  Factor exchange maps the \((m_x,1)\) algebra to the
\((1,m_x)\) algebra and becomes an internal involution only for \(m_x=1\).
The explicit decomposition, the dimension count \(6^2=36\), and the
diagonal-symmetry boundary are proved in
Paper~III-G, Appendix~\ref{p3a:app:vector-scalar-obstruction-details}.
\end{proof}

\begin{corollary}[No full-product-intrinsic scalar evolution variable inside the compact vector sector]
\label{p3:cor:no-intrinsic-scalar-pair}
The compact doubled vector sector, considered with its full pre-asymmetry
product vector symmetry, contains no globally defined intrinsic scalar evolution variable together
with a conjugate scalar generator.  At finite level, invariant endomorphisms
may include non-scalar block projectors, but their algebra is abelian and
therefore contains no exact canonical pair.  The required scalar phase data cannot be hidden inside the vector product.  The separately prior-generated scalar role is therefore the relevant available scalar input to the Born completion.
\end{corollary}

\begin{proof}
Theorem~\ref{p3:thm:vector-sector-not-symplectic} excludes a global compact Hamiltonian structure on
the doubled vector parent-role space.  Lemma~\ref{p3:lem:no-intrinsic-vector-scalar} gives the sharper local
obstruction relevant here: even before asking for a global symplectic form, the homogeneous vector parent-role space has no nonconstant scalar variables invariant under the full pre-asymmetry product vector
symmetry that could satisfy a local canonical relation.  Hence the canonical
scalar pair cannot be internal to \(M_{\rm vec}\); it belongs to the separately prior-generated and then Born-completed scalar torus.  The finite statement follows from
the explicit commutant calculation in the same lemma.
\end{proof}


\begin{remark}[Transported scalar labels and remaining scope]
\label{p3:rem:provisional-scalar-labels}
The scalar labels \(t\) and \(E\) are no longer assigned by a free downstream
sign convention.  Paper~I-A fixes the positive primitive Reeb lattice and
Paper~II-A transports it to the actual scalar-succession role through
\(\mathsf{HS}_t\).  The theorem below fixes the unique reciprocal energy
primitive.  This oriented pair supplies neither a thermodynamic arrow nor a
choice of Hamiltonian, actual trajectory, duration unit, or energy unit.
\end{remark}

\begin{theorem}[Transported scalar succession and dual energy orientation]
\label{p3:thm:transported-scalar-orientation}
Import the Paper-I-A positive primitive Reeb lattice and the Paper-II-A
Hopf-succession provenance theorem.  In Paper~III-A's scalar presentation,
identify the imported \(H_{\rm succ}\) with \(A_{\rm s}=V/\Lambda\).  Then
the scalar A-lattice has the fixed positive primitive
\[
 e_t=(\kappa_t)_*\tau_H^+.
\]
Its full dual lattice has a unique primitive \(e_E\) satisfying
\[
 \boxed{e_E(e_t)=1.}
\]
Consequently the Poincar\'e evaluation form has
\[
 \Omega(e_E,e_t)=1,
 \qquad
 \Omega(e_t,e_E)=-1,
\]
and on the universal cover or a local torus chart
\[
 \boxed{\omega_{\rm sc}=dE\wedge dt,\qquad\{t,E\}=1.}
\]
The positive \(t\)-direction is the transported Paper-I-A Reeb direction.  The
sign is derived; the duration unit, energy unit, phase-cell calibration,
choice of Hamiltonian, and assertion that a system follows a particular
trajectory remain separate.
\end{theorem}

\begin{proof}
Corollary~\ref{p1:cor:positive-reeb-lattice} supplies the positive primitive
\(\tau_H^+\), and
Theorem~\ref{p2:thm:transported-succession-orientation} supplies the retained
isomorphism of period lattices and its primitive image \(e_t\).  Because
\(\Lambda\) has rank one, its full dual \(\Lambda^*\) contains exactly one
primitive \(e_E\) with \(e_E(e_t)=1\).  Substitution into the alternating
evaluation form
\[
 \Omega\bigl((x,\xi),(y,\eta)\bigr)=\xi(y)-\eta(x)
\]
gives the two displayed signs.  The associated translation-invariant local
form is \(dE\wedge dt\), with inverse Poisson relation
\(\{t,E\}=1\).  Positive rescaling changes calibration but cannot reverse the
transported primitive.
\end{proof}

\begin{theorem}[Canonical Poincar\'e phase class of the scalar A/B pair]
\label{p3:thm:scalar-born-residue}
Let
\[
        A_{\rm s}=V/\Lambda,
        \qquad
        B_{\rm s}=V^*/\Lambda^*
\]
be the scalar transformation torus and its full reciprocal test-dual torus from
Theorem~\ref{p3:thm:scalar_classification_import}.  The evaluation pairing
\[
        \langle\cdot,\cdot\rangle:
        \Lambda^*\times\Lambda\longrightarrow\ZZ
\]
is perfect and determines the canonical Poincar\'e line bundle
\(\mathcal P\to A_{\rm s}\times B_{\rm s}\).  We use the
evaluation-character convention whose alternating lattice representative is
\((x,\xi),(y,\eta)\mapsto\xi(y)-\eta(x)\).  Its first Chern class
\[
        c_{\rm sc}:=c_1(\mathcal P)
        \in H^2(A_{\rm s}\times B_{\rm s};\ZZ)
\]
is primitive.  Theorem~\ref{p3:thm:transported-scalar-orientation} supplies
the positive scalar primitive and its unique full-dual primitive, so the
ordered basis \((e_E,e_t)\) is positive with \(e_E(e_t)=1\).
Therefore the scalar Born product carries a canonical nontrivial primitive
integral phase class relative to the evaluation pairing.
\end{theorem}

\begin{proof}
The scalar classification and transported orientation above supply
the full dual lattice and $e_E(e_t)=1$. In ordered $(t,E)$ coordinates,
Paper~I-C uses $\eta=[d\theta_t\wedge d\theta_E/(2\pi)^2]$ and the
dual evaluation convention. Thus its primitive phase result,
Proposition~\ref{p1b:prop:scalar-moments}~\cite{Kugelmann2026Paper1B},
applies here with
\[
 \mathcal P_{\mathrm{III-A}}\cong\mathcal P_{\mathrm{1B}}^*,
 \qquad c_{\rm sc}=-\eta.
\]
It evaluates to $+1$ on $[e_E\times e_t]$ and to $-1$ on
$[e_t\times e_E]$. The physical primitive is the transported
$\mathsf{HS}_t$ primitive already fixed here: its full-dual evaluation
generates the primitive class, rather than an arbitrary integral multiple.
GP2 exact/faithful realization retains that generated class.
Paper~III-G, Appendix~\ref{p3a:app:v166-phase-gluing} constructs the bundle and its
loop holonomy directly, with \(c_{\rm sc}=[dv\wedge du]\) in normalized
cover coordinates. The phase is encoded by bundle gluing, not a global
torus bicharacter.
\end{proof}

\begin{theorem}[Canonical scalar lattice pairing from reciprocal closure]
\label{p3:thm:minimal-symplectic-completion}
Let \(e_t\) and \(e_E\) be the dual primitive lattice generators above.  The
alternating form on
\[
        \Lambda_{\rm sc}:=\Lambda\oplus\Lambda^*
\]
defined by
\[
        \Omega\bigl((x,\xi),(y,\eta)\bigr)
        :=\xi(y)-\eta(x)
\]
has matrix
\[
        J=
        \begin{pmatrix}0&-1\\[2pt]1&0\end{pmatrix}
\]
in the basis \((e_t,e_E)\).  It is integral and unimodular, so it descends to
the scalar product torus as the primitive nondegenerate phase class.  In
particular,
\[
 \Omega(e_t,e_E)=-1,
 \qquad
 \Omega(e_E,e_t)=+1,
\]
which is the reciprocal orientation fixed above.  On its
universal cover or in a local torus chart it has Darboux form
\[
        \omega_{\rm sc}=dE\wedge dt,
        \qquad
        \{t,E\}=1.
\]
Because the GP2-derived scalar Born product has exactly the two rank-one A/B factors, this form has rank two and hence gives one scalar canonical plane.
\end{theorem}

\begin{proof}
Perfectness of the lattice evaluation pairing gives
\(\det J=1\), hence the induced map
\(\Lambda_{\rm sc}\to\Lambda_{\rm sc}^*\) is an integral isomorphism.  The
corresponding translation-invariant two-form is therefore closed,
nondegenerate, and primitive on the product torus.  Any multiplication by
\(|N|>1\) would make the lattice map have determinant \(N^2\) and nontrivial
cokernel, so it is not the Poincar\'e evaluation form.  The zero multiple is
degenerate.  Darboux's theorem gives the displayed local relation, while the
rank count gives exactly one canonical plane.
Theorem~\ref{p3:thm:transported-scalar-orientation} identifies its oriented
coordinates as scalar succession and reciprocal energy.  Selecting a
Hamiltonian and an actual evolution remains a dynamical step.
\end{proof}

\begin{theorem}[Canonical reciprocal generator-germ localization]
\label{p3a:thm:canonical-generator-germ-localization}
Let \(b\in S_B^3\), let \(a=\iota_1^{-1}(b)\), and let \(t\in S_t^1\).
Paper~II-C, Theorem~\ref{p2c:thm:first-symbol-pairing}, gives the canonical
nondegenerate pairing
\[
 \mathcal P_a:T_b^*S_B^3\times\mathfrak p_a\longrightarrow\mathbb R,
\]
where \(\mathfrak p_a\) is the generated momentum-transvection role at \(a\).
Hence evaluation induces a canonical isomorphism
\[
 \Phi_{v,b}:\mathfrak p_a^*\xrightarrow{\ \sim\ }T_b^*S_B^3.
\]
The transported scalar primitive and its full dual give, independently,
\[
 \Phi_{s,t}:\Lambda^*\otimes_\mathbb Z\mathbb R
 \xrightarrow{\ \sim\ }T_t^*S_t^1,
 \qquad e_E\longmapsto\theta_t,
\]
where \(\theta_t(e_t)=1\).  Therefore on every product chart
\(U\subset S_B^3\times S_t^1\), the direct sum
\begin{equation}
\label{p3a:eq:generator-germ-cotangent}
 \boxed{
 \Phi_U:
 (\iota_1^{-1})^*(\mathfrak p^*)\oplus
 (\Lambda^*\otimes_\mathbb Z\mathbb R)
 \xrightarrow{\ \sim\ }T^*U }
\end{equation}
is canonical, full-rank, evaluation-preserving, restriction-compatible, and
preserves the vector/scalar provenance.  Under coordinate change it transforms
by the ordinary cotangent representation; no metric curry, frame, or
result-dependent intertwiner is used.

Consequently the local cotangent generator role is a finitely presented
generated-role descendant of the already occurring reciprocal generator role.
By Paper~II-A, Theorem~\ref{p2:thm:generated-role-corealization}, it is
co-realized without a second G0 localization occurrence.  This theorem
identifies only the local generator/readout germs.  Together with Paper~II-C,
Corollary~\ref{p2c:cor:compact-transvection-cotangent-germ}, it also fixes the
exact first-order compact-transvection/cotangent-germ correspondence.  It does
not identify the global compact carrier \(S_p^3\times S_E^1\) with the
noncompact cotangent bundle and does not discharge any stronger global or
all-resolution finite-band operator intertwiner, nor any dimensionful
reciprocal calibration.
\end{theorem}

\begin{proof}
Nondegeneracy of \(\mathcal P_a\) identifies \(T_b^*S_B^3\) with the linear
dual of \(\mathfrak p_a\) uniquely: a covector is sent to its evaluation
functional on the transvection module.  The rank-one lattice pairing is
perfect and the transported positive primitive fixes the unique dual energy
primitive, so the scalar generator line is the cotangent line to the
succession circle.  Cotangent spaces of a product split canonically, giving
\eqref{p3a:eq:generator-germ-cotangent}.  Naturality of the II-C pairing and
of lattice evaluation gives restriction and coordinate compatibility; the
unique evaluation-preserving change of generator coordinates is the dual
pullback.  The construction is finite-rank, canonical and nonredundant, so it
satisfies the generated-role-descendant criteria imported from Paper~II-A and
inherits co-realization from the same reciprocal occurrence.
\end{proof}

\begin{corollary}[Calibrated reciprocal-position momentum unit]
\label{p3a:cor:calibrated-reciprocal-position-momentum-unit}
Assume $(\mathrm{CAL}_{xp})$, so that the transported reciprocal-position
radius is
\[
 R_x=\alpha_{xp}\Rmom,
 \qquad \alpha_{xp}>0.
\]
The natural dimensionful momentum unit conjugate to an orthonormal
reciprocal-position frame is then
\begin{equation}
 \label{p3a:eq:calibrated-reciprocal-position-momentum-unit}
 \boxed{
 p_{x,\rm unit}:=\frac{\hbar}{R_x}
 =\frac1{\alpha_{xp}}\frac{\hbar}{\Rmom}
 =\frac{\Punit}{\alpha_{xp}}.}
\end{equation}
Thus the microscopic compact normalization $\Punit=\hbar/\Rmom$ and the
reciprocal-position cotangent normalization are related by the inverse
covector scaling forced by $(\mathrm{CAL}_{xp})$.  No equality
$\alpha_{xp}=1$, state, momentum value, or readout token follows.
\end{corollary}

\begin{proof}
The reciprocal-position calibration rescales spatial lengths by
$\alpha_{xp}$.  Cotangent components therefore rescale inversely, and the
physical action unit $\hbar$ gives $\hbar/R_x$.  Substitution of
$R_x=\alpha_{xp}\Rmom$ and $\Punit=\hbar/\Rmom$ proves the identity.
\end{proof}

\begin{corollary}[Local cotangent compatibility with the reciprocal tangent limit]
\label{p3a:cor:local-cotangent-born-compatibility}
On the retained common-model identification
\(S_B^3\times S_t^1\cong K=SU(2)\times U(1)\), every generated chart
\(U\subset S_B^3\times S_t^1\) identifies with an open
\(U_K\subset K\), and its cotangent lift gives a canonical symplectomorphism
\[
 T^*U\cong T^*K|_{U_K}.
\]
The mathematical Born involution induced by the Paper-I-C tangent deformation
acts on this zero-fibre chart as
\[
 (q,\xi)\longmapsto(q,-\xi),
 \qquad
 R^*\omega_{\rm can}=-\omega_{\rm can}.
\]
This is a compatibility statement: the local cotangent role was already
generated by Theorem~\ref{p3a:thm:canonical-generator-germ-localization}.
Neither this corollary nor the tangent deformation identifies the compact
reciprocal host \(M_8\) with \(T^*K\).
\end{corollary}
\begin{proof}
A diffeomorphism of base charts has a canonical cotangent lift preserving the
Liouville form and therefore \(\omega_{\rm can}\).  Paper~I-C's pair-groupoid
swap descends at the tangent-groupoid zero fibre to vector reversal and, after
fibre Fourier transform, to covector reversal.  The latter is
anti-symplectic.  Restriction to \(U_K\) gives the displayed local statement.
\end{proof}


\begin{corollary}[Scalar dynamical sector: exactly one conjugate pair]
\label{p3:cor:missing-sector-one-pair}
The Paper-II scalar classification, the premise-minimal GP2-derived reciprocal
closure, and the Poincar\'e
evaluation pairing give precisely
\[
        S^1_t\times S^1_E
\]
with one local/universal-cover canonical pair inside that constitution.
\end{corollary}

\begin{proof}
Theorem~\ref{p3:thm:scalar_classification_import} gives exactly two scalar
circle factors, and Theorem~\ref{p3:thm:minimal-symplectic-completion} gives
their unique primitive nondegenerate integral pairing.  Therefore the scalar
sector is exactly one conjugate pair.
\end{proof}

\begin{lemma}[Calibrated scalar phase-cell normalization]
\label{p3:lem:minimal-phase-torus-quantization}
Assume the physical phase-unit calibration \((\mathrm{CAL}_{\hbar})\), namely
\[
   \left[\frac{\omega_{\rm sc}}{2\pi\hbar}\right]=c_{\rm sc}.
\]
Choose angular coordinates
\(\theta_t,\theta_E\in\RR/2\pi\ZZ\), calibrated periods
\(\tper,\Eper>0\), and global closed one-forms
\[
        \Theta_t=\frac{\tper}{2\pi}d\theta_t,
        \qquad
        \Theta_E=\frac{\Eper}{2\pi}d\theta_E.
\]
Orient the scalar torus by
\([S^1_E\times S^1_t]\), equivalently by the positive ordered coframe
\(d\theta_E\wedge d\theta_t\).
For the primitive scalar pairing of
Theorem~\ref{p3:thm:minimal-symplectic-completion},
\[
        \omega_{\rm sc}=\Theta_E\wedge\Theta_t,
\]
and the calibration together with
\(\langle c_{\rm sc},[S^1_E\times S^1_t]\rangle=+1\) gives
\[
 \left\langle
 \left[\frac{\omega_{\rm sc}}{2\pi\hbar}\right],
 [S^1_E\times S^1_t]
 \right\rangle=+1.
\]
Therefore
\[
        \boxed{\tper\Eper=2\pi\hbar}.
\]
Here \(\hbar\) is the declared physical phase unit of the already derived primitive
phase class; it does not select the topology or the integer class.  Equivalently,
\[
 \boxed{
 c_{\rm sc}\text{ primitive is generated},
 \qquad
 [\omega_{\rm sc}]=2\pi\hbar\,c_{\rm sc}
 \text{ uses }(\mathrm{CAL}_{\hbar}).}
\]
\end{lemma}

\begin{proof}
The integral of the displayed form is
\[
        \int_{[S^1_E\times S^1_t]}\omega_{\rm sc}=\tper\Eper.
\]
Theorem~\ref{p3:thm:minimal-symplectic-completion} fixes the positive primitive
Chern number to one.  The separate calibration \(\mathrm{CAL}_{\hbar}\) identifies its physical representative with \(2\pi\hbar\) times that class, giving
\[
        \frac{\tper\Eper}{2\pi\hbar}=1,
\]
which is the claimed relation.
\end{proof}




\begin{lemma}[The scalar torus is the GP2-derived Born product]
\label{p3:def:separable-scalar-compactification}
\label{p3:lem:separable-scalar-compactification}
The compact scalar role space is
\[
        T^2=S^1_t\times S^1_E
\]
because the Paper-II scalar A-role and its generated contravariant B-role
satisfy the intrinsic lattice/evaluation version of the Paper-II reciprocal
antecedents; Paper~II-A's GP2 theorem objectifies the B-role and supplies their universal relation-free
commuting product.  It is not
obtained by choosing a lattice inside an
otherwise unmarked two-plane.  A nontrivial diagonal quotient has only the
proper invariant subalgebra of the full product tests and is excluded by
GP2 relation-free exactness.  Independently, if such a quotient is an oriented
finite covering and the primitive Poincar\'e class is required to descend
unchanged, the degree argument forces absolute degree one.  A twisted or
higher-dimensional realization would require additional constitutive data.
\end{lemma}

\begin{proof}
This is the scalar specialization of
Paper~II-A's GP2 reciprocal closure, Lemma~\ref{p3:lem:rank-one-lattice-dual-spectrum},
and Theorem~\ref{p3:thm:scalar_classification_import}.  The product lattice is the
direct sum of the two already constituted circle lattices.  No independent
choice of a two-dimensional compactification lattice occurs.
\end{proof}

\begin{theorem}[Scalar phase torus and universal-cover pair]
\label{p3:thm:scalars}
Under the Paper-II scalar classification and the GP2-derived master
reciprocal closure, including the G0 phase-unit calibration
\((\mathrm{CAL}_{\hbar})\),
the scalar closure is
\[
        (t,E)\in S^1_t\times S^1_E,
        \qquad
        \tper\Eper=2\pi\hbar,
\]
with one primitive integral phase class and, under $(\mathrm{CAL}_{\hbar})$, the displayed physical cell normalization.  On local charts and on the universal
covers
\[
        \widetilde t,\widetilde E\in\RR,
\]
the phase form is
\[
        d\widetilde E\wedge d\widetilde t,
        \qquad
        \{\widetilde t,\widetilde E\}=1.
\]
The compact circles form the microscopic scalar role space; causal evolution and
unbounded generator values use their universal covers.  The scalar succession
orientation is inherited from Paper~I-A through \(\mathsf{HS}_t\); it is not a
thermodynamic arrow and does not select an actual evolution.  The unit-dependent
numerical presentation remains a calibration.  Once one primitive turn is
assigned duration \(T_t\), Theorem~\ref{p3:thm:gp1-speed-calibration} derives
the nonstationary clocked Hopf history and its finite operational ceiling; no
independent propagation-family occurrence is required.
\end{theorem}

\begin{proof}
Theorem~\ref{p3:thm:scalar_classification_import} gives the product torus from
the GP2-objectified dual role in the master interface domain,
Theorem~\ref{p3:thm:minimal-symplectic-completion} gives its positive
primitive integral pairing, and
Lemma~\ref{p3:lem:minimal-phase-torus-quantization} restores the physical phase
unit.  Passing to the universal cover replaces the angular phases by real
coordinates without changing the compact microscopic role space.
\end{proof}


\begin{theorem}[Integral Born-eigenphase cover of the scalar parent]
\label{p3a:thm:v190-scalar-born-eigenphase-cover}
Let $z_t=e^{i\theta_t}$ and $z_E=e^{i\theta_E}$ be the two primitive
scalar factors of $T^2_{tE}=S_t^1\times S_E^1$.  Define
\[
 q_+:=z_tz_E,\qquad q_-:=z_tz_E^{-1},
\]
and
\[
 A=\begin{pmatrix}1&1\\1&-1\end{pmatrix},\qquad
 \binom{\alpha}{\beta}=A\binom{\theta_t}{\theta_E}.
\]
Then
\[
 \det A=-2,\qquad A^{-1}=\frac12A,\qquad A^2=2I,
\]
and
\[
 Q:T^2_{tE}\longrightarrow T^2_{\pm},\qquad
 Q(z_t,z_E)=(q_+,q_-)
\]
is a two-fold cover with
\[
 \boxed{\ker Q=\{(1,1),(-1,-1)\}\cong\mathbb Z_2.}
\]
Born exchange is diagonalized on the quotient variables,
\[
 q_+\longmapsto q_+,\qquad q_-\longmapsto q_-^{-1},
\]
and the normalized infinitesimal generators are
\[
 \boxed{N_+=\frac12(N_t+N_E),\qquad
 N_-=\frac12(N_t-N_E).}
\]
If $J=\begin{psmallmatrix}0&-1\\1&0\end{psmallmatrix}$ is the primitive
Poincar\'e lattice form, then
\[
 \boxed{A^TJA=-2J.}
\]
Thus the Born eigenlattice has index two and symplectic cell two, whereas the
primitive $(t,E)$ lattice remains unimodular.  The quotient variables are
generated readouts of the scalar parent, not a replacement of the primitive
scalar torus.
\end{theorem}

\begin{proof}
Apply Paper~I-C, Theorem~\ref{p1b:thm:eigenphase}, with the same
ordered angles $(\theta_t,\theta_E)$. It gives the covering kernel,
character map, exchange action, and inverse-angle generators. The
alternating matrix used here is the negative of the ordered positive
form in that convention; nevertheless $A^TJA=(\det A)J=-2J$.
The quotient readouts therefore retain the primitive parent and its
two-sheeted cover.
\end{proof}

\begin{corollary}[Primitive phase-class non-descent and scalar deck grading]
\label{p3a:cor:v190-scalar-deck-grading}
Choose the orientation of $T^2_\pm$ so that its primitive class $c_\pm$ obeys
\[
 \boxed{Q^*c_\pm=2c_{\rm sc}.}
\]
Consequently $c_{\rm sc}$ itself does not descend through $Q$, whereas
$2c_{\rm sc}$ does.  On character lattices,
\[
 Q^*(r,s)=(r+s,r-s),
\]
so there is an exact sequence
\[
 \boxed{0\longrightarrow\mathbb Z^2_{\pm}
 \xrightarrow{\ Q^*\ }\mathbb Z^2_{tE}
 \xrightarrow{\ \epsilon_{\rm sc}\ }\mathbb Z_2\longrightarrow0,}
 \qquad \epsilon_{\rm sc}(n,w)=n+w\pmod2.
\]
Equivalently the parent character algebra is strongly $\mathbb Z_2$-graded:
its deck-even part is generated by $q_\pm$, while the deck-odd part is the
other lattice coset and squares into the even algebra.  In particular every
character descending through either single quotient $q_+$ or $q_-$ is
deck-even.  This grading is distinct from the relative half-turn parity
$(-1)^n$.
\end{corollary}

\begin{proof}
Paper~I-C, Theorem~\ref{p1b:thm:eigenphase} and
Corollary~\ref{p1b:cor:relative-characters}, give the exact sequence and
strong character grading. With $c_{\rm sc}=-\eta$ from
Theorem~\ref{p3:thm:scalar-born-residue}, choose the standard ordered
target generator $c_\pm=\eta_\pm$. Its pullback is
$-2\eta=2c_{\rm sc}$, as stated. The quotient axes have even deck parity;
the separately defined relative half-turn still acts by $(-1)^n$.
\end{proof}

\begin{corollary}[Mutual scalar circle-valued moments]
\label{p3a:cor:v190-scalar-mutual-moments}
On the calibrated scalar phase form
\[
 \omega_{\rm sc}=\frac{\hbar}{2\pi}\,d\theta_E\wedge d\theta_t,
\]
let $V_+=\partial_{\theta_t}+\partial_{\theta_E}$ and
$V_-=\partial_{\theta_t}-\partial_{\theta_E}$.  Then
\[
 \boxed{\iota_{V_+}\omega_{\rm sc}
 =\frac{\hbar}{2\pi}d(\theta_t-\theta_E),\qquad
 \iota_{V_-}\omega_{\rm sc}
 =-\frac{\hbar}{2\pi}d(\theta_t+\theta_E).}
\]
Thus $q_-$ is the circle-valued moment of the diagonal action and $q_+$ is
its anti-diagonal companion.  No Hamiltonian member or actual trajectory is
selected by this structural statement.
In the notation of Paper~I-C, Proposition~\ref{p1b:prop:scalar-moments},
\[
 \omega_{\mathrm{III-A}}=-2\pi\hbar\,\omega_{\mathrm{1B}}.
\]
Multiplying its contractions by this factor gives the displayed signs:
the diagonal primitive is $q_-$ in this calibration, while the
anti-diagonal primitive has the opposite signed angular differential of
$q_+$.
\end{corollary}

\begin{remark}[Scalar-parent retention firewall]
\label{p3a:rem:v190-scalar-parent-retention}
The two Born quotient phases do not contain the primitive Poincar\'e class or
the deck-odd character coset.  They therefore cannot replace
$S_t^1\times S_E^1$ in the physical constitution.  In particular the
relative $m_{\rm cmp}=2$ comparison window samples only the deck-even quotient
sector and does not determine the full scalar parent lattice.
\end{remark}

\begin{definition}[Clocked Hopf composite and generated propagation doctrine]
\label{p3a:def:clocked-hopf-propagation-doctrine}
\label{p3a:def:full-nonstationary-propagation-occurrence}
Let \(\kappa_t:U(1)_H\to H_{\rm succ}\) be the retained
Hopf-to-succession isomorphism, let
\(\iota_H^{\rm op}:U(1)_H\hookrightarrow A_{\rm v}^{\rm op}\) be the typed
inclusion of Paper~II-A,
Lemma~\ref{p2:lem:chosen-hopf-vector-parent-inclusion}, and let
\(B_{\rm v}=B_{A_{\rm v}}^{(1)}\) be the premise-minimal reciprocal role of
Paper~II-C, Theorem~\ref{p2c:thm:coordinate-generated-reciprocal-role}.
Use that theorem's contragredient convention
\(T^\vee=\iota_1T^{-1}\iota_1^{-1}\) and put \(R_g(a)=ag\). Define
\[
 \rho_{\rm v}^{(1)}:A_{\rm v}^{\rm op}
 \longrightarrow\operatorname{Isom}(B_{\rm v},d_{B_{\rm v}}),
 \qquad \rho_{\rm v}^{(1)}(g):=(R_{g^{-1}})^\vee.
\]
Then \(\rho_{\rm v}^{(1)}(g)\iota_1(a)=\iota_1(ag)\) and
\(\rho_{\rm v}^{(1)}(g)\rho_{\rm v}^{(1)}(h)=
\rho_{\rm v}^{(1)}(hg)\), as required for \(A_{\rm v}^{\rm op}\).
Evaluation at the identity proves faithfulness. Inverse translation uses the
already constituted group inverse and preserves the intended Hopf primitive;
literal restriction \((R_g)^\vee\) would reverse that primitive.
This action requires no metric-curry/full-profile premise.  The typed
composite is
\begin{equation}
\label{p3a:eq:clocked-hopf-composite}
 \boxed{
 \gamma_{tp}^{\rm op}:=\iota_H^{\rm op}\circ\kappa_t^{-1}:
 H_{\rm succ}\longrightarrow A_{\rm v}^{\rm op},\qquad
 \beta_t:=\rho_{\rm v}^{(1)}\circ\gamma_{tp}^{\rm op}:
 H_{\rm succ}\longrightarrow
 \operatorname{Isom}(B_{\rm v},d_{B_{\rm v}}).}
\end{equation}
The intermediate \(\gamma_{tp}^{\rm op}\) is therefore an actually used
operation/provenance morphism, not a support graph: it fixes the inherited
Hopf axis, handedness, primitive, and winding of the clock-generated
vector-parent operation but does not assert \(p=p(t)\) for the local cotangent
generator role.

Assume the explicit clock calibration \((\mathrm{CAL}_{t})\), assigning the
positive duration \(T_t\) to one turn of the already derived positive
primitive \(e_t\).  On the elapsed universal cover define the seed history
\begin{equation}
\label{p3a:eq:clocked-hopf-seed}
 \Phi^H_{\widetilde t}
 :=\beta_t\!\left(
 \exp_{H_{\rm succ}}\!\left(\frac{\widetilde t}{T_t}e_t\right)
 \right).
\end{equation}
Before evaluation, declare the physical history/doctrine target sort
\(\mathsf{HistDoctrine}_{tp}\) and the typed constructor
\[
 \mathsf{CLK}_{tp}:
 D_{tp}:=\bigl(\mathsf{HS}_t,\mathrm{CAL}_t,
 A_{\rm v},B_{\rm v},\iota_H^{\rm op},\rho_{\rm v}^{(1)}\bigr)
 \longmapsto
 \bigl(\Phi^H,\mathscr O_{\rm prop}^{\rm phys}\bigr).
\]
At this stage the frozen clock-propagation grammar already licenses the
displayed exponential, the retained composites
\(\gamma_{tp}^{\rm op},\beta_t\), and the stated finite closure operations.
Let \(w_{\rm CLK}\) denote that word.  Its inputs are precisely \(D_{tp}\);
all intermediates are retained in their native sorts, and
\[
 \operatorname{Orig}_{\Sigma_{tp}^{\rm CLK}}
 \bigl(w_{\rm CLK};D_{tp},\mathsf{HistDoctrine}_{tp}\bigr)
\]
holds before the seed is evaluated.
Its physical essential image
\(\mathscr O_{\rm prop}^{\rm phys}\) is the least collection of
interval-parametrized spatial-isometry histories containing
\(\Phi^H\) and the stationary history and closed under restriction, elapsed
time translation, conjugation by already generated round spatial isometries,
and finite compatible sequential concatenation.  These typed constructors
neither add a channel label nor permit an independent cover reparametrization,
a higher winding \(z\mapsto z^n\), or a rate multiplier.

Separately, let \(\mathscr V_x^{\rm phys}\subset T_xS_x^3\) be the set of
metric-derivative velocity germs at \(x\) supplied by
\(\mathscr O_{\rm prop}^{\rm phys}\), and define the \emph{mathematical
relaxed velocity envelope}
\[
 \mathscr V_x^{\rm rel}
 :=\overline{\operatorname{co}}\bigl(\mathscr V_x^{\rm phys}\bigr)
 \subset T_xS_x^3.
\]
Write \(\mathscr O_{\rm prop}^{\rm rel}\) for the absolutely continuous
point histories \(\gamma:I\to S_x^3\) whose metric derivative belongs to
\(\mathscr V_{\gamma(t)}^{\rm rel}\) for almost every \(t\).  Thus
``time-sharing'' means convex relaxation of tangent controls (equivalently,
limits of finite chattering controls), not affine averaging of points or maps
on the sphere.  This mathematical curve family is used below only to describe
the closed convex causal cone.  Convex relaxation and local-uniform completion
are not physical constructors or additional occurring histories.

The occurrence notation
\((\mathrm{OCC}_{\rm prop}^{\rm full,\ne0})\) is not an
independent occurrence premise.  By
Lemma~\ref{p3a:lem:generated-role-occurrence-inheritance}, the complete joint
root tuple \(D_{tp}\) transports through the choice-free faithful composite to
the nonstationary seed and co-realizes its exact physical essential image.
This does not assert that every member history of
\(\mathscr O_{\rm prop}^{\rm phys}\) occurs as a separate physical token, and
it asserts no physical occurrence for members added only by
\(\mathscr O_{\rm prop}^{\rm rel}\).
\end{definition}

\begin{corollary}[Co-realization of the clocked Hopf propagation doctrine]
\label{p3a:cor:clocked-hopf-corealization}
On the complete joint domain \(D_{tp}\), the nonstationary seed history and
its least generated physical propagation doctrine inherit the same occurring
provenance/calibration root tuple.  In the constructor-relative sense of
Lemma~\ref{p3a:lem:generated-role-occurrence-inheritance},
\[
 \boxed{
\operatorname{CoReal}
 \left(
 \Phi^H,\,\mathscr O_{\rm prop}^{\rm phys}
 \mid D_{tp};\mathsf{CLK}_{tp}
 \right).}
\]
Every retained restriction, elapsed-translation, conjugation, and sequential
composition law is an incidence edge of this same rooted process component.
On each generated separating evaluation role, its pullback/test dynamics is
the unique co-realized dual family of
Theorem~\ref{p2a:thm:evaluation-dual-corealization}; it is not a second
physical propagation.
No selector among member histories is introduced by this statement: it
constitutes the generated propagation role/doctrine, while a particular
physical trajectory remains a separate token-level occurrence.  No member of
the relaxed envelope is promoted to a physical token by this corollary.
\end{corollary}

\begin{proof}
All noncalibration entries of \(D_{tp}\) are jointly generated descendants of
the frozen Paper-I-A/II transfer and reciprocal package; \(\mathrm{CAL}_t\) is
the displayed occurring comparison datum.  The composite is natural under
isomorphisms preserving the role tags and the Hopf provenance, and is faithful
because \(\kappa_t\), \(\iota_H^{\rm op}\), and
\(\rho_{\rm v}^{(1)}\) are faithful.  Its target sort and closure operations
are declared before evaluation, its essential image is exactly the least
closure displayed in the definition, and it introduces no scale beyond
\(\mathrm{CAL}_t\), no channel, relation, or member selector.  The standard
clocked action is a GP1-bounded faithful realization and the complete tuple has
zero simultaneous insertion, omission, relation, type, and closure defect.
Paper~II-A,
Theorem~\ref{p2a:thm:rooted-generated-image-corealization}, applied to
\(w_{\rm CLK}\), therefore gives the stated role-level co-realization.  The
naturality, faithfulness, and zero-defect checks discharge the obstruction;
they do not supply the constructor's license.
\end{proof}

\begin{corollary}[Central half-turn and antipodal clock operation]
\label{p3a:cor:clock-half-turn-antipode}
On the \(\mathsf{HS}_t+(\mathrm{CAL}_t)\) domain, the primitive-preserving
provenance map satisfies
\[
 \gamma_{tp}^{\rm op}\!\left(\exp_{H_{\rm succ}}(\tfrac12 e_t)\right)=-1
 \in SU(2)^{\rm op}.
\]
Consequently the clocked Hopf seed obeys the exact finite relations
\[
 \boxed{\Phi^H_{T_t/2}(x)=-x,\qquad \Phi^H_{T_t}(x)=x.}
\]
If \(E_k^x\) denotes the scalar round-\(S^3_x\) harmonic band of degree
\(k\), then
\[
 \boxed{\Phi^H_{T_t/2}|_{E_k^x}=(-1)^k.}
\]
Thus the half-clock operation, the vector-parent center element, and spatial
antipodal harmonic parity are one generated finite-operation chain.
\end{corollary}

\begin{proof}
One primitive Hopf turn is \(2\pi\), hence its halfway element is the unique
order-two element \(-1\in SU(2)\).  The right-regular reciprocal action sends
that element to \(x\mapsto x(-1)=-x\).  Paper~II-D's antipodal harmonic
theorem acts on degree-\(k\) scalar harmonics by \((-1)^k\).
\end{proof}

\begin{corollary}[Clock comparison and joint-orbit monodromy]
\label{p3a:cor:clock-joint-orbit-monodromy}
Use the existing primitive transport \(\kappa_t\), opposite Hopf
inclusion \(\iota_H^{\mathrm{op}}\) and reciprocal action.  Their
comparison is
\[
 \gamma_{tp}^{\mathrm{op}}=\iota_H^{\mathrm{op}}\kappa_t^{-1},
 \qquad \beta_t=\rho_v^{(1)}\gamma_{tp}^{\mathrm{op}},
 \qquad \chi_n^t\circ\kappa_t=\chi_n^H.
\]
Under the conjugation readout of a nonparallel marked joint tuple,
the primitive Hopf half-turn traces the nontrivial loop of
\(\mathcal O=SO(3)\).  Its quadratic coefficient has winding \(-1\)
in the rotating frame convention, while the full primitive turn
traces that loop twice and has even root winding.
\end{corollary}
\begin{proof}
Primitive-preserving transport identifies the integer character
lattices, giving the third formula and hence the stated composite
action.  Lift the conjugation path to \(g(\theta)=e^{i\theta}\)
in \(SU(2)\).  At \(\theta=\pi\) the lift is \(-1\), so the
conjugation tuple has returned but its spin line has changed sign.
The comparison coefficient transforms as \(e^{-2i\theta}\).
It winds once negatively on this half-turn.  At \(2\pi\) the spin
lift returns to \(+1\), doubling the winding.  This is exactly the
sign line and square descent of
Theorem~\ref{p1:thm:joint-orbit-quadratic-descent}.
\end{proof}

The scalar dual is retained as \(e_E(e_t)=1\) with its primitive
Poincare phase line on \(S_t^1\times S_E^1\), as in the scalar
Born-residue theorem.  The above loop compares operations; it is
not an assertion that momentum follows a trajectory on its sphere.
The position antipode is a separate representation of the same
clock operation.  No energy-cycle map into the joint orbit, or
choice of either global scalar-cycle spin sign, follows here.



\begin{theorem}[Derived clocked-Hopf operational speed ceiling]
\label{p3:thm:gp1-speed-calibration}
Choose a round metric representative \(h\) of radius \(R_x>0\) on the
reciprocal position parent (so \(R_x=\Rcs\) on the calibrated notation of
\((\mathrm{CAL}_{xp})\)) and the clock calibration \((\mathrm{CAL}_{t})\) above.
For a history \(\Phi:I\times S^3_x\to S^3_x\) in
\(\mathscr O_{\rm prop}^{\rm phys}\), put
\[
 v(\Phi)
 :=\sup_{\substack{x\in S^3_x,\ s,t\in I,\ s\ne t}}
 \frac{d_h(\Phi_t x,\Phi_s x)}{|t-s|}.
\]
Then the seed is a faithful smooth nonstationary isometric history of constant
pointwise speed
\[
 c_H:=\frac{2\pi R_x}{T_t},
\]
every generated history obeys
\begin{equation}
\label{p3:eq:gp1-calibrated-speed-bound}
 d_h(\Phi_t x,\Phi_s x)\le c_H|t-s|,
\end{equation}
and the sharp unlabelled operational ceiling is
\[
 \boxed{
 c_{\rm op}:=\sup_{\Phi\in\mathscr O_{\rm prop}^{\rm phys}}v(\Phi)
 =c_H=\frac{2\pi R_x}{T_t}.}
\]
Thus nonstationarity, bounded propagation content, and the value of the
ceiling relative to the spatial-radius and primitive-time calibrations are
derived.  Every point history in the mathematical relaxed curve family
\(\mathscr O_{\rm prop}^{\rm rel}\) is also \(c_H\)-Lipschitz.  The numerical witnesses for
\(R_x\) and \(T_t\), comparison with \(\Rmom\), and proof that a later field
channel has this undeformed principal owner remain separate.  Once that
principal ancestry is proved for an unlabelled neutral descendant,
Paper~II-A,
Corollary~\ref{p2a:cor:later-light-scalar-speed-inheritance}, forces its
front speed to equal \(c_{\rm op}\); no second neutral support rate may then
be inserted.  Deformed or medium-dependent effective channels may propagate
more slowly.
\end{theorem}

\begin{proof}
Because \(\kappa_t\) is an isomorphism, \(\iota_H^{\rm op}\) is injective,
and \(\rho_{\rm v}^{(1)}\) is faithful, \(\beta_t\) is a faithful smooth circle
action.  Pulling \eqref{p3a:eq:clocked-hopf-seed} back through \(\kappa_t\)
gives the provenance-fixed positive Hopf one-parameter subgroup
\[
 \exp_H\!\left(\frac{\widetilde t}{T_t}\tau_H^+\right).
\]
Its right-regular orbit through every point of the radius-\(R_x\) round
three-sphere is a great circle traversed once in time \(T_t\), hence has
constant speed \(2\pi R_x/T_t\).  In particular the seed is nonstationary and
attains \(c_H\).

The stationary history has rate zero.  Restriction and elapsed translation do
not increase a Lipschitz constant; isometric conjugation preserves it; and
the triangle inequality preserves the same bound under compatible
concatenation.  Therefore every member of the physical essential image obeys
\eqref{p3:eq:gp1-calibrated-speed-bound}.  Hence every physical velocity germ
lies in the closed norm ball of radius \(c_H\); convexity and closure put every
relaxed velocity germ in the same ball.  If
\(\gamma\in\mathscr O_{\rm prop}^{\rm rel}\), the metric length estimate for
absolutely continuous curves gives
\[
 d_h(\gamma(t),\gamma(s))
 \le \int_s^t |\dot\gamma(u)|_h\,du
 \le c_H|t-s|.
\]
This proves the auxiliary relaxed-family assertion.  The seed gives
the reverse inequality in the supremum, proving equality.  Paper~II's GP1
bounded-rate theorem is thereby discharged by an explicitly generated
antecedent rather than by an independently occurring propagation family.
\end{proof}

\begin{corollary}[Exact closed convex envelope of the physical clocked doctrine]
\label{p3:cor:gp1-round-causal-envelope}
Let \(\mathscr V_x^{\rm phys}\subset T_xS^3_x\) be the metric-derivative
velocity germs at \(x\) generated by
\(\mathscr O_{\rm prop}^{\rm phys}\).  Then
\[
 \overline{\operatorname{co}}(\mathscr V_x^{\rm phys})
 =\{v\in T_xS^3_x:h(v,v)\le c_{\rm op}^2\}.
\]
Equivalently, on the scalar cover the closed future operation cone is exactly
\[
 \mathcal C_x^+
 =\{(\tau,\xi):\tau\ge0,\ h(\xi,\xi)\le c_{\rm op}^2\tau^2\}.
\]
The boundary sphere is attained by the clocked Hopf seed and its physical
spatial-isometry conjugates.  The interior is the mathematical convex/closed
relaxation represented by \(\mathscr O_{\rm prop}^{\rm rel}\).  Thus the cone
is the exact closed convex kinematic envelope of the co-realized doctrine; it
does not assert physical occurrence of every interior history or velocity.
\end{corollary}

\begin{proof}
Theorem~\ref{p3:thm:gp1-speed-calibration} places every velocity germ in the
closed \(h\)-ball of radius \(c_{\rm op}\).  The seed supplies one germ of
norm \(c_{\rm op}\) at every point.  The round-isometry isotropy from
Paper~II, Theorem~\ref{p2:thm:rank-one-actual-isotropy}, acts transitively on
the unit tangent sphere, so spatial-isometry conjugation supplies every germ
with norm \(c_{\rm op}\).  The stationary history supplies zero.  By the
definition of closed convex hull, the envelope therefore contains the convex
hull of the boundary sphere, which is the whole
closed ball.  The speed bound gives the reverse inclusion.  Homogeneous
rescaling by \(\tau\ge0\) gives the displayed future cone.
\end{proof}

\begin{corollary}[Three-dimensional causal determinant]
\label{p3a:cor:v177-three-dimensional-causal-determinant}
On the exact round causal-envelope domain of
Corollary~\ref{p3:cor:gp1-round-causal-envelope}, the spatial tangent carrier
has real dimension three.  Define
\[
 C_c:=c_{\rm op}\,\operatorname{id}_{T_xS_x^3}.
\]
Then
\begin{equation}
\label{p3a:eq:v177-three-dimensional-causal-determinant}
 \boxed{J_c:=\det C_c=c_{\rm op}^{3}.}
\end{equation}
Equivalently, for the metric unit ball \(B_1\subset T_xS_x^3\) and
\(B_c=c_{\rm op}B_1\),
\[
 \frac{\operatorname{Vol}_h(B_c)}{\operatorname{Vol}_h(B_1)}=c_{\rm op}^{3}.
\]
Thus the cubic power is the determinant/volume character of isotropic causal
propagation on the generated three-dimensional spatial carrier; it is not an
independently fitted speed exponent.
\end{corollary}
\begin{proof}
For a scalar multiple \(aI\) on a three-dimensional real vector space,
\(\det(aI)=a^3\).  The volume statement is the determinant transformation
law.  Corollary~\ref{p3:cor:gp1-round-causal-envelope} supplies the isotropic
carrier and the radius \(c_{\rm op}\).
\end{proof}

\begin{theorem}[Exact polar dual of the GP1 future operation cone]
\label{p3:thm:gp1-operation-generator-polar-dual}
At a point \(x\in S^3_x\), use the same transported metric representative
and generated clocked-Hopf doctrine as
Theorem~\ref{p3:thm:gp1-speed-calibration}.  Put
\[
 W_x:=\mathbb R\oplus T_xS^3_x,
 \qquad
 W_x^*:=\mathbb R\oplus T_x^*S^3_x,
\]
and pair a scalar/vector operation increment \((\tau,\xi)\) with a
scalar/covector generator \((E,p)\) by
\[
 \langle(E,p),(\tau,\xi)\rangle
 :=E\tau-p(\xi).
\]
For the closed future operation cone
\[
 \mathcal C_x^+
 =\{(\tau,\xi):\tau\ge0,\ \|\xi\|_h\le c_{\rm op}\tau\},
\]
its positive polar cone is exactly
\[
 (\mathcal C_x^+)^\vee
 =\{(E,p):E\ge c_{\rm op}\|p\|_{h^{-1}}\},
\]
where
\[
 \|p\|_{h^{-1}}
 :=\sup_{\|\xi\|_h\le1}p(\xi).
\]
Moreover the duality is involutive:
\[
 \bigl((\mathcal C_x^+)^\vee\bigr)^\vee=\mathcal C_x^+.
\]
Thus the constituted local timed-operation geometry canonically determines a
crossed scalar/vector generator dual of type
\[
 (\tau,\xi)\longleftrightarrow(E,p).
\]
The boundary relation
\[
 E=c_{\rm op}\|p\|_{h^{-1}}
\]
is the boundary of this polar cone.  It is not, by itself, a declaration that
every physical field mode lies on that boundary.
\end{theorem}

\begin{proof}
Suppose first that \(E\ge c_{\rm op}\|p\|_{h^{-1}}\).  For every
\((\tau,\xi)\in\mathcal C_x^+\),
\[
 p(\xi)
 \le \|p\|_{h^{-1}}\|\xi\|_h
 \le c_{\rm op}\tau\|p\|_{h^{-1}}
 \le E\tau,
\]
so \(E\tau-p(\xi)\ge0\).  Hence the displayed generator cone is contained in
the positive polar.

Conversely, let \((E,p)\) lie in the positive polar.  In finite dimension the
dual norm is attained, so choose \(u\in T_xS^3_x\) with
\(\|u\|_h=1\) and \(p(u)=\|p\|_{h^{-1}}\).  For any \(\tau>0\), the
boundary vector \(\xi=c_{\rm op}\tau u\) lies in \(\mathcal C_x^+\).
Polar positivity gives
\[
 0\le E\tau-p(c_{\rm op}\tau u)
 =\tau\bigl(E-c_{\rm op}\|p\|_{h^{-1}}\bigr),
\]
hence \(E\ge c_{\rm op}\|p\|_{h^{-1}}\).  This proves the polar formula.

For the bipolar statement, let \((\tau,\xi)\) be nonnegative on every
\((E,p)\) in the displayed polar cone.  Taking \(p=0\), \(E>0\), gives
\(\tau\ge0\).  For any \(p\) with \(\|p\|_{h^{-1}}\le1\), choose the
boundary value \(E=c_{\rm op}\).  Then
\[
 c_{\rm op}\tau-p(\xi)\ge0.
\]
Taking the supremum over the dual unit ball gives
\(\|\xi\|_h\le c_{\rm op}\tau\).  Thus
\((\tau,\xi)\in\mathcal C_x^+\); the reverse inclusion follows from the
first half of the proof.  Hence the cone is exactly its bipolar.
\end{proof}

\begin{theorem}[Determinant root of the equivariant causal-cone group]
\label{p3a:thm:causal-cone-determinant-root-isomorphism}
Let
\[
 \mathcal G_x^0
 :=\operatorname{Aut}_{O(3)}^0
       \bigl(W_x^*,(\mathcal C_x^+)^\vee\bigr)
\]
be the identity component of the group of linear, \(O(3)\)-equivariant
automorphisms of the positive generator cone in
Theorem~\ref{p3:thm:gp1-operation-generator-polar-dual}.  Then
\[
 \det:\mathcal G_x^0\longrightarrow\mathbb R_{>0}
\]
is a Lie-group isomorphism, with inverse
\[
 H\longmapsto H^{1/4}I_{W_x^*}.
\]
Consequently a positive determinant character \(H\) has a unique
identity-connected equivariant rank-four causal action, namely the scalar
action \(H^{1/4}I\).
\end{theorem}

\begin{proof}
The scalar summand of \(W_x^*=\mathbb R\oplus T_x^*S_x^3\) is trivial and
the cotangent summand is the irreducible real standard \(O(3)\)-module.
Equivariance therefore forces every linear automorphism to have the form
\[
             (E,p)\longmapsto(aE,bp),
             \qquad a,b\in\mathbb R^\times.
\]
Preservation of the cone and its boundary gives \(a>0\) and \(|b|=a\).
The identity component is the branch \(b=a>0\), so its elements are exactly
\(aI_{W_x^*}\).  Since \(\dim W_x^*=4\), their determinant is \(a^4\).
The positive fourth root is smooth and multiplicative, proving the claimed
isomorphism and inverse.
\end{proof}

\begin{corollary}[Quadratic descendant of the causal determinant root]
\label{p3a:cor:causal-root-quadratic-half-character}
Under Theorem~\ref{p3a:thm:causal-cone-determinant-root-isomorphism}, put
\(U=W_x^*\) and let \(Q\in\operatorname{Sym}^2U\) be an already
retained quadratic tensor descendant of the causal-generator representation.
The direct associated symmetric-square action is
\[
 \boxed{\operatorname{Sym}^2(H^{1/4}I_U)Q=H^{1/2}Q.}
\]
Thus the tensor square carries the positive determinant half-character,
canonically within this generated associated representation. It does not
identify an independently constituted curvature-response or action
coefficient line with that descendant.
\end{corollary}
\begin{proof}
On a decomposable symmetric tensor,
\[
 (H^{1/4}u)\odot(H^{1/4}v)=H^{1/2}(u\odot v).
\]
Such tensors span \(\operatorname{Sym}^2U\), proving the formula.
The preceding theorem makes the positive quarter-root action unique.
A tensor in \(\operatorname{Sym}^2U\) pairs with two elements of
\(U^*\), not two elements of \(U\). A bilinear form
\(B\in\operatorname{Sym}^2U^*\) has the same numerical weight under
pullback \(B(Au,Av)\), whereas the inverse-dual representation induces
\(H^{-1/2}\). These are distinct typed operations; the present corollary
uses the direct tensor action.
\end{proof}

\begin{remark}[Bundle and occurrence scope]
\label{p3a:rem:determinant-root-bundle-scope}
The construction is fibrewise natural under the already constituted
orthogonal frame changes, hence assembles over the local generator bundle.
It classifies how a supplied positive determinant character acts; it does not
by itself supply a response-line occurrence or identify that character with
any later physical owner.
\end{remark}

\begin{lemma}[Causal homothety and calibration independence]
\label{p3a:lem:causal-homothety}
Let \(h_\alpha=\alpha^2h_0\) with \(\alpha>0\), and compute the operational
ceiling from the same elapsed-time unit and generated clocked-Hopf doctrine.  Then
\(c_\alpha=\alpha c_0\), and
\[
 h_\alpha(\xi,\xi)\le c_\alpha^2\tau^2
 \quad\Longleftrightarrow\quad
 h_0(\xi,\xi)\le c_0^2\tau^2.
\]
Moreover
\[
 c_\alpha\|p\|_{h_\alpha^{-1}}
 =c_0\|p\|_{h_0^{-1}},
 \qquad
 -c_\alpha^2d\widetilde t^{\,2}+h_\alpha
 =\alpha^2\bigl(-c_0^2d\widetilde t^{\,2}+h_0\bigr).
\]
Hence the operation cone, its positive polar, Lorentzian sign, and conformal
class do not require \((\mathrm{CAL}_{xp})\).  That calibration survives only
for dimensionful cross-role claims, such as a numerical radius or the ratio
\(R_x/\Rmom\).
\end{lemma}

\begin{proof}
Every spatial distance in the rate definition scales by \(\alpha\), so its
supremum scales by the same factor.  The dual norm scales by
\(\alpha^{-1}\).  Substitution gives all three displayed identities.
\end{proof}

\begin{remark}[Local polar duality and the derived generator-germ localization]
\label{p3:rem:polar-dual-not-global-pair}
Theorem~\ref{p3:thm:gp1-operation-generator-polar-dual} is a local theorem
about the dual of the generated timed-operation cone.  It does not identify
the global compact reciprocal role \(S^3_B\times S^1_B\) with \(T^*U\), and
it does not manufacture a pointwise graph relation between independently
covariant compact vector factors.  The local generator/readout realization on
\(T^*U\), however, is no longer a separate bridge: it is the generated
co-realization of Theorem~\ref{p3a:thm:canonical-generator-germ-localization}.
Paper~II-C's first-symbol pairing, all-band shared-generator action and
scalar-harmonic bridge
\(J_{xp}f=f\circ\iota_1^{-1}\) are already generated, the latter by
Corollary~\ref{p2c:cor:harmonic-pairxp-generated}. It is unitary for the
transported metric and measure and intertwines scalar Laplacians and complete
bands. Proposition~\ref{p2c:prop:v233-unit-and-positive-maps} additionally
fixes the finite unit-only map \(z\mapsto zI\). What remains separate is
a specified stronger observable-to-operator map beyond these outputs and any
continuum/local--global comparison; the positive extensions proved upstream
are explicitly nonunique.  By
Lemma~\ref{p3a:lem:causal-homothety}, reciprocal scale calibration is not a
premise of either the cotangent germ or cone/polar theorem; it returns only for
dimensionful cross-role comparisons.
\end{remark}

\begin{remark}[Notation for \(dt\) and \(dE\) on the scalar circles]
On the compact scalar torus, the symbols \(dt\) and \(dE\) are local Darboux
representatives of the globally defined circle one-forms
\[
        \Theta_t=\frac{\tper}{2\pi}d\theta_t,
        \qquad
        \Theta_E=\frac{\Eper}{2\pi}d\theta_E .
\]
Thus expressions such as \(\omega_{\rm sc}\) on \(S^1_t\times S^1_E\) mean
\(\Theta_E\wedge\Theta_t\).  On the universal covers
\(\widetilde t,\widetilde E\in\mathbb R\), the ordinary exact differentials
\(d\widetilde t\) and \(d\widetilde E\) are used.
\end{remark}

\begin{remark}[Why the scalar compact role space is a torus]
The scalar torus is not selected by compactifying an otherwise unstructured
symplectic plane.  Paper~II first derives the scalar A-role
\(A_{\rm s}\cong U(1)\).  Pullback contravariance gives the opposite test action.
Theorem~\ref{p3:thm:reciprocal-status-discharge} certifies both GP2 antecedents;
Paper~II-A's GP2 theorem objectifies
\(B_{\rm s}\cong A_{\rm s}^{\rm op}\cong U(1)\) and fixes the universal
relation-free joint doctrine.  Their joint role space is therefore
\[
        A_{\rm s}\times B_{\rm s}\cong S^1_t\times S^1_E.
\]
The Poincar\'e lattice pairing then supplies its primitive integral phase
class.  Other compact surfaces require a different generated datum.
\end{remark}


\begin{theorem}[No scalar metric modulus in the adopted phase datum]
\label{p3:thm:scalar-pair-no-modulus}
Assume $(\mathrm{CAL}_{\hbar})$.  Let the scalar completion be the compact canonical phase pair
\[
(t,E)\in S^1_t\times S^1_E,
\qquad \tper\Eper=2\pi\hbar .
\]
Then, inside the later spectral constitution \(\Dcl\), provided no scalar
circle metric is added, the scalar sector contains no independent running
metric modulus.  Its invariant phase datum is the cell \(2\pi\hbar\).
This is a statement about the declared datum: topology and the symplectic form
alone do not forbid adding an aspect-ratio metric modulus in a larger theory.
\end{theorem}

\begin{proof}
In normalized angular coordinates
\(\theta=2\pi t/\tper\) and \(\phi=2\pi E/\Eper\), the calibrated form is
\[
 \omega_{\rm sc}=\frac{\hbar}{2\pi}\,d\phi\wedge d\theta .
\]
It depends only on the product
\(\tper\Eper=2\pi\hbar\) and is invariant under
\((\tper,\Eper)\mapsto(\lambda\tper,\lambda^{-1}\Eper)\).  The declared
spectral constitution contains the dimensionless integer generators, but no
scalar-circle metric or dimensionful derivative that could distinguish the
two periods separately.  Hence their ratio is presentation freedom, not a
running modulus of the adopted phase datum.  The detailed normalization and
the boundary created by adjoining a scalar metric are recorded in
Paper~III-G, Appendix~\ref{p3a:app:scalar-no-modulus-details}.
\end{proof}

\begin{corollary}[CAL\(_t\)-adapted scalar periods and calibrated \(M_8\) scale chain]
\label{p3a:cor:calibrated-scalar-m8-scale-chain}
On the domain of Theorem~\ref{p3:thm:scalar-pair-no-modulus}, add
\((\mathrm{CAL}_t)\), which assigns the physical duration \(T_t\) to one
positive primitive succession turn.  Fix the resulting scalar-presentation
freedom by the CAL\(_t\)-adapted representative
\[
 \boxed{\tper:=T_t.}
\]
Then the primitive Poincar\'e cell gives
\[
 \boxed{\Eper=\frac{2\pi\hbar}{T_t}.}
\]
Together with the clocked-Hopf ceiling
\(c_{\rm op}=2\pi R_x/T_t\),
\[
 \boxed{\Eper R_x=\hbar c_{\rm op}.}
\]
If in addition \((\mathrm{CAL}_{xp})\) is used so that
\(R_x=\alpha_{xp}\Rmom\) and \(\Punit=\hbar/\Rmom\), then
\[
 \boxed{\alpha_{xp}\Eper=c_{\rm op}\Punit.}
\]
These are calibrated normal-form identities; before \((\mathrm{CAL}_t)\),
\(\tper\) and \(\Eper\) still have the reciprocal presentation freedom of
Theorem~\ref{p3:thm:scalar-pair-no-modulus}.
\end{corollary}

\begin{corollary}[Momentum-unit normal form]
\label{p3a:cor:momentum-unit-normal-form}
The calibrated scalar chain obeys
\[
 \boxed{\frac{c_{\rm op}T_t}{R_{\rm mom}}=2\pi\alpha_{xp},\qquad
 \frac{E_{\rm per}R_{\rm mom}}{\hbar c_{\rm op}}=\alpha_{xp}^{-1}.}
\]
Thus \(\mathrm{CAL}_{xp}+\mathrm{CAL}_t\) introduces no second dimensional
geometric modulus.  On any later branch that independently proves
\(R_x=R_{\rm mom}\), these reduce to
\(T_t=2\pi R_{\rm mom}/c_{\rm op}\) and
\(E_{\rm per}=\hbar c_{\rm op}/R_{\rm mom}\).  No such equality is asserted
for a generic reciprocal-position radius in Paper~III-A.
\end{corollary}



\begin{theorem}[Massless transverse gauge realization of the scalar energy lattice]
\label{p3a:thm:v245-massless-gauge-clock-lock}
Assume the calibrated branch of
Corollary~\ref{p3a:cor:calibrated-scalar-m8-scale-chain}, so that
\[
 \Eper=\frac{2\pi\hbar}{T_t},
 \qquad
 \Eper R_x=\hbar c_{\rm op}.
\]
Let a later action-owned stationary source-free transverse gauge sector have
positive-frequency Hamiltonian
\[
 H_0=\hbar c_{\rm op}\sqrt{\Delta_1^{\rm co}}
\]
on the round reciprocal-position three-sphere of physical radius \(R_x\).
Paper~I-A gives
\[
 \sqrt{\Delta_1^{\rm co}}\,Y_{L,\pm}
 =\frac{L+1}{R_x}Y_{L,\pm},\qquad L\ge1.
\]
Writing \(n:=L+1\in\{2,3,\ldots\}\) and
\[
 N_x:=R_x\sqrt{\Delta_1^{\rm co}},
 \qquad
 N_t:=\frac{iT_t}{2\pi}\partial_t
\]
on positive-frequency solutions, one has
\[
 \boxed{H_0=\Eper N_x},\qquad
 \boxed{N_t=N_x},\qquad
 \boxed{E_n=n\Eper}.
\]
Consequently the complete free massless transverse tower satisfies
\[
 \boxed{U_0(T_t)=e^{-iH_0T_t/\hbar}=I}.
\]
Thus the already constituted scalar character lattice becomes an actual
Hamiltonian energy lattice on this action-generated massless transverse gauge
branch.  This theorem does not make an arbitrary massive or interacting
Hamiltonian periodic and does not identify the compact scalar circle with
unbounded elapsed universal-cover time.
\end{theorem}

\begin{proof}
The coexact spectrum gives \(N_xY_{L,\pm}=(L+1)Y_{L,\pm}=nY_{L,\pm}\).
The calibrated identity \(\Eper R_x=\hbar c_{\rm op}\) then gives
\[
 H_0Y_{L,\pm}
 =\frac{\hbar c_{\rm op}n}{R_x}Y_{L,\pm}
 =n\Eper Y_{L,\pm}.
\]
On positive-frequency solutions \(i\hbar\partial_t=H_0\), and
\(T_t\Eper=2\pi\hbar\), so \(N_t=H_0/\Eper=N_x\).  Finally
\(e^{-iH_0T_t/\hbar}=e^{-2\pi iN_x}=I\) because \(N_x\) has integer spectrum.
\end{proof}

\begin{remark}[Second quantization and the vacuum-energy phase]
On a free bosonic Fock realization with
\(H_F=d\Gamma(H_0)+E_0I\), the same calculation gives
\[
 U_F(T_t)=e^{-iE_0T_t/\hbar}I.
\]
Hence rays and observables recur exactly even if an arbitrary scalar vacuum
energy reference is retained; literal identity follows for the normal-ordered
choice \(E_0=0\).  Neither statement asserts recurrence of the full interacting
Yang--Mills Hamiltonian.
\end{remark}



\begin{definition}[Affine spectral defect relative to the primitive clock]
\label{p3a:def:v247-affine-spectral-defect}
On the calibrated stationary round branch of
Corollary~\ref{p3a:cor:calibrated-scalar-m8-scale-chain}, let a free resolved
mode have physical one-particle energy \(E_a\) and let \(\nu_a>0\) denote the
affine harmonic number fixed by its round-\(S^3\) representation branch.  Define
\[
 \boxed{\mathfrak D_a
 :=\frac{E_a^2}{\Eper^2}-\nu_a^2
 =\frac{R_x^2\omega_a^2}{c_{\rm op}^2}-\nu_a^2.}
\]
We call \(\mathfrak D_a\) the \emph{affine spectral defect}; it is a physical
clock defect only on branches carrying the present \((\mathrm{CAL}_t)\) and
round-radius calibration.  Exact conformal spectral matching means
\(\mathfrak D_a=0\).  This definition does not assert that every massless
sector has zero defect; the value remains sector- and action-dependent.
\end{definition}

\begin{theorem}[Conformal scalar and Dirac realization of the primitive clock]
\label{p3a:thm:v247-conformal-clock-hierarchy}
Assume the calibrated stationary Einstein-cylinder branch with
\(\Eper R_x=\hbar c_{\rm op}\) and four-dimensional scalar curvature
\(R_4=6/R_x^2\).

First, restrict the native scalar action of
Paper~III-F, Theorem~\ref{p3f:thm:v175-native-local-scalar-action}, to a
canonically normalized stable scalar Hessian eigendirection whose
\emph{noncurvature} algebraic contribution is
\(m_A^2c_{\rm op}^4/\hbar^2\).  Because that action carries the fixed
conformal term \(R_4\Phi^2/6\), its free stationary operator is
\[
 \boxed{K_{\phi,A}
 =\partial_t^2+c_{\rm op}^2\!\left(\Delta_0+\frac{R_4}{6}\right)
 +\frac{m_A^2c_{\rm op}^4}{\hbar^2}.}
\]
For \(\Delta_0Y_k=k(k+2)Y_k/R_x^2\),
\[
 \boxed{E_{\phi,A,k}^2
 =(k+1)^2\Eper^2+m_A^2c_{\rm op}^4.}
\]
Hence the massless conformal member \(m_A=0\) has
\(\nu_k=k+1\), \(\mathfrak D_k=0\), and
\[
 \boxed{U_\phi(T_t)=I.}
\]
The nonzero lowest energy \(\Eper\) on finite \(S^3\) is a conformal-curvature
spectral gap, not a rest mass.

Second, restrict the local uninserted fermion jet of Paper~III-F,
Theorem~\ref{p3f:thm:v177-local-clocked-fermion-jet}, to a member whose
Yukawa mass vanishes on the same stationary background.  Paper~I-A,
Theorem~\ref{p1:thm:representation-harmonic-package}, gives the spatial
Dirac eigenvalues \(\pm(k+\tfrac32)/R_x\).  Therefore the positive-energy
one-particle levels are
\[
 \boxed{E_{D,k}=\left(k+\frac32\right)\Eper,\qquad k\ge0,}
\]
and the complete massless one-particle Dirac evolution obeys
\[
 \boxed{U_D(T_t)=-I,\qquad U_D(2T_t)=I.}
\]
Thus the primitive clock distinguishes integer conformal towers from the
half-integer spinor tower through generated spectral monodromy rather than an
inserted temporal sign convention.

For the transverse gauge branch of
Theorem~\ref{p3a:thm:v245-massless-gauge-clock-lock},
\(\nu_n=n\) and
\[
 \boxed{\mathfrak D_{n,X}
 =\left(\frac{m_Xc_{\rm op}^2}{\Eper}\right)^2.}
\]
Consequently zero mass is equivalent to complete primitive
clock locking on that transverse gauge sector.  The broader
criterion is \(\mathfrak D=0\), not masslessness alone.
\end{theorem}

\begin{proof}
For the scalar, the normalized Hessian of
Paper~III-F, Theorem~\ref{p3f:thm:v179-normalized-scalar-hessian}, contains
the curvature contribution \(+R_4/6\) as a scalar internal identity.  It may
therefore be separated from the noncurvature eigenvalue on any fixed
normalized eigendirection.  Since
\(k(k+2)+1=(k+1)^2\) and
\(\hbar c_{\rm op}/R_x=\Eper\), the scalar formula follows.  For the
massless Dirac branch, separation on the static product identifies the
positive Hamiltonian with \(\hbar c_{\rm op}|D_{S^3}|\); the round Dirac
spectrum therefore gives the displayed half-integer ladder.  Multiplication
by \(T_t/\hbar\), together with \(T_t\Eper=2\pi\hbar\), gives phase \(+1\)
for the integer scalar tower and phase \(-1\) for every half-integer Dirac
level.  The final gauge formula is the transverse spectrum already proved in
Paper~V-C rewritten using Definition~\ref{p3a:def:v247-affine-spectral-defect}.
\end{proof}

\begin{remark}[Sector scope of spectral matching]
\label{p3a:rem:v247-conformal-matching-scope}
The statement \(\mathfrak D=0\) is a spectral property of a specified free
operator.  It neither selects a scalar/Yukawa member nor promotes compact
clock characters into a universal periodic-time axiom.  In particular,
Paper~VI-A exhibits a massless tensor-gravity branch with nonzero geometric
defect.  A minimally coupled scalar comparison likewise differs from the
native conformally coupled scalar and is not substituted into the present
positive chain.
\end{remark}

\begin{corollary}[Calibrated normalized local \(M_8\) cone/evaluation form]
\label{p3a:cor:normalized-m8-cone-evaluation}
On the preceding calibrated domain, define
\[
 \vartheta:=\frac{2\pi\tau}{T_t},\qquad X:=\frac{\xi}{R_x},\qquad
 \varepsilon:=\frac{E}{\Eper},\qquad P:=\frac{p}{\Punit}.
\]
Then the local operation/generator evaluation becomes
\[
 \boxed{\frac1\hbar\langle(E,p),(\tau,\xi)\rangle
 =\varepsilon\vartheta-\alpha_{xp}P(X),}
\]
while the operation cone and its positive polar are
\[
 \boxed{\vartheta\ge0,\quad \|X\|\le\vartheta,\qquad
 \varepsilon\ge\alpha_{xp}\|P\|.}
\]
Equivalently, with the dimensionless covector
\(\bar P_x:=\alpha_{xp}P=(R_x/\hbar)p\),
\[
 \frac1\hbar\langle q,v\rangle=\varepsilon\vartheta-\bar P_x(X),\qquad
 \|X\|\le\vartheta,\qquad \varepsilon\ge\|\bar P_x\|.
\]
The notation $\Pi_i^{(a)}$ is reserved for the dimensionful finite spatial
transport operators of Paper~II-D and Paper~VII-C.
No equality \(\alpha_{xp}=1\) follows; such a unit-identification would require
an additional normalization bridge.
\end{corollary}

\begin{theorem}[Complete simultaneous zero-defect locus for the Paper-III-A interfaces]
\label{p3a:thm:complete-joint-zero-defect-locus}
Fix the frozen Paper-II role package, the actual provenance datum
\(\mathsf{HS}_t\), and an activation pattern (not prescribed numerical
members) for whichever of
\(\mathrm{CAL}_t\), \(\mathrm{CAL}_{\hbar}\),
\(\mathrm{CAL}_{xp}\), \(\mathrm{HIST}_x^{\rm fin}\), and
\(\mathsf P_{\rm mc}\) are explicitly active.  Relative to the interface
search basis used in this paper, let a candidate tuple record
\[
 Q=(S_{\rm v},S_{\rm s},\epsilon_t,N_{\rm sc},m_p,m_x^{\rm cand},
     n_{\rm clk},\mathscr O; \text{active calibration data}),
\]
where \(S_{\rm v}\subseteq A_{\rm v}\times B_{\rm v}\) and
\(S_{\rm s}\subseteq A_{\rm s}\times B_{\rm s}\) are proposed joint
supports, \(m_x^{\rm cand}\) is the candidate reciprocal-position
least-resolution label, \(\epsilon_t\in\{\pm1\}\) is the scalar primitive relative to the
Paper-I-A positive primitive, \(N_{\rm sc}\in\mathbb Z\) multiplies the
Poincar\'e class, \(n_{\rm clk}\in\mathbb Z\) is the clock winding, and
\(\mathscr O\) is a proposed physical history doctrine.  Let
\(\Delta_{3A}(Q)\) be the standard simultaneous GP2 defect vector
\[
 \Delta_{3A}(Q)
 =(D_{\rm ins},D_{\rm omit},D_{\rm rel},D_{\rm type},D_{\rm clos}),
\]
where provenance and calibration compatibility are recorded in the relation,
type, and closure components rather than treated as new geometric principles.

Then the joint zero-defect locus
\(\mathfrak Z_{3A}:=\{Q:\Delta_{3A}(Q)=0\}\) is nonempty and consists
exactly of the tuples satisfying all applicable conditions
\begin{align}
 S_{\rm v}&=A_{\rm v}\times B_{\rm v},
 &S_{\rm s}&=A_{\rm s}\times B_{\rm s},
 \label{p3a:eq:joint-locus-products}\\
 \epsilon_t&=+1,
 &N_{\rm sc}&=1,
 &m_p&=1,
 \label{p3a:eq:joint-locus-provenance}\\
 n_{\rm clk}&=1,
 &\mathscr O&=\mathscr O_{\rm prop}^{\rm phys}.
 \label{p3a:eq:joint-locus-clock}
\end{align}
If no finite position history occurs, \(m_x^{\rm cand}\) remains the unresolved member
of the already derived position-stage family and is not selected.  If
\(\mathrm{HIST}_x^{\rm fin}\) occurs with frontier \(F_x\), joint
compatibility adds the necessary and sufficient condition
\begin{equation}
 m_x^{\rm cand}=m_x^{\rm res}(F_x).
 \label{p3a:eq:joint-locus-frontier}
\end{equation}
When the corresponding calibrations are active, the remaining joint
conditions are
\begin{equation}
 \tper=T_t,
 \qquad T_t\Eper=2\pi\hbar,
 \qquad R_x=\alpha_{xp}\Rmom,
 \qquad c_{\rm op}=\frac{2\pi R_x}{T_t};
 \label{p3a:eq:joint-locus-calibrations}
\end{equation}
with \(T_t,\Eper,\hbar,R_x,\Rmom,\alpha_{xp}>0\), as already required by
the native calibration sorts.  Conditions involving an inactive calibration are simply absent, rather than
silently assumed.  For a fiber in which numerical calibration members have
already been fixed externally, the fiber is empty unless these equations hold.
If \(\mathsf P_{\rm mc}\) is active, its physical essential image is exactly
the typed profile/evaluation/action package
\((B_A^{\rm prof},E,\rho_A)\); if it is inactive,
the full profile algebra/spectrum remains mathematical and its absence causes
no defect in the premise-minimal reciprocal role.

Thus coordinatewise admissibility is not sufficient.  For example, the
oppositely oriented scalar circle, a higher-winding circle action, a positive
but incompatible value of \(\Eper\), a graph-supported A/B pair, or an
arbitrary admissible position stage may each exist as a mathematical object
or in a weaker separately typed domain.  Combined with the frozen provenance,
relation-free product, calibrations, or an occurring frontier, respectively,
each lies off \(\mathfrak Z_{3A}\).  Only simultaneous combinations on the
displayed locus survive.  This is an exhaustive statement relative to the
listed Paper-III-A interface coordinates, not an absolute classification of
all possible later dynamics.
\end{theorem}

\begin{proof}
The vector and scalar support equalities are necessary and sufficient by the
two total evaluation constructors and Paper~II-A's relation-free reciprocal
closure: a proper support inserts a cross-role relation and omits generated
off-support evaluations.  The transported-orientation theorem forces
\(\epsilon_t=+1\), while the perfect lattice evaluation generates the
primitive class \(N_{\rm sc}=1\).  Paper~II-B supplies \(m_p=1\).
Definition~\ref{p3a:def:clocked-hopf-propagation-doctrine} and
Corollary~\ref{p3a:cor:clocked-hopf-corealization} give winding one and the
exact physical essential image
\(\mathscr O_{\rm prop}^{\rm phys}\); changing either creates provenance,
insertion, omission, or closure defect.

The Paper-II-C least-resolution theorem leaves the generic position family
unselected and, when a frontier occurs, evaluates its already declared
statistic to the unique compatible \(m_x^{\rm res}(F_x)\).  The calibrated
phase-cell lemma, the CAL\(_t\)-adapted normal form, reciprocal metric
calibration, and the clocked-Hopf speed theorem give precisely
\eqref{p3a:eq:joint-locus-calibrations}.  Finally, the physical-profile
firewall in Definition~\ref{p3:def:reciprocal-interface-hypotheses} proves the
last conditional statement.

All displayed conditions are realized together by the constructed reciprocal
products, transported primitives, primitive Poincar\'e class, frozen momentum
seed, clocked history constructor, and compatible members of every active
calibration role, so
the locus is nonempty.  Conversely each failed equality is exactly one of the
defects just identified.  Hence their conjunction is both necessary and
sufficient for \(\Delta_{3A}(Q)=0\) on the stated search basis.
\end{proof}

\begin{corollary}[Prospective neutral-channel extension of the joint locus]
\label{p3a:cor:joint-locus-neutral-speed-extension}
Extend a zero-defect tuple by a later unlabelled neutral field channel.  If
that later paper proves that its kinetic operator is an undeformed descendant
of the unique clocked-host principal owner, the extended tuple has zero
principal-ownership defect only if
\[
 c_{\rm front}=c_{\rm op}.
\]
Support ancestry alone is insufficient; a medium term, mass term, or other
lower-order deformation may change dispersion without changing this front
speed, while a genuine principal-symbol deformation defines a separately
typed branch and must remain explicit.
\end{corollary}

\begin{proof}
Apply Paper~II-A,
Theorem~\ref{p2a:thm:undeformed-host-principal-inheritance}, and
Corollary~\ref{p2a:cor:later-light-scalar-speed-inheritance}.  The theorem's
countermodel shows why unique support without the principal-owner hypothesis
would not suffice.
\end{proof}

\begin{remark}[Hopf-character and scalar dual-lattice inheritance]
\label{p3a:rem:hopf-character-energy-lattice}
The provenance isomorphism \(\kappa_t:U(1)_H\to S_t^1\) preserves the
primitive and therefore identifies the integer Hopf character lattice with
the integer Fourier/dual lattice of the succession circle.  If \(e_E\) is the
positive dual primitive, the charge label \(q\in\mathbb Z\) corresponds to
\(q e_E\) in that discrete dual lattice and
\[
 \chi_q(T_t/2)=(-1)^q.
\]
This is a lattice/character statement, not a pointwise identification inside
the compact \(S_E^1\) factor and not a physical energy-spectrum theorem.
Identifying \(q e_E\) with an actual Hamiltonian energy requires a later DYN
bridge.
\end{remark}

\begin{remark}[Interpretation status of the transported scalar labels]
\label{p3:rem:time-energy-upgrade}
On the \(\mathsf{HS}_t\) domain, \(t\) is the scalar-succession coordinate
with the transported Paper-I-A orientation, and \(E\) is its positively paired
full-dual coordinate.  The scalar completion supplies their compact phase
torus, local Darboux charts, and universal covers.  What remains later is
dimensionful calibration, Hamiltonian selection, the actual trajectory, and
any thermodynamic arrow; none follows from compact topology alone.
\end{remark}

\begin{remark}[Why the scalar closure is required]
\label{p3:rem:scalar-closure}
The scalar $S^1\times S^1$ sector is not introduced as an additional embellishment. It is the Born completion of the separately prior-generated scalar role, used after the reciprocal vector parent is shown not to contain the required full-product-intrinsic canonical scalar pair: the branch
\[
S^3_x\times S^3_p
\]
has no degree-two symplectic class and no nonconstant full-product-intrinsic scalar observables, so it cannot itself supply the scalar canonical pair needed for a completed Born/contact geometry.  It is not a claim that the compact product itself is a global symplectic manifold.
\end{remark}

\begin{theorem}[No global symplectic form on the complete compact geometry]
\label{p3:thm:full-compact-carrier-no-symplectic-a}
Let
\[
 \mathcal M_8=S^3_A\times S^3_B\times T^2
\]
be the compact manifold supplied by the premise-minimal Paper-II-C/GP2
reciprocal closure and the scalar completion.  Then
\(\mathcal M_8\) admits no symplectic form.  The primitive scalar two-form on
\(T^2\) is therefore a rank-two phase form, not a symplectic form on the full
eight-manifold.
\end{theorem}

\begin{proof}
The reciprocal closure and scalar completion identify the compact
manifold with the product in Paper~I-D,
Proposition~\ref{p1b:prop:lcs}. Its cohomological obstruction excludes a
symplectic form on that eight-manifold. The scalar factor still carries
the calibrated rank-two phase form of
Theorem~\ref{p3:thm:scalar-born-residue}.
\end{proof}

\subsection{Primitive scalar descent and the unchanged real cover}
\label{p3a:sec:scalar-cover-faithfulness}

\begin{corollary}[Unchanged primitive scalar descent has degree one]
\label{p3a:cor:primitive-scalar-unchanged-descent}
Retain the complete primitive scalar-pair datum of
Lemmas~\ref{p3:lem:minimal-phase-torus-quantization}
and~\ref{p3:lem:separable-scalar-compactification}, including its line,
orientation, calibrated phase cell and tests. Its scalar torus has
\[
 c_1=1,\qquad T_*E_*=2\pi\hbar.
\]
If it is claimed to descend unchanged through an orientation-preserving
finite cover \(q:T^2\to\overline T^2\) of degree \(d\), then \(d=1\).
An elapsed-real-time presentation retains the deck and line data of this
same owner.
\end{corollary}
\begin{proof}
For unchanged descent the primitive line is the pullback of a line on the
quotient. Integral Chern-class pairing gives
\(1=\langle q^*c_1,[T^2]\rangle
=d\langle c_1,[\overline T^2]\rangle\).
The second pairing is an integer, so a positive degree is one. Pullback to
a real cover presents sections with the inherited deck identifications;
forgetting that identification changes the retained datum rather than
compressing its unchanged descent.
\end{proof}

\begin{corollary}[Faithful marked clock on elapsed real time]
\label{p3a:cor:faithful-clock-real-cover}
On the fixed clocked Hopf-propagation doctrine of
Definition~\ref{p3a:def:clocked-hopf-propagation-doctrine}, a proposed constant
rate replacement \(\theta\mapsto\lambda\theta\) on the unchanged marked
clock circle descends only for integer \(\lambda\). A faithful such map
has \(|\lambda|=1\); the retained positive primitive orientation forces
\(\lambda=1\). Passing to the elapsed cover \(\mathbb R\) therefore does
not insert another rate multiplier into
\[
 \boxed{c_{\rm op}=2\pi R_x/T_t.}
\]
\end{corollary}
\begin{proof}
The deck translation \(\theta\mapsto\theta+2\pi\) is preserved precisely
when \(2\pi\lambda\) is an integer turn. The induced circle map has degree
\(\lambda\), and is injective precisely for degree \(\pm1\). The frozen
orientation and unit winding select \(+1\). The orbit-length calculation
in Theorem~\ref{p3:thm:gp1-speed-calibration} then retains the same speed.
\end{proof}

These results preserve the global line and clock calibration. They do not
identify a scalar phase-cell area, whose dimension is action, with proper
spatial volume or a maximum Hamiltonian energy. Physical compression bounds
consume the metric/action readouts evaluated in Paper~VI-A. Paper~III-G, Appendix~\ref{p3a:app:scalar-cell-energy-scope} gives the phase-patch and energy-period
diagnostics supporting this scope statement.


\subsection{Relative-solder localization of the local generator germ}
\label{p3a:sec:v187-m8-relative-solder-localization}
The complete parent coframe of Paper~III-B has
$\mathfrak k=\mathfrak{su}(2)\oplus\mathfrak u(1)$ and therefore dualizes to
$\mathfrak k^*$.  % PROPOSED REPLACEMENT of the vector-descent paragraph immediately before
\paragraph{Exact relative-solder dual normalization.}
Let $\mathfrak g=\mathfrak{su}(2)$ and let
$\mathfrak p=\{(X,-X):X\in\mathfrak g\}$ be the inversion-odd
complement of the diagonal stabilizer. The source II-C evaluation map
$\operatorname{ev}|_{\mathfrak p}(X,-X)=2X$ has the canonical inverse
$\tau_{\rm ev}(X)=(X/2,-X/2)$. For the \emph{relative solder}
coframe of Paper I-C, use the equally canonical anti-diagonal
inclusion
\[
 \boxed{\tau_{\mathcal E}(X)=(X,-X)=2\tau_{\rm ev}(X).}
\]
Since $\mathcal E=(\theta_x-\theta_p)/2$,
$\mathcal E|_{\mathfrak p}=\operatorname{ev}|_{\mathfrak p}/2$.
Let $\mathcal I_a:T_a^*A\to T_b^*S_B^3$ denote the II-C
first-symbol cotangent transport, including the evaluation-spectrum
diffeomorphism and the natural local left-translation tangent
identification. Define the vector-dual descent by
\[
 \boxed{\mathfrak g^*
 \xrightarrow{(\tau_{\mathcal E}^*)^{-1}}\mathfrak p^*
 \xrightarrow{\Phi_v=\mathcal I_a\circ(\operatorname{ev}^*)^{-1}}
 T_b^*S_B^3.}
\]
Here every inverse dual is taken on the three-dimensional first-symbol
fibres, not on the ambient six-dimensional $\mathfrak g\oplus\mathfrak g$.
The identity
$\tau_{\mathcal E}^*\circ\operatorname{ev}^*=2I_{\mathfrak g^*}$
proves the \emph{exact} relative-solder result
\[
 \boxed{\Phi_v\circ(\tau_{\mathcal E}^*)^{-1}
 =\tfrac12\mathcal I_a.}
\]
The original use of an unspecified $\tau$ is thereby replaced by
a fully normalized arrow; $\tau_{\rm ev}$ would instead produce
the full evaluation/Maurer--Cartan coframe and may not be silently
substituted for the relative solder. The scalar $\mathfrak u(1)^*$
normalization is independently inherited as printed. The local block
therefore still has ranks $3+1=4$ and the same determinant character
$a\mapsto a^4$; no physical clock, source response $H$, or G0
occurrence is inferred from the normalization.

The separately normalized relative scalar character and scalar Poincare dual
give the $\mathfrak u(1)^*$ descendant in $T^*S_t^1$.  Thus
\[
 \boxed{\Phi_{M8,U}:\mathfrak k^*
 \xrightarrow{\sim}T^*S_B^3\oplus T^*S_t^1=T^*U.}
\]

\begin{corollary}[Rank-four parent of the local determinant root]
\label{p3a:cor:v187-rank-four-parent-root}
The local generator block is the descended coframe of
$\mathcal E_{M8}$ and hence has rank four, while its vector child has rank
three. Paper~I-D, Proposition~\ref{p1b:prop:local-coframe}, now gives
the determinant character $a\mapsto a^4$ on this full local block.  Whenever the separately generated CRP4
owner supplies response character $H$, the already derived positive root is
\[
 a=H^{1/4}.
\]
This identifies the geometric parent of the exponent $1/4$; it does not
supply the response $H$, a clock orientation, or a physical history.
\end{corollary}



\section{Clocked Hopf dynamics and exact continuation}
\label{paper3a:sec:organized-continuation}
\subsection{Universal generated exact-continuation architecture}
\label{p3a:sec:universal-generated-continuation}

The universal object in this section is a single \emph{construction rule}
applied to sector-dependent parent data.  It is not one common state space,
clock, Hamiltonian, semigroup, or generator shared by every process.

\begin{definition}[Generated exact-continuation signature]
\label{p3a:def:generated-exact-continuation-signature}
Let \(D=(d_1,\ldots,d_n)\) be a jointly occurring/generated tuple with
its complete licensed roots in a frozen presentation with grammar \(\Sigma\).
A comparison specification, fixed before candidate evaluation, consists of
\[
 \Sigma_{\rm spec}(D)=
 (\mathcal S_D,X_D,\mathcal C_D^\tau,\Delta_D,
  \mathcal G_D,\operatorname{Cal}_D,\mathsf{DynGraph}_D^\tau),
\]
where \(\mathcal S_D\) is a generated succession category or preorder;
\(X_D\) assigns the GP1-finitely-presented objects \(X_s\);
\(\mathcal C_D^\tau\) declares the arrow types; \(\Delta_D\) contains
all prior generated comparisons, owner relations, equalities, inequalities,
defects, conservation and endpoint conditions; \(\mathcal G_D\) is the
unmarked presentation groupoid; \(\operatorname{Cal}_D\) is the trivial
or already generated calibration congruence; and
\(\mathsf{DynGraph}_D^\tau\) is the predeclared process-family/doctrine
output sort. The object assignment does not yet assign continuation arrows.

A \emph{generated exact-continuation signature} is the certified package
\[
 \boxed{\Sigma_{\rm cont}(D)=(\Sigma_{\rm spec}(D),\Omega_D),}
\]
where \(\Omega_D\) records a finite result-independent word \(w_D\)
satisfying Paper~II-A's
Definition~\ref{p2a:def:frozen-constructor-origin-certificate}. Its required
family output is precisely the complete exact continuation family in
Definition~\ref{p3a:def:exact-continuation-moduli}:
\[
 \operatorname{Orig}_{\Sigma}(w_D;D,\mathsf{DynGraph}_D^\tau),
 \qquad \operatorname{ReqOut}_{\Sigma}(w_D,D)
       =\operatorname{Cont}^{\rm ex}(\Sigma_{\rm spec}(D)).
\]
Every construction, functional-calculus, completion and retained intermediate
actually used is already licensed in its native sort. Structured-isomorphism
and calibration quotient steps are likewise licensed and fixed before the
quotient is formed; the family is retained as an intermediate if the final
word returns its quotient presentation. The complete retained family and
quotient must have the stated GP1 finite/exact, bounded presentation.

The datum \(\Omega_D\) records an existing license; it does not add a
universal solution constructor. Generated comparisons and a declared target
alone give a specification, not this certificate. Without it, the family
below remains a mathematical classification. This is the current upstream
constructor boundary, also made explicit in
Remark~\ref{p2a:rem:v233-native-constructor-boundary}.
\end{definition}

\begin{definition}[Exact continuation moduli]
\label{p3a:def:exact-continuation-moduli}
For a comparison specification define mathematically
\[
\begin{split}
 \operatorname{Cont}^{\rm ex}(\Sigma_{\rm spec}(D))
 :=\bigl\{\Phi:\mathcal S_D\longrightarrow\mathcal C_D^\tau:\;&
       \Phi(s)=X_s,\quad
       \Phi(\operatorname{id}_s)=\operatorname{id}_{X_s},\\
     & \Phi(\beta\alpha)=\Phi(\beta)\Phi(\alpha),\\
     & \Phi\text{ is }\mathcal G_D\text{-natural},\\
     & \Phi\models\Delta_D,\quad
       \mathfrak D_{\rm joint}^{\rm GP2}(D,\Phi)=0,\\
     & \Phi\text{ has a GP1-finite exact presentation}
       \bigr\}.
\end{split}
\]
For a certified package we use the consistent abbreviation
\[
 \operatorname{Cont}^{\rm ex}(\Sigma_{\rm cont}(D))
 :=\operatorname{Cont}^{\rm ex}(\Sigma_{\rm spec}(D)).
\]
The certificate changes the construction status, not the mathematical
membership conditions. Its structured moduli are
\[
 \mathfrak M_{\rm cont}(D)
 :=
 \operatorname{Cont}^{\rm ex}(\Sigma_{\rm spec}(D))
 \big/(\cong,\operatorname{Cal}_D).
\]
For a certified signature these family and quotient operations are the
licensed steps recorded by \(\Omega_D\). Only its already generated
calibrations are quotiented. Independently owned clocks or rates are not
identified merely because two formulas have the same shape.
\end{definition}

\begin{theorem}[Universal generated exact-continuation and selection theorem]
\label{p3a:thm:universal-generated-exact-continuation}
For every certified generated signature of
Definition~\ref{p3a:def:generated-exact-continuation-signature},
\(\mathfrak M_{\rm cont}(D)\) is the specialization of Paper~II-A,
Theorem~\ref{p2a:thm:generated-exact-solution-moduli}, to its declared
process-family sort. The licensed construction has the outcome
\[
 \boxed{\mathfrak U_{\rm dyn}(\Sigma_{\rm cont}(D))=
 \begin{cases}
 \bot_{\rm FRC},&\mathfrak M_{\rm cont}(D)=\varnothing,\\
 [\Phi_*],&[\Phi_*]\text{ is the unique GP3-least structured class},\\
 \operatorname{Irr}\mathfrak M_{\rm cont}(D),&
   \text{otherwise on the nonempty branch}.
 \end{cases}}
\]
On the nonempty, complete GP1/GP2 domain, the family/doctrine output is
co-realized with the same complete root tuple. GP3 selects its unique least
class only when one exists; it supplies neither a missing constructor nor an
absent candidate. No member of an unresolved process family, initial state
or actual history is thereby selected. History evaluation retains its own
licensed evaluation and actual root token.

Without \(\Omega_D\), the exact family and quotient are mathematical
classification only. Their emptiness is not obstruction of an unlicensed
physical constructor. Different certified sectors may have different
carriers, succession categories, arrows, generators and calibrations; the
universality is that of the architecture.

When joint admissibility factors coordinatewise, no generated cross-role
relation exists, and no newly enabled common constructor adds an obligation,
joint moduli are the relation-free product. A generated common-value
incidence contributes its equality while every other joint condition remains.
The complete fibre product requires redundancy of those other conditions;
a parent-image definition additionally requires the joint-lift criterion of
Lemma~\ref{p2a:lem:repair-joint-lift-completeness}.
\end{theorem}
\begin{proof}
Collect object assignment, identities, composition, naturality squares and
all frozen comparisons into the upstream diagram. Its exact solution family
is exactly \(\operatorname{Cont}^{\rm ex}(\Sigma_{\rm spec}(D))\).
The word and quotient licenses in \(\Omega_D\), together with the
complete GP1/GP2 realization domain, supply all antecedents of Paper~II-A's
exact-solution theorem. Its empty/least/unresolved alternatives, rooted
co-realization and joint-incidence statements give the conclusions. The
declared-output boundary separates family content from history tokens.
No extra dynamics principle or duplicate G0 root is needed on this certified
domain. Without the origin certificate, the physical application of the
upstream theorem has not been justified.
\end{proof}


\begin{definition}[Dynamics-atlas chronology]
\label{p3a:def:v236-dynamics-atlas-status}
For a certified continuation family, the status labels of
Definition~\ref{p2a:def:v236-atlas-status} are applied prospectively:
\(\mathrm{CND}\) denotes a classified continuation branch whose occurrence is
not asserted; \(\mathrm{GEN}\) a generated process-family output;
\(\mathrm{DL}\) a branch/member whose occurrence is established only by a
later independently constituted theorem; \(\mathrm{ROOT}\) an actual process,
initial-state or history member not discharged upstream; and
\(\mathrm{OBS}\) an empty/incompatible continuation branch.  A later
\(\mathrm{DL}\) theorem never changes an earlier unresolved
\(\operatorname{Irr}\mathfrak M_{\rm cont}\) into an earlier GP3 least class.
\end{definition}

\begin{corollary}[Common-content descent from unresolved dynamics]
\label{p3a:cor:v236-unresolved-common-content}
Let \(\mathcal F\subset\operatorname{Irr}\mathfrak M_{\rm cont}(D)\) be a
retained unresolved family.  Any licensed readout or generated descendant
which is equal on every member of \(\mathcal F\) descends to the family
without selecting a member.  A member-dependent trajectory, initial state or
history remains conditional until a later theorem or G0 token supplies it.
\end{corollary}


\begin{lemma}[Transport of a certified continuation family]
\label{p3a:lem:v166-continuation-certificate-transport}
Suppose a frozen word returns a retained family \(K_D\), and the same
grammar licenses a result-independent map \(N_D\), its required family
image, and the stated quotient. If
\[
 N_D(K_D)=\operatorname{Cont}^{\rm ex}(\Sigma_{\rm spec}(D)),
\]
then composition of these words certifies the complete continuation family
and quotient. An equivariant bijection respecting exactly the licensed
equivalences identifies their moduli.
\end{lemma}
\begin{proof}
Concatenate the finite words. Every input port is an existing retained
output, every operation and reuse is licensed, and the required target was
fixed before evaluation. The displayed equality proves completeness of the
image; it does not license the image operation. Equivariance and a compatible
inverse prove the quotient identification.
\end{proof}

\begin{corollary}[Rooted dynamical-parent cascade and exact stop test]
\label{p3a:cor:rooted-dynamical-parent-cascade}
Let \(D_\star\) be an occurring/generated root and
\(D_a=\mathsf F_a(D_\star)\) a finite family admitted by Paper~II-A,
Theorem~\ref{p2a:thm:rooted-generated-parent-cascade}. Each sector may
use the same continuation architecture after its own certificate
\(\Omega_{D_a}\) and complete realization domain have been supplied.
Its carrier, clock, arrow type and generator need not agree with another's.

A retained process output \(Q_a\) may generate a later parent
\[
 D_{a+1}=\operatorname{Cl}^{\min}_{\rm GP2}
       (D_a,\operatorname{EssIm}Q_a)
\]
without a second root on its licensed nonobstructed image. This promotion
uses a unique retained output or a licensed invariant common to the whole
surviving family, never an arbitrary unresolved member.

A nontrivial unique physical law cannot be claimed from this architecture
when a constructor certificate or complete realization is missing, the exact
family is empty, stationarity survives, the readout fails its structured
descent test, or the required nonstationarity/occurrence comparison is absent.
On a finite homogeneous linear branch, one positive generator direction gives
the family of allowed rates in
Theorem~\ref{p3a:thm:finite-generator-one-ray-rigidity}. Uniqueness by that
argument additionally requires a singleton quotient of those rates by the
actual structured isomorphisms and calibration congruence. Otherwise an
independent whole-candidate GP3 leastness proof is required; absence of a
listed obstruction is not itself such a proof. No common cross-sector rate
follows from shared root provenance.
\end{corollary}
\begin{proof}
Apply the rooted-parent theorem and the certified continuation theorem at
each admitted step. The identity-survival and common-invariant corollaries
give the stationary and family/member boundaries. The projectability theorem
gives the category-specific readout condition, and the finite-generator
theorem gives precisely the positive-direction and rate-quotient conclusion.
All conclusions use the complete sector data and its own licensed words.
\end{proof}



\subsection{Generated law, certified flow, and actual history}
\begin{definition}[Generated law, certified flow, and actual history]
\label{p3a:def:v149-law-flow-history}
For a generated exact-continuation problem, a \emph{law object} is any
already licensed relation, vector field, derivation, Hamiltonian/constraint
generator, cocycle generator, or finite linear generator retained in the
process signature before a particular history is chosen.  A \emph{flow
certification} is a proved mathematical theorem which, on its stated domain,
constructs a process family from that law object.  Its use is classified by
Paper~II-A's distinction between internal theorem closure and analytic
certification.  An \emph{actual history} is a member trajectory/process graph
obtained only after the relevant G0 initial/root token is present or after a
previously occurring root is conservatively transported by a licensed
solution constructor.
\end{definition}

\begin{theorem}[Law--flow--history separation]
\label{p3a:thm:v149-law-flow-history-separation}
Let \(D\) be a frozen parent and let a licensed construction retain a law
object \(\mathscr L_D\) in the declared dynamics signature.  Then:
\begin{enumerate}[label=(\roman*),leftmargin=2em]
\item retention of \(\mathscr L_D\) does not by itself assert existence,
      uniqueness, or actuality of a process flow;
\item if a proved generation theorem \(\mathfrak F\) constructs a process
      family \(\Phi_{t,s}=\mathfrak F(\mathscr L_D)\), then that use is an
      internal theorem closure only when it introduces no additional analytic
      certification hypotheses beyond the already frozen/generated
      mathematical structure.  If it requires regularity, domain, closedness,
      hyperbolicity, compactness, existence, uniqueness, continuous
      dependence, resolvent estimates, or analogous hypotheses not implied by
      correct provenance, its use is an analytic certification and those
      hypotheses must first be proved on the claimed domain;
\item once the certified process family exists, choosing an actual initial
      state or root token \(x_0\) and evaluating
      \(\gamma(t)=\Phi_{t,t_0}(x_0)\) is occurrence content unless the same
      token is already inherited by the conservative G0 witness-transport
      clause; and
\item neither theorem closure nor analytic certification selects a member of
      an unresolved initial-state/history family.
\end{enumerate}
Consequently the framework-wide implication has the typed form
\[
 \boxed{
 \text{generated law/generator}
 \longrightarrow
 \text{certified flow/process family}
 \longrightarrow
 \text{actual G0 history},}
\]
with no reversal of the arrows.
\end{theorem}

\begin{proof}
The first and fourth claims are the declared-output and member-occurrence
boundaries of
Theorem~\ref{p3a:thm:universal-generated-exact-continuation}.  The second is
Paper~II-A's theorem-status distinction applied to a process-generation
constructor: correct provenance of \(\mathscr L_D\) does not prove analytic
hypotheses which are not consequences of that provenance.  For the third,
a certified process family is still a family of admissible arrows until an
actual root token is supplied; conservative witness transport may co-realize
the resulting graph with that already occurring root, but cannot actualize
all possible roots or histories.  Thus law, certified flow, and actual history
remain distinct typed outputs.
\end{proof}


\begin{theorem}[Exact presentation transport of an existing action and law]
\label{p3a:thm:v204-action-law-transport}
Let $T:W\to V$ be the fixed isometry onto its image supplied by
Paper~II-C, Corollary~\ref{p2c:cor:v204-sampled-realization}, or a
restriction of that presentation to an existing invariant subspace.
For a generated finite law $H=H^*$ on $W$, put $\widetilde H=THT^*$.
Then, on $V_0=\operatorname{im}T$ (whose ambient identity is $TT^*$),
\[
 T e^{-itH/\hbar}T^*=e^{-it\widetilde H/\hbar},\qquad
 T[H,A]T^*=[\widetilde H,TAT^*].
\]
For real action variables use the underlying real spaces and their adjoints.
For a differentiable existing action $S(q)$ on a configuration space carried
by $W$ (with other variables unchanged), define its coordinate presentation
$\widetilde S(v)=S(T^*v)$ for $v\in V_0$. All retained variations satisfy
\[
 D^k\widetilde S(Tq)[Tu_1,\ldots,Tu_k]=D^kS(q)[u_1,\ldots,u_k].
\]
The Euler equations, stationary families, Hessian, its spectrum in the
transported kinetic metric, and the same-action Noether identities are
therefore equivalent. Write $B$ for the Hessian. If $W=W_\partial\oplus W_I$ and
$B_{II}$ is invertible, transport of both subspaces gives
\[
 \widetilde K_\partial=T_\partial
 (B_{\partial\partial}-B_{\partial I}B_{II}^{-1}B_{I\partial})T_\partial^*.
\]
\end{theorem}
\begin{proof}
The realized sampling of Paper~II-C supplies the isometry $T$.
Apply Paper~I-C, Theorems~\ref{p1b:thm:image-transport}
and~\ref{p1b:thm:action-transport}, to the already generated $H$ and
action $S$, rescaling the evolution parameter by $1/\hbar$. They
transport all existing derivatives, constraint kernels and the Schur
reduction on $V_0$. The action, sources and symmetries are those already
constructed here, so their transported invariance yields the same
Noether identities of Theorem~\ref{p3a:thm:v149-native-noether-frontier}.
\end{proof}
This theorem uses the existing action or generator. It creates no new choice
of Hamiltonian and supplies no missing interacting term. Where those objects
are already derived, representation transport adds no further physical
premise. Its hypotheses concern their stated finite domain, and do not
assume a continuum field theory has been restricted to this carrier.

\begin{corollary}[Moving presentations carry the induced connection]
\label{p3a:cor:v204-moving-presentation}
For a smooth family $T(s):W\to V_s$ of isometries onto their images, define
$\nabla_s^T v=T(s)\partial_s(T(s)^*v)$. Then
$\nabla_s^T(T\psi)=T\partial_s\psi$, and
$i\hbar\partial_s\psi=H(s)\psi$ is equivalent to
$i\hbar\nabla_s^T v=THT^*v$. The same pullback connection governs
source-dependent presentation changes.
\end{corollary}
\begin{proof}
Use the induced image-bundle connection in Paper~I-C,
Theorem~\ref{p1b:thm:action-transport}, for the present family $T(s)$.
Its intertwining identity applied to the already generated $H(s)$ gives
the displayed Schr\"odinger and source-variation laws.
\end{proof}
Ordinary derivatives of a moving frame include its derivative terms; the
fixed-presentation formula cannot discard them. This is coordinate transport
of an existing law, consistent with law--flow--history separation.

\begin{remark}[Generator--flow theorem families]
\label{p3a:rem:v149-generator-flow-families}
Theorem~\ref{p3a:thm:v149-law-flow-history-separation} is category-independent.
For a smooth finite-dimensional vector field, local ODE existence/uniqueness
is an analytic certification of its local flow.  For a constrained
hyperbolic field theory, a Cauchy-development theorem plays the analogous
role.  On a finite-dimensional norm-continuous linear semigroup, the
exponential representation is proved internally below.  Infinite-dimensional
semigroup generation requires its own domain/resolvent hypotheses.  None of
these mathematical generation theorems supplies the physical root token.
\end{remark}



\subsection{Action constructors and Noether constraints}
\subsubsection{Claim-complete action and generated evolution}

\begin{definition}[Claim-complete action owner]
\label{p3a:def:claim-complete-action-owner}
For a claim \(C\), fix its variation space \(V_C\), the maximal consumed
action-jet order \(r_C\), transformations \(G_C\), sources, marked
coefficients, boundary data and any reduction map.  An action owner is
complete for \(C\) if it includes every sector contributing to those
jets or boundary responses and is closed under the couplings and
transformation paths consumed by \(C\).  Equivalently, the collectively
omitted part has vanishing consumed jets and boundary contributions
on the stated domain and cannot re-enter through the retained
couplings or transformations.  An action-value or partition-function
claim also consumes the appropriate zero-order data.
\end{definition}

\begin{theorem}[Claim-complete action descent]
\label{p3a:thm:claim-complete-action-descent}
Fix such an owner and its differentiability domain.  Its Euler
equations, Hessian and vertices through order \(r_C\), and its
licensed Noether identities, are unchanged by omission of the
collectively invisible part specified in the definition.
For an actual complete-owner symmetry, the resulting identities
are constraints in the generated exact-continuation diagram.
A quotient evolution additionally requires the process
projectability condition.
\end{theorem}
\begin{proof}
Differentiation and boundary variation are linear in the action.
The omitted part has zero consumed responses, proving equality of
the claimed derivatives.  Differentiating invariance of the entire
retained action, including its transformed sources and marks,
gives the Noether identity on its admitted domain.  The native
Noether and exact-continuation theorems then apply to that owner.
For a reduced evolution \(q\circ U_t=\bar U_t\circ q\), the
right-hand side is well-defined precisely when the relevant process
maps \(q\)-fibres into \(q\)-fibres, as required by the projectability
theorem.  An Euler expression alone does not prove this condition.
\end{proof}

\begin{corollary}[Boundary-trivial same-action reclosure]
\label{p3a:cor:boundary-trivial-same-action-reclosure}
Let a claim-complete action owner have boundary configuration space
\(Q_{\partial}\) and convention-independent intrinsic boundary one-form
\(\Theta_{\partial}^{\rm intr}\).  Suppose a rooted reclosure satisfies both
of the following conditions:
\begin{enumerate}
\item the inherited boundary functional is locally constant on the reclosed
boundary role, so its first variation vanishes;
\item the frozen pre-action signature contains no independently generated new
boundary source, impulse, phase, or boundary functional with nonzero first
variation.
\end{enumerate}
Then the reclosed branch has
\begin{equation}
\label{p3a:eq:boundary-trivial-zero-cotangent}
 \boxed{\Theta_{\partial}^{\rm intr}=0.}
\end{equation}
In any regular coordinate chart \(q^A\) on \(Q_{\partial}\),
\(\Theta_{\partial}^{\rm intr}=\pi_A^{\rm intr}dq^A\), hence
\(\pi_A^{\rm intr}=0\) for every independent boundary coordinate.  The
statement is invariant under addition of a total derivative to the
Lagrangian: the raw canonical momenta and the explicit boundary functional
shift together, while \(\Theta_{\partial}^{\rm intr}\) is unchanged.
\end{corollary}

\begin{proof}
By condition (i), the inherited contribution to the boundary differential is
zero.  By Definition~\ref{p3a:def:claim-complete-action-owner}, every term with
nonzero consumed boundary variation belongs to the action owner.  Condition
(ii), together with the generated-action chronology, therefore forbids a new
nonzero boundary differential on the same-action reclosure branch.  Their sum
is Eq.~\eqref{p3a:eq:boundary-trivial-zero-cotangent}.  Independence of the
\(dq^A\) gives the component statement.  Under
\(L\mapsto L+dB/dt\), raw momenta shift by \(\partial_A B\) and the explicit
boundary functional shifts by the same exact one-form, so the intrinsic
combination is invariant.
\end{proof}

\begin{definition}[Frozen pre-action signature]
\label{p3a:def:frozen-preaction-signature}
For an action role $R$, the frozen pre-action signature
$\Sigma_R^{\rm pre}$ is the complete constituted information available before
a candidate action for $R$ is evaluated.  It may include generated field and
configuration carriers, operators, process laws, trace/readout grammar,
sources, boundary ports, orientations, gradings, and retained marks.  It does
not include a candidate-specific Euler equation, Hessian, vacuum, fitted
spectrum, stationary point, measured endpoint, or any other consequence that
first exists only after that candidate action is chosen.
\end{definition}

\begin{theorem}[Generated action-constructor closure]
\label{p3a:thm:generated-action-constructor-closure}
Let $\Sigma_R^{\rm pre}$ be frozen as in
Definition~\ref{p3a:def:frozen-preaction-signature}.  Suppose GP2 licenses a
result-independent typed constructor
\[
 \mathcal C_{\rm act}:\Sigma_R^{\rm pre}\longrightarrow
 \mathfrak S_R
 =(\mathcal C_R,S_R,\mathcal P_{\rm src},\mathcal P_{\partial},
   \mathcal T_R)
\]
whose output is a claim-complete action owner in the sense of
Definition~\ref{p3a:def:claim-complete-action-owner}.  Then every term,
coefficient, process profile, source port, boundary term, and normalization
that is part of $\mathcal C_{\rm act}(\Sigma_R^{\rm pre})$ is generated
descendant content and may not be independently deleted or rescaled while
claiming the same generated action owner.  Conversely, an additional term,
weight, profile, projector, or mark absent from the frozen signature requires
its own generated origin before it may enter that owner.

GP3 acts only on faithful realizations of the already generated action owner.
If those realizations possess a unique least structured class, that class is
retained.  If two inequivalent GP2-generated action constructors or
incomparable faithful minima survive, the framework retains the unresolved
action family and selects none.  In particular, this theorem is not a
``simplest Lagrangian'' principle and introduces no new constitutional GP.
\end{theorem}

\begin{proof}
Paper~II-A Corollary~\ref{p2a:cor:generated-constructor-chronology} makes the
output of a licensed frozen constructor descendant content.  GP2 exactness
therefore forbids deleting generated incidences and forbids inserting
candidate data with no frozen origin.  GP3 compares complete faithful
realizations only after that exactification and has no independent numerical
complexity functional.  The unresolved-family conclusion is the ordinary GP3
abstention rule.
\end{proof}

\begin{proposition}[Rooted singlet preservation under split-preserving reductions]
\label{p3a:prop:rooted-c3-split-preserving-closure}
Let $G=\mathbb C^3$, $c=(1,1,1)^{\mathsf T}/\sqrt3$,
$Q_0=cc^\dagger$, $Q_\perp=I-Q_0$, and let $X_G$ be the rooted
cyclic permutation.  Let $\mathscr A$ be the unital $*$-algebra
on a family-bearing carrier generated by family-blind operators,
$I_G$, and $X_G^{\pm1}$, all with their original declared sorts.
Every element of $\mathscr A$ commutes with $I\otimes Q_0$.
The same is true of inverses where defined, resolvents and analytic
functional calculus, and of Schur reductions whose retained and
complementary projectors commute with $I\otimes Q_0$.
In particular, every resulting family coefficient $M$ obeys
\[
 Q_\perp M Q_0=0=Q_0 M Q_\perp.
\]
This conclusion does not apply after inserting an independently generated
operator that breaks the rooted split or a projection that fails to
preserve it.
\end{proposition}
\begin{proof}
Since $X_Gc=c$, $[X_G,Q_0]=0$.  The commutant of
$I\otimes Q_0$ is a unital $*$-algebra, closed under the
listed operations. Inverses commute by multiplying the commuting
identity by the inverse on both sides. For Schur reduction, each
block and the inverse of the invertible complementary block
commutes with the induced rooted projection, so their Schur
combination does too.
\end{proof}


\begin{corollary}[Pre-action no-self-witness rule]
\label{p3a:cor:preaction-no-self-witness}
A candidate action cannot acquire GP2 origin from an Euler equation, Hessian,
vacuum, mass spectrum, fitted observable, or other consequence that is first
derived from that same candidate.  Such outputs may test a previously
generated action, but they cannot select it retroactively.
\end{corollary}

\begin{proof}
Those data are excluded from $\Sigma_R^{\rm pre}$ by definition.  Admitting
them as selectors would reverse the framework chronology and make the
candidate witness its own origin.
\end{proof}

\begin{corollary}[Bare/effective-action firewall]
\label{p3a:cor:bare-effective-action-firewall}
Uniqueness of a generated microscopic or bare action $S_R$ does not by itself
generate a quantum effective action $\Gamma_R$, a path/configuration measure,
a renormalized vacuum, or an occurring stationary member.  Those descend only
when the corresponding measure/quantum/refinement constructor is separately
licensed.  Thus
\[
 S_R\text{ generated}\quad\not\Rightarrow\quad
 \Gamma_R\text{ or its physical vacuum already generated}.
\]
\end{corollary}

\begin{corollary}[Born-equivariant constructor descent: Paper-II-A specialization]
\label{p3a:cor:born-equivariant-constructor-descent}
Let $B_\Sigma$ be the generated involutive Born automorphism of the complete
frozen input presentation of a licensed constructor
$\mathcal C:\Sigma\to Y$, including all marks, sources, boundary data,
calibrations and normalizations consumed by the claim.  If
$B_\Sigma\in\operatorname{Aut}_0(P;\mathcal C)$, then
Paper~II-A, Theorem~\ref{p2a:thm:generated-born-inheritance}, gives
\[
 \boxed{\mathcal C\circ B_\Sigma=B_Y\circ\mathcal C.}
\]
The target action $B_Y$ need not be trivial: it may fix a scalar, exchange
components, reverse a parity, conjugate an operator, or transport an action
graph in its declared target sort.  Only a trivial target action gives
ordinary invariance.  If Born changes a retained level or other mark, the
appropriate statement is instead the typed reciprocal-branch transport of
Paper~II-A, Proposition~\ref{p2a:prop:typed-born-branch-transport}.
\end{corollary}

\begin{proof}
This is exactly Paper~II-A,
Theorem~\ref{p2a:thm:generated-born-inheritance}, with the III-A constructor
as the licensed descendant word.  The marked-branch clause is the cited
transport proposition.  Thus III-A no longer owns an independent general
reciprocity principle.
\end{proof}

\begin{theorem}[Variational-sort firewall]
\label{p3a:thm:variational-sort-firewall}
Let a generated action/effective-action owner be written
\[
 \Gamma[\phi;\lambda],
\]
where $\phi$ belongs to the generated variational configuration sort while
$\lambda$ belongs only to a frozen external/source/mark sort.  Then the
generated stationarity equations include $d_\phi\Gamma=0$ but do not include
$\partial_\lambda\Gamma=0$.  The latter becomes a licensed variational
equation only after an independently generated upstream constructor has
promoted $\lambda$ into the declared configuration sort of the same action
owner.
\end{theorem}

\begin{proof}
The allowed variations are part of the claim-complete action owner and hence
of the frozen pre-action signature.  Varying an external mark would enlarge
that signature after the action has been constituted and would violate the
pre-action no-self-witness/no-insertion chronology.  Once a separate upstream
constructor generates the enlarged configuration role, ordinary action descent
supplies the additional Euler equation.
\end{proof}

\begin{theorem}[Unique solution-history descent]
\label{p3a:thm:unique-solution-history-descent}
Let a generated evolution law on a declared state space be written abstractly
as
\[
 \mathcal E(x)=0
\]
with an occurring initial/boundary datum $d$.  Assume that, before
substituting $d$, the law has a frozen, licensed exact-solution constructor
whose declared output is the solution graph/family in the same typed state
and time sorts, with its gauge/equivalence relation fixed.  Suppose that the
corresponding Cauchy/boundary-value problem is well posed on the claim
domain: a solution exists and is unique modulo that generated relation.
Then the solution history class $[x_d]$ is a generated descendant of the law,
the already licensed solution constructor, and the occurring datum.  No second history selector is licensed.  This
theorem does not generate the initial/boundary datum itself.
\end{theorem}

\begin{proof}
The assumed frozen exact-solution word produces the compatible solution
family from the constituted law and occurring datum; well-posedness is the
separate analytic certificate that this family is nonempty and has one
class modulo the generated equivalence.  Thus Paper~II-A
Theorem~\ref{p2a:thm:compatible-family-observable-closure} applies to its
already typed output.  Well-posedness alone would not have supplied a new
GP2-licensed output sort or a physical initial-state token.  The occurrence
firewall remains because the theorem conditions on $d$.
\end{proof}

\begin{corollary}[Transverse crossing descendant]
\label{p3a:cor:transverse-crossing-descendant}
Let $x_d(t)$ be the unique history of
Theorem~\ref{p3a:thm:unique-solution-history-descent} and let a generated
interface be the zero set of a $C^1$ scalar $h$.  If
\[
 h(x_d(t_*))=0,
 \qquad
 \frac{d}{dt}h(x_d(t_*))\neq0,
\]
then the crossing time is locally unique and is a generated descendant of the
history/interface data.  If $h(x_d(t))$ is strictly monotone on the entire
claim interval, the crossing is globally unique there.  Existence and
transversality are hypotheses, not consequences of the theorem.
\end{corollary}

\begin{proof}
Apply the implicit-function theorem to the scalar function
$t\mapsto h(x_d(t))$.  Strict monotonicity gives global uniqueness on the
stated interval.
\end{proof}

The hierarchy is therefore: geometry and configurations constitute
the action owner; its actual symmetries give Noether identities;
the generated process and any analytic certificate determine the
available evolution; an actual history is a further member in that
role.  None of these implications manufactures an absent tensor or
interaction.  A missing action is identified by the specified owner and domain
of each claim.  This does not assert one action containing every later physical sector.

\begin{lemma}[Regular constraints and infinitesimal ideal preservation]
\label{p3a:lem:regular-constraint-ideal}
Let $C=(C_1,\ldots,C_m):U\to\mathbb R^m$ be a smooth constraint map on a
finite-dimensional smooth configuration chart $U$, and let $X$ be a smooth
vector field there.
If $X[C_A]=M_A{}^BC_B$ with coefficients continuous along an integral curve,
then that curve preserves $C^{-1}(0)$. Conversely, if $C$ is a submersion
along $C^{-1}(0)$ and $X$ is tangent there, smooth local coefficients $M$
exist. Without regularity or a separate ideal-preservation certificate,
zero-set preservation alone does not imply this coefficient formula.
\end{lemma}
\begin{proof}
Along the curve, $C$ solves a homogeneous linear system, proving the first
claim. In regular coordinates $(u,z)$ with $z=C$, the function
$F_A=X[C_A]$ vanishes at $z=0$, so
\[
 F_A(u,z)=\sum_Bz_B\int_0^1\partial_{z_B}F_A(u,tz)\,dt.
\]
This gives the smooth local coefficients on the regular chart.
\end{proof}

\begin{theorem}[Native same-owner Noether frontier for exact continuation]
\label{p3a:thm:v149-native-noether-frontier}
\label{p3a:thm:v176-action-generated-noether-frontier}
Let \(D\) be a frozen parent with a licensed scalar action descendant \(S_D\)
and let \(s\mapsto g_s\) be a constituted continuous one-parameter family of
unmarked automorphisms of the \emph{complete action-owning presentation},
with $g_0=1$. Assume the first-variation/local integration-by-parts constructors required on
the declared domain are already licensed. At every configuration used,
require that $s\mapsto g_sD$ is differentiable at zero and that $S_D$ has
the corresponding curve chain rule; Fr\'echet differentiability on the
declared configuration chart is sufficient. Write
$R_\xi=\left.\frac{d}{ds}g_sD\right|_{s=0}$. A continuous finite-dimensional
linear group satisfies the orbit requirement by
Lemma~\ref{p2a:lem:repair-finite-generator-regularity}; a general continuous
family requires its own derivative certificate. Paper~II-A's same-owner
automorphism theorem then gives
\[
 S_D(g_sD)=S_D(D),
 \qquad
 dS_D(R_\xi)=0,
\]
so the corresponding rigid Noether-I consequences on that domain are an internal
theorem closure rather than an additional physical premise.  Noether-II
identities require the further constituted symmetry family with arbitrary
compactly supported local parameter functions, its differentiable finite-jet
action on the fields, and invariance for every such function. A rigid
one-parameter family alone supplies no such local test family. Let their typed
outputs be
\[
 \Delta_D^{N}
 =\left\{\begin{gathered}
 \text{Euler--Lagrange identities, currents, moment maps,}\\
 \text{constraints or Ward relations}
 \end{gathered}\right\}.
\]
Once these outputs satisfy Paper~II-A's licensed theorem-frontier conditions,
the exact-continuation signature is sharpened by
\[
 \boxed{\Delta_D\longmapsto
 \Delta_D^{+N}:=\Delta_D\cup\Delta_D^{N}.}
\]
Every admissible continuation must preserve the enlarged diagram.  In
particular, if the generated relations define a constraint locus
\(\mathcal C^{-1}(0)\subset X_D\), then
\[
 \boxed{
 \Psi_{s\to t}(\mathcal C^{-1}(0))
 \subseteq \mathcal C^{-1}(0)}
\]
for every admissible process arrow. A generated candidate smooth vector
field $X_{\rm dyn}$ must generate a flow preserving that zero locus.
On a finite-dimensional regular constraint chart with finitely many
constraints, where $\mathcal C$ is a submersion, tangency gives smooth
multiplication coefficients with
\[
 \boxed{X_{\rm dyn}[\mathcal C_A]=M_A{}^B\mathcal C_B}
\]
by Lemma~\ref{p3a:lem:regular-constraint-ideal}. At singular constraints this
formula requires a separate ideal-preservation certificate; zero-set
preservation alone is insufficient. An operator-valued propagation equation
requires its own typed operator domain, uniqueness theorem and boundary data.
An arbitrary configuration differential operator cannot replace this
propagation argument. If the
constraints generate gauge or diffeomorphism fibres, admissibility also
requires preservation of that vertical distribution, so that the process is
projectable to the corresponding reduced role whenever the quotient itself is
licensed.

If the action or symmetry is externally supplied, or if \(g_s\) changes a
generated source, coupling, boundary/interface datum, orientation, grading,
or any other retained mark of the complete action owner, the same-owner
invariance theorem does not apply and action invariance must be separately
proved before any Noether output may enter \(\Delta_D\).  Integrated global
charges and quantum Ward identities retain their separate boundary/global and
anomaly/measure hypotheses.
\end{theorem}

\begin{proof}
Paper~II-A's GP2 same-owner equivariance theorem makes a scalar action
descendant invariant under every unmarked automorphism of its complete frozen
input. On the stated orbit-derivative and chain-rule domain, differentiation
along \(g_s\) gives \(dS_D(R_\xi)=0\); the already
licensed local variational calculus then yields the corresponding rigid
Noether-I current relation without another physical principle. Explicitly,
if $\delta L=E_i\delta\phi^i+\partial_\mu\Theta^\mu$ and the rigid
symmetry has $\delta L=\partial_\mu B^\mu$, then
$\partial_\mu(\Theta^\mu-B^\mu)=-E_i\delta\phi^i$. If the additional
local symmetry family is constituted, write
\[
 \delta_\epsilon\phi^i=
 \sum_{|\alpha|\le r}R_A^{i,\alpha}\partial_\alpha\epsilon^A.
\]
Integrating by parts in $\delta_\epsilon S=0$ for arbitrary compactly
supported $\epsilon^A$ gives, by the fundamental lemma,
\[
 \sum_{|\alpha|\le r}(-\partial)_\alpha
       (E_iR_A^{i,\alpha})=0.
\]
These are the Noether-II identities on that stronger domain. Every varied
source and boundary term belongs to this calculation; an external source
variation is retained explicitly rather than set to zero by a field Euler
equation. The native local bosonic branch supplies the local symmetry
family in Paper~III-F, Theorem~\ref{p3f:thm:v176-local-bosonic-euler-noether}.
Paper~II-A's theorem-frontier rule registers those typed conclusions as
internal theorem closure.  Definition~\ref{p3a:def:generated-exact-continuation-signature}
already places prior generated equalities, owner relations, conservation laws,
and endpoint conditions in \(\Delta_D\), and
Theorem~\ref{p3a:thm:universal-generated-exact-continuation} requires every
candidate to satisfy that diagram. On a regular constraint chart the
coefficient-function form follows from Lemma~\ref{p3a:lem:regular-constraint-ideal};
the singular and operator branches retain the additional certificates just
stated. Preservation of the vertical symmetry distribution is the local
projectability requirement. Global smooth descent additionally uses the
licensed quotient and the full fibre condition; connected fibres or an
explicit component comparison discharge the distinction as in
Corollary~\ref{p3a:cor:v166-smooth-process-descent}.  The final scope
clauses are the same-owner and analytic-certification firewalls: changing a
retained antecedent mark changes the physical input, while global charges and
quantum identities require mathematical hypotheses not contained in the
local scalar invariance relation.
\end{proof}

\begin{corollary}[Noether-constrained one-ray reduction]
\label{p3a:cor:v176-noether-constrained-one-ray}
Let \(K_D^\tau\) be the exact admissible generator set of
Theorem~\ref{p3a:thm:finite-generator-one-ray-rigidity} and set
\[
 K_{D,N}^\tau=\{L\in K_D^\tau:
 e^{tL}\text{ preserves every relation in }\Delta_D^{+N}\}.
\]
Noether data can reduce this set without selecting a state or history.
If its nonzero part has one positive real projective direction and the
required nonstationarity witness excludes zero, its laws are exactly
\(e^{\kappa tL_*}\) for the surviving rates \(\kappa\in I_{D,N}\).
The law is unique under the one-ray argument precisely when the induced
structured-isomorphism/calibration quotient of \(I_{D,N}\) is a singleton.
Several inequivalent rates or directions are not selected by Noether alone.
A homogeneous candidate cone may be used for those preliminary conditions
whose scaling closure is proved; endpoint and strict conditions remain in
the exact set.
\end{corollary}
\begin{proof}
Intersect the complete admissibility conditions with the generated Noether
relations and apply the positive-ray and rate-quotient clauses of the cited
theorem. Relations already in \(\Delta_D\) are not additional data.
\end{proof}

\begin{remark}[Constraint dynamics does not select physical time]
\label{p3a:rem:v176-constraint-history-clock-firewall}
A first-class generally covariant bulk Hamiltonian may be a combination of
constraints.  Preservation of that constraint system therefore does not by
itself identify a physical clock, deparametrize the reduced dynamics, choose
an initial state, or select an actual history.  Such data remain separately
typed by their own continuation and occurrence hypotheses.
\end{remark}


\begin{corollary}[Identity survival and oriented nonstationarity]
\label{p3a:cor:identity-survival-oriented-nonstationarity}
Suppose a stationary continuation \(I_D\) exists in the declared process type.
If
\[
 I_D\in\operatorname{Cont}^{\rm ex}(\Sigma_{\rm cont}(D)),
\]
then GP2 and GP3 do not force a nontrivial process.  In particular, the
existence of a nonzero intrinsic defect or generated distinction is
insufficient when \(I_D\) preserves it.

A nontriviality conclusion therefore requires a generated
\emph{oriented nonstationarity witness}: a typed comparison
\(\mathcal N_\alpha\), fixed before candidate evaluation, for which every
admissible physical continuation must satisfy
\[
 \mathcal N_\alpha(\Phi)\succ0,
 \qquad
 \mathcal N_\alpha(I_D)=0,
\]
or an equivalent strict decay, propagation, enlargement, crossing, matching,
or endpoint condition.  Such a witness removes the stationary candidate; it
does not select among several inequivalent nonzero continuations.
\end{corollary}

\begin{proof}
If \(I_D\) preserves every frozen relation, type, intrinsic defect signature,
and constructor closure, it has zero GP2 realization defect.  Paper~II-A
Corollary~\ref{p2a:cor:zero-member-occurrence-test} then forbids GP2 from
manufacturing a nonzero member, while GP3 may compare only the surviving
complete candidates and has no independent numerical tie-breaker.  A strict
generated comparison failed by \(I_D\) excludes it from the joint zero-defect
locus.  Excluding one member plainly does not prove that the remaining locus
has a unique least class.
\end{proof}

\begin{corollary}[Common invariants and attractors without microscopic selection]
\label{p3a:cor:common-process-invariant}
Let the surviving continuation moduli be nonempty, and let
\[
 R:\mathfrak M_{\rm cont}(D)\longrightarrow\mathsf Y
\]
be a licensed generated readout into a predeclared target sort.  If
\[
 R([\Phi])=y_*
 \qquad
 \text{for every surviving }[\Phi],
\]
then \(y_*\) is common generated content even when no process member is
selected.  The statement applies to a stationary role, fixed-point subspace,
endpoint projector, or attractor only when its defining fixed-point or limit
operation belongs to the licensed readout schema and exists on every
surviving member.  Actual convergence along a particular history remains a
token-level statement.
\end{corollary}

\begin{proof}
Apply part~(iii) of
Theorem~\ref{p2a:thm:generated-exact-solution-moduli}.  The image singleton is
common family content, while the graph of \(R\) projects isomorphically to the
unresolved source family and therefore does not select it.  The final
qualification is the declared-output and role/token boundary.
\end{proof}

\subsection{Finite generators, projectability, and threshold descent}
\begin{theorem}[Finite-dimensional generator and positive-ray classification]
\label{p3a:thm:finite-generator-one-ray-rigidity}
Let a comparison specification have one fixed finite-dimensional real or
complex carrier \(V\), succession \((\mathbb R_{\ge0},+)\), and
norm-continuous homogeneous linear candidates with \(T_0=I\). Every
candidate has a unique generator \(L\) and \(T_t=e^{tL}\). Its exact
admissible generator set is
\[
 \boxed{K_D^\tau:=\{L\in\operatorname{End}_{\mathcal G_D}(V):
   (e^{tL})_{t\ge0}\in
   \operatorname{Cont}^{\rm ex}(\Sigma_{\rm spec}(D))\}.}
\]
This membership includes the full process type, GP1 domain, joint GP2
conditions and frozen diagram, including endpoints and strict comparisons.
The map \(L\mapsto(e^{tL})\) is a bijection onto that exact family.
Physical generatedness additionally uses its certified word as in
Theorem~\ref{p3a:thm:universal-generated-exact-continuation}.

Define positive projectivization by \(L\sim_+L'\) when \(L'=aL\)
for \(a>0\). If
\[
 K_D^\tau\setminus\{0\}\ne\varnothing,\qquad
 \mathbb P_+(K_D^\tau\setminus\{0\})=\{[L_*]_+\},
\]
put \(I_D=\{\kappa>0:\kappa L_*\in K_D^\tau\}\). All nonstationary
laws are exactly
\[
 \boxed{T_t=e^{\kappa tL_*},\qquad\kappa\in I_D.}
\]
This is one positive direction, not automatically one rate. Uniqueness by
this argument holds exactly when \(I_D/\!\sim_D\) is a singleton, where
\(\sim_D\) is induced by the actual structured isomorphisms and
\(\operatorname{Cal}_D\) on the complete presentation. A strict
nonstationarity test may exclude zero without selecting a sign; the sign
must follow from its actual oriented inequality or the displayed positive-ray
hypothesis. Ordinary real or complex projectivization is not used here.

For compact symmetry, write the actual representation as
\(V\simeq\bigoplus_\lambda S_\lambda^{\oplus m_\lambda}\). Then
\[
 \operatorname{End}_{G_D}(V)\simeq
 \prod_\lambda M_{m_\lambda}(\mathbb D_\lambda),\qquad
 \mathbb D_\lambda=\operatorname{End}_{G_D}(S_\lambda).
\]
Here \(\mathbb D_\lambda=\mathbb C\) on a complex irreducible, and
\(\mathbb R,\mathbb C\) or \(\mathbb H\) on a real irreducible, with
the corresponding multiplicity-module convention. The complete comparisons
cut a set in this commutant. It may be called a cone only after proving the
relevant scaling closure, and a convex cone only after proving convexity.
An empty exact set gives no admissible law; \(K_D^\tau=\{0\}\) gives
stationarity unless a mandatory strict condition makes the full set empty.
Several surviving directions or rates are not resolved by this argument.
\end{theorem}
\begin{proof}
For small \(h>0\), the matrix
\[
 B_h=\frac1h\int_0^hT_s\,ds
\]
is close to \(I\) and is invertible. Exact composition gives
\[
 T_tB_h=\frac1h\int_t^{t+h}T_r\,dr.
\]
The right side is differentiable, including a right derivative at zero.
Thus
\[
 T'_t=T_tL,\qquad L=\frac{T_h-I}{h}B_h^{-1}.
\]
The finite matrix ODE gives \(T_t=e^{tL}\); the derivative at zero proves
uniqueness. Differentiating naturality gives commutation with the symmetry,
and the exponential series proves its converse. The full membership
condition therefore gives the stated bijection, without replacing global
conditions by their derivatives at zero. The positive-projective hypothesis
gives exactly the displayed set of multiples; quotient uniqueness is exactly
the singleton rate-quotient condition.

Every continuous additive orientation-preserving time change is
\(t\mapsto at\), \(a>0\): prove the formula on nonnegative rationals
using additivity and extend by continuity. With a permitted presentation
isomorphism \(S\),
\[
 S e^{atL}S^{-1}=e^{t(aSLS^{-1})}.
\]
Such a map identifies rates only when it is licensed and preserves or
correctly transports every retained clock and endpoint datum. Mathematical
availability is not a calibration license. Schur's lemma gives the displayed
commutant decomposition. No scaling property of the complete constraint set
follows from that decomposition or from time homogeneity alone.
\end{proof}

\begin{proposition}[Exact common-clock affinity certificate]
\label{p3a:prop:v2211-clock-affinity-certificate}
Let $K(s),N(s)$ be two $C^2$ readouts of one actual history on a connected
interval, with $N'(s)>0$. Their relation is affine, $K=cN+b$, if and only if
\[
 \boxed{K''N'-K'N''=0.}
\]
The initial values and one slope then determine $b,c$. If instead both
readouts have already been proved continuous additive characters of the
\emph{same} oriented succession monoid, their common zero and the additive
Cauchy equation give $b=0$ and constant $c$ directly. A one-dimensional
homogeneous carrier, monotonicity or scale homothety does not prove either
hypothesis. In particular the native trace-depth readout gives the exact
counterexample $K=56(1-e^{-s})$, $N=s$, whose affinity residual is
$-56e^{-s}\ne0$.
\end{proposition}
\begin{proof}
The derivative of $K'/N'$ is
$(K''N'-K'N'')/(N')^2$. Its vanishing is equivalent to a constant ratio;
one integration gives the affine relation and the converse is immediate.
For common additive characters, continuity gives linear functions of their
shared additive parameter, as in the generator theorem. The displayed native
visibility readout verifies the counterexample by direct differentiation.
Thus the test must be evaluated on the actual two clock readouts before
exporting a constant rate to a physical e-fold or other history coordinate.
\end{proof}

\begin{corollary}[Finite normal-form certificate transport]
\label{p3a:cor:v166-finite-normal-form-certificate}
If \(K_D^\tau\), the required exponential-family image and its quotient
are already licensed retained outputs, their composite word certifies the
complete finite homogeneous continuation family. The exponential bijection
above supplies the exact-image proof for
Lemma~\ref{p3a:lem:v166-continuation-certificate-transport}.
Rewriting an existing semigroup in this normal form introduces no new process
occurrence. Conversely, an ambient set of matrices does not acquire physical
generatedness merely from this representation theorem.
\end{corollary}


\begin{corollary}[Finite unitary generator without a separate Stone premise]
\label{p3a:cor:v149-finite-unitary-generator}
Let \(\mathcal H\) be finite dimensional and let
\(U:\mathbb R\to U(\mathcal H)\) be an already constituted norm-continuous
one-parameter unitary group.  Applying
Lemma~\ref{p2a:lem:repair-finite-generator-regularity} gives a unique
\(L\in\operatorname{End}(\mathcal H)\)
with
\[
 U(t)=e^{tL}.
\]
Differentiating \(U(t)^*U(t)=I\) at \(t=0\) yields
\(L^*=-L\).  Hence
\[
 \boxed{G:=iL=G^*,\qquad U(t)=e^{-itG}.}
\]
Thus on the finite branch the self-adjoint generator representation is an
internal consequence of the already occurring unitary group law and does not
supply that group, its physical Hamiltonian interpretation, or an actual
history.
\end{corollary}

\begin{remark}[Finite versus continuum semigroup generation]
\label{p3a:rem:v149-finite-continuum-semigroup-status}
No Hille--Yosida-type analytic input is needed for the finite-dimensional
norm-continuous exponential statement above.  On an infinite-dimensional
carrier, density of the generator domain, closedness, dissipativity or
resolvent bounds, and the relevant continuity assumptions are genuine
analytic-certification hypotheses and are not consequences of G0 or
GP1--GP3 merely from correct provenance.
\end{remark}

\begin{remark}[Discrete, nonautonomous, nonlinear, and ordered branches]
\label{p3a:rem:continuation-generator-boundary}
Exact composition is universally functorial but does not universally imply
one autonomous exponential.  On \(\mathbb N\), one obtains
\(T_n=T_1^n\) and no canonical logarithm.  A differentiable two-parameter
cocycle \(U(t,s)\) gives
\(\partial_tU(t,s)=L(t)U(t,s)\); autonomy additionally requires generated
translation homogeneity.  A differentiable nonlinear semiflow gives a vector
field \(\dot x=F(x)\), whose Hamiltonian or gradient form requires the
corresponding generated symplectic or metric/action data.  A directed support
poset may give only a filtration by inclusions.  Regularity is therefore part
of the signature or the strongest invariant regularity of an independently
forced exact object; it is not selected after the desired equation is known.
\end{remark}

\begin{theorem}[Projectability of a generated parent process]
\label{p3a:thm:generated-process-projectability}
Let \(\Phi_\alpha^X:X_s\to X_t\) be an exact continuation and let
\(\pi_s:X_s\to Y_s\) be generated readouts surjective on underlying
sets. A unique descended set continuation satisfying
\[
 \boxed{\pi_t\Phi_\alpha^X=\Phi_\alpha^Y\pi_s}
\]
exists exactly when
\[
 \pi_s(x)=\pi_s(x')\quad\Longrightarrow\quad
 \pi_t\Phi_\alpha^X(x)=\pi_t\Phi_\alpha^X(x')
\]
for every arrow. It inherits identities and exact composition.

For descent in a concrete typed category, additionally require the relevant
factorization certificate: each morphism \(\pi_t\Phi_\alpha^X\)
constant on the fibres must factor through \(\pi_s\) by a morphism of
that category. A quotient with that universal property suffices;
underlying surjectivity alone does not. Smooth surjective submersions on
smooth manifolds without boundary give the smooth instance in
Corollary~\ref{p3a:cor:v166-smooth-process-descent}. Topological quotient
maps give continuous individual arrows; joint parameter continuity requires
the corresponding parameterized quotient condition. Other arrow properties
retain their actual factorization domains.

On fixed finite linear carriers, descent through a surjective linear
\(\pi\) is equivalent to invariance of \(\ker\pi\). For the
homogeneous generators it is equivalent to
\[
 \pi L_X=L_Y\pi,
 \qquad L_Y=\pi L_Xs,
\]
where \(s\) is any linear right inverse and the result is independent of
that proof-device choice. Failure of fibre constancy prevents a closed
memoryless process on the displayed set role; failure only of structured
regularity must not be described as failure of the set factorization.
\end{theorem}
\begin{proof}
Necessity follows from the intertwining square. Under fibre constancy define
\(\Phi_\alpha^Y(\pi_sx)=\pi_t\Phi_\alpha^X(x)\). It is independent
of the lift, and surjectivity gives uniqueness. Precompose each identity and
composition equation with its surjection to reduce it to the parent equation.
The additional factorization property supplies exactly the missing morphism
claim. In the linear case, \(L_Y(\pi x)=\pi L_Xx\) is well-defined
exactly when \(L_X\ker\pi\subseteq\ker\pi\). Two right inverses
differ into that kernel. Induction gives \(\pi L_X^n=L_Y^n\pi\), so
the exponential series intertwines; differentiation gives the converse.
\end{proof}

\begin{corollary}[Smooth process and infinitesimal descent]
\label{p3a:cor:v166-smooth-process-descent}
Let the parent and target be smooth manifolds without boundary and each
\(\pi_s\) a smooth surjective submersion. Fibre constancy in the preceding
theorem is necessary and sufficient for smooth descended arrows. A jointly
smooth parameterized parent descends jointly smoothly.

For a fixed such map \(\pi:X\to Y\) and a smooth vector field \(V_X\),
a unique smooth field \(V_Y\) is related by
\[
 d\pi_x(V_X(x))=V_Y(\pi(x))
\]
exactly when the left side agrees at every point of each fibre. This relation
is equivalent to the flow intertwining on common flow domains. A complete
parent field yields a complete descended field. In submersion charts,
vertical-bracket preservation is the local infinitesimal test; to infer the
global fibre condition from it, require connected fibres or compare their
components explicitly. Boundaries and corners use an appropriate local-section
certificate rather than an automatic extension of this domain.
\end{corollary}
\begin{proof}
A submersion has local smooth sections \(\sigma:O\to X_s\). The set
descendant restricts to \(\pi_t\Phi_\alpha^X\sigma\), so it is smooth;
the formula with the same section proves joint parameter smoothness. Fibre
constancy makes section choices agree. Applying this argument to
\(d\pi\circ V_X\) gives the smooth vector field. The chain rule and
uniqueness for smooth ODEs prove flow intertwining; differentiation proves
its converse. A lift of an initial point and its complete parent trajectory
give a global projected solution. In coordinates \((y,z)\), write
\(V_X=a(y,z)\partial_y+b(y,z)\partial_z\). Vertical-bracket preservation
is \(\partial_z a=0\), proving constancy locally along fibres; connectedness
or an inter-component check supplies global constancy.
\end{proof}

\begin{theorem}[Generated threshold-event descent and clock-loss obstruction]
\label{p3a:thm:threshold-event-descent}
Let \(I=[t_0,t_1]\), let \(x:I\to J\subset\mathbb R\) be a generated
continuous nondecreasing parent history in the inherited orientation, and let
\[
 b_{k_0+1}<b_{k_0+2}<\cdots
\]
be a generated locally finite threshold family.  On the range where the
maximum is finite, define
\[
 q(z):=\max\bigl(\{k_0\}\cup\{k\ge k_0+1:z\ge b_k\}\bigr),
 \qquad q(t):=q(x(t)).
\]
Put \(m_0=q(t_0)\).  For every \(k>m_0\) reached by the history, define its
first-entry time
\[
 \tau_k:=\inf\{t\in I:x(t)\ge b_k\}.
\]
Then:
\begin{enumerate}[label=(\roman*),leftmargin=2em]
\item every \(\tau_k\) is attained,
      \(x(\tau_k)=b_k\), and consecutive reached thresholds satisfy
      \(\tau_k<\tau_{k+1}\);
\item the threshold-event readout is the exact discrete successor process
      \[
       q(\tau_k)=k,
       \qquad q(\tau_k^-)=k-1,
       \qquad S_n(k)=k+n
      \]
      wherever the indicated later events exist.  If generated carriers
      \(Y_k\) have their unique retained inclusions
      \(\iota_{\ell\leftarrow k}:Y_k\hookrightarrow Y_\ell\), then
      \[
       \iota_{r\leftarrow\ell}\iota_{\ell\leftarrow k}
       =\iota_{r\leftarrow k}
      \]
      is the corresponding exact event-indexed continuation;
\item if \(x\) is strictly increasing through one reached threshold on an
      open neighborhood of its entry time, there is no
      translation-homogeneous deterministic physical-time family
      \(T_h\) on the integer role satisfying
      \[
       q(t+h)=T_h(q(t))
      \]
      for all admissible \(t,h\ge0\).  In particular, the integer readout is
      not a closed homogeneous semigroup state for that parent history;
\item on every nonterminal threshold cell, the generated residual coordinate
      \[
       r(z):=b_{q(z)+1}-z
      \]
      makes \(R(z):=(q(z),r(z))\) injective, since
      \(z=b_{q(z)+1}-r(z)\).  Hence an exact parent continuation already closed on the displayed
      \(z\)-role transports by set conjugacy to the image of \(R\).
      If \(z\) is itself a readout of a larger state, its projectability
      must first be proved. The quotient \(q\) alone transports exactly
      only when it passes
      Theorem~\ref{p3a:thm:generated-process-projectability}.
      With the ordinary discrete topology on the integer component,
      \(R\) jumps at thresholds; no continuous or smooth product-coordinate
      conjugacy is asserted. Cellwise smooth dynamics uses the corresponding
      event reset, or a separately licensed transported target structure.
\end{enumerate}
Thus a continuous parent generates a unique ordering of its threshold events
and an exact discrete successor law, but neither an event rate nor a
physical-time law on the coarse integer alone.  The coordinate \(r\) is a
canonical generated fibre-separating completion; no uniqueness among all
possible coordinates and no parent dynamics is asserted.
\end{theorem}

\begin{proof}
For a reached \(b_k\), the set in the definition of \(\tau_k\) is nonempty
and closed in the compact interval \(I\), so the infimum is attained.
Continuity excludes \(x(\tau_k)>b_k\), while monotonicity and the first-entry
definition exclude \(x(\tau_k)<b_k\).  Since \(b_{k+1}>b_k\), continuity
also gives \(\tau_k<\tau_{k+1}\).  Immediately before \(\tau_k\), after the
preceding entry, the threshold readout is \(k-1\); at \(\tau_k\) it is \(k\).
This proves the successor formula.  The inclusion identity is transitivity of
the already generated retained inclusions, hence gives exact composition.

For the third claim, choose a sufficiently small \(\delta>0\) about a strict
crossing \(\tau_k\) so that no other threshold is crossed there.  Put
\[
 s_1=\tau_k-2\delta,
 \qquad s_2=\tau_k-\delta,
 \qquad h=\frac{3\delta}{2}.
\]
Then \(q(s_1)=q(s_2)=k-1\), but
\(q(s_1+h)=k-1\) and \(q(s_2+h)=k\).  A single map \(T_h\) cannot send the
same input to both outputs.  This is precisely failure of fibre preservation
in Theorem~\ref{p3a:thm:generated-process-projectability}.  Finally, the
displayed reconstruction proves injectivity of \(R\); defining
\(R\Phi R^{-1}\) on its image transports identities and composition exactly
as set maps. At a crossing the limiting residual zero is replaced by the next
cell width, while the integer increases. This proves the stated topological
qualification. Paper~III-G, Appendix~\ref{p3a:app:v166-descent-diagnostics} records
explicit smooth-quotient and threshold examples.
\end{proof}

\subsection{K16 source localization and projected process descendants}
\label{p3a:sec:v179-k16-source-localization}

Fix the generated oriented terminal direction $u$ and let $a$ lie in the
primitive scalar source line.  Define the self-adjoint balanced operator
\[
  X_{a,u}=|u\rangle\langle a|+|a\rangle\langle u|.
\]
Using the retained pseudoscalar $\omega$ and the already generated Clifford
first-order evaluation, the balanced $K=16$ source role has the local
cotangent representative
\[
 \boxed{\rho_{16}^{\rm src}(2X_{a,u})=(a\omega)_A e^A\in T^*U_{\mathbb C}.}
\]
This is a source/readout localization.  It does not add a new bare one-form
field to the minimal configuration tangent: a nongauge metric perturbation,
a gauge variation with no internal generator, or a scalar/spinor variation
requires its own generated action tangent.

The generated-process projectability theorem of Paper~II-A applies to any
later scalar readout of a process family generated from this same rooted role.
In particular, downstream contractions may be transported to a scalar defect
coordinate when fibre constancy is proved; the resulting scalar process is a
descendant, not a new law coefficient.


\begin{remark}[K16 response grades do not constitute a primitive fermion source]
\label{p3a:rem:v182-k16-fermion-sort-firewall}
An identity among hierarchy/response weights, including
$E_{16}^2/E_0=E_{32}$, is an action-response statement in its declared
cotangent/vector/readout sort.  It does not by itself create an odd finite
fermion $\mathsf{SrcVal}$.  A same-action constructor with that target sort is
required before K16 data can enter the primitive fermion source chain.
\end{remark}



\subsection{Normalized successor-cell residual}
\label{p3a:sec:v184-successor-cell-residual}

\begin{theorem}[Normalized successor-cell residual]
\label{p3a:thm:v184-normalized-successor-residual}
Let a generated continuous history traverse two consecutive protected
thresholds $b_j<b_{j+1}$ and let the already-generated residual coordinate on
the open cell be
\[
 r_j(K)=b_{j+1}-K,
 \qquad b_j\le K<b_{j+1}.
\]
Then the dimensionless normalized residual
\[
 \boxed{\rho_j(K)=\frac{b_{j+1}-K}{b_{j+1}-b_j}}
\]
is generated without an additional clock or scale.  For
$b_j\le K_a<K_b<K_c<b_{j+1}$ define
\[
 \rho_j(K_b\mid K_a)
 :=\frac{r_j(K_b)}{r_j(K_a)}.
\]
Then
\[
 \boxed{\rho_j(K_c\mid K_a)
 =\rho_j(K_c\mid K_b)\rho_j(K_b\mid K_a)}.
\]
Consequently
\[
 \boxed{\ell_j(K_b\mid K_a):=-\log\rho_j(K_b\mid K_a)}
\]
is additive under concatenation.
\end{theorem}

\begin{corollary}[Protected width-eight cells]
\label{p3a:cor:v184-width-eight-cell}
For the protected hierarchy $b_j=8j$, write
\[
 s_j=\frac{K-8j}{8}\in[0,1),
 \qquad
 \rho_j=1-s_j.
\]
No independent exponential rate is licensed: any exponential representation
$\rho=e^{-\ell}$ is only the logarithmic reparameterization of the generated
threshold residual.
\end{corollary}

\subsection{Exact continuation of the coordinate-generated native algebra}
\label{p3a:sec:v226-coordinate-continuation}
\begin{corollary}[Polynomial transport retains the complete native operator sort]
\label{p3a:cor:v226-coordinate-continuation}
Let $U$ be an actual unitary continuation on the intrinsic momentum fibre.
Transporting the four coordinates transports their whole algebra, the six
round generators and the separate duality/chirality operators:
\[
 \mathcal A_D\longmapsto U\mathcal A_DU^*,\qquad
 p(Y)\longmapsto Up(Y)U^*.
\]
All native matrices admit coordinate-word representatives of degree at most
six. This does not turn an isometric nonreducing compression into a unitary
continuation, nor make its omitted-image products vanish.
\end{corollary}
\begin{proof}
Unitary conjugation preserves every product and adjoint. Apply
Theorems~\ref{p1:thm:v226-coordinate-only}--\ref{p1:thm:v226-chirality}.
The fixed and covariantly transported coefficient jets are
Corollary~\ref{p1b:cor:v226-source-jets}; a compression retains the separate
word-defect identities of Paper II-B.
\end{proof}
The shared-variable chronology of the complete action still precedes its
scalar readout. A transported basis provides exact tests of an actual
native response, rather than a new selection of its action, source or history.



\section{Co-realization, common action, and source chronology}
\label{paper3a:sec:organized-action}
\subsection{Claim-complete action co-consumption}
\label{p3a:sec:v183-action-coconsumption}

\begin{corollary}[Claim-complete action co-consumption]
\label{p3a:cor:v183-action-coconsumption}
Let the generated action constructor of
Theorem~\ref{p3a:thm:generated-action-constructor-closure} have a
claim-complete output on its declared differentiability domain.  Every
nonredundant generated contribution required by the consumed action jet is
co-consumed with its inherited coefficient, incidence, source/current port,
kinetic metric and normalization.  No second binding or occurrence premise is
licensed for those same generated constituents.

In particular, for a complete generated resolution $\sum_aP_a=I$ of an actual
finite source operator $\mathcal D$,
\[
 \mathcal D=\sum_{a,b}P_a\mathcal D P_b,
\]
a claim consuming $\mathcal D$ retains every nonzero block
$P_a\mathcal D P_b$, including the mixed blocks.  For a nonlinear generated
restriction $x=f(y)$ the consumed second jet is
\[
 D^2(S\circ f)=D^2S(df,df)+dS(D^2f),
\]
and the second contribution may be omitted only after its vanishing has been
proved on the stated domain.
\end{corollary}

\begin{proof}
Theorem~\ref{p3a:thm:generated-action-constructor-closure} retains every term,
coefficient, source port, boundary term and normalization in the generated
owner.  Paper~II-A Theorem~\ref{p2a:thm:v3101-generated-coconsumption} makes
that retention joint for the complete constructor output.  The block identity
is insertion of a complete resolution, and the displayed Hessian is the
ordinary chain rule for the already licensed restriction.  Neither operation
creates a new physical constructor.
\end{proof}

\begin{remark}[Action co-consumption firewall]
\label{p3a:rem:v183-action-coconsumption-firewall}
Co-consumption closes binding only after a typed action constructor exists.  It
does not turn a scalar norm, diagnostic trace, common host, or separately
available operator into an action term.  A missing particular action word
remains a constructor-provenance problem rather than a co-consumption problem.
\end{remark}
\begin{corollary}[Reuse of the complete parent before restriction]
\label{p3a:cor:complete-parent-descendant-reuse}
Let a licensed action/source constructor have an occurring, nonobstructed,
claim-complete output \(Q\) on its declared source-variation domain.
Let \(F(Q)\) be a licensed descendant into a predeclared physical target,
using the retained action/source jets, state or boundary input actually
required by that claim.  Then \(F(Q)\) is co-realized from the same parent:
its first later consumer adds no independent occurrence, binding coefficient
or normalization for that already generated output.
If a licensed restriction \(\pi(Q)\) discards data needed by \(F\), construct
\(F\) from the retained \(Q\) before applying that restriction.  Reconstructing
it from \(\pi(Q)\) additionally requires an independently licensed
\(\widetilde F\) and a proved factorization
\(F=\widetilde F\circ\pi\) on the claimed domain.
For a generated nonempty complete parent family, a licensed whole-family
\(F\) retains its compatible image; common values close without a selector.
\end{corollary}
\begin{proof}
Compose the certified parent word with the licensed \(F\).
Paper~II-A rooted co-realization and
Theorem~\ref{p2a:thm:v3101-generated-coconsumption} retain its inputs,
incidences and output.  The existing claim-complete action descent supplies
exactly the consumed jets; no omitted term is declared absent merely because
a restricted presentation drops it.  A readout through \(\pi\) equals \(F\)
precisely when its displayed factorization holds.  Whole-family reuse and
common-value closure follow from
Corollary~\ref{p2a:cor:v3101-coconsumed-family-readout}.
\end{proof}
This result preserves an existing complete source word.  Quantum measure,
state, symmetry/Ward antecedents, a new interaction incidence, global history
and consumed limits enter that word only when separately generated or proved.
Critical directions are retained in the full relation of Paper~II-A,
Theorem~\ref{p2a:thm:v3110-finite-critical-stationary-relation}, before any
noncritical elimination.



\subsection{Co-realized actions and their common physical variations}
\label{p3a:sec:v2216-common-action}

GP2 co-realization retains each action's owner and incidence. Several
co-realized functionals need not be equal. Their contribution to one physical
field equation requires actual field/source maps into a common configuration,
unit, density and boundary chart. If those maps and the interaction word have
been generated, the physical functional has the declared form
\[
 S_{\rm phys}=\sum_i\iota_i^*S_i+S_{\rm int}.
\]
The sum is an operation supplied by that word, not a consequence of merely
putting its constituents in a common host. Each Einstein, vacuum and source
contribution is allocated once. In particular a bare parent term and its
effective descendant are not both added unless the actual matching map has
explicitly subtracted their overlap. A whole-action change of units transports
all relative source and interaction coefficients together.

\begin{theorem}[Common-action consumption and shared-interior reduction]
\label{p3a:thm:v2216-common-action-consumption}
Suppose the declared physical pullbacks identify one retained real variable
$u$ and one \emph{shared} interior variable $z$ in finite real Hilbert spaces.
At a common background let a finite family of complete quadratic contributions,
after their unit and density maps, be
\[
 S_i^{(2)}(u,z)=\tfrac12\langle u,A_i u\rangle+
 \langle u,B_i z\rangle+\tfrac12\langle z,C_i z\rangle,
 \qquad C_i=C_i^*>0.
\]
Set
\[
 K_i=A_i-B_iC_i^{-1}B_i^*,\quad L_i=C_i^{-1}B_i^*,\quad
 C_\Sigma=\sum_iC_i,\quad
 \overline L=C_\Sigma^{-1}\sum_iC_iL_i.
\]
The stationary reduction of the physical sum has retained Hessian
\begin{equation}
 \boxed{K_\Sigma=\sum_iK_i+
 \sum_i(L_i-\overline L)^*C_i(L_i-\overline L).}
 \label{p3a:eq:v2216-shared-variance}
\end{equation}
Thus $K_\Sigma\succeq\sum_iK_i$, with equality exactly when all the
$L_i$ agree. If the field maps instead supply independent interiors $z_i$,
the reduced Hessian is exactly $\sum_iK_i$.

If $K_i(C)\succeq0$ at a terminal curvature port, total terminal nullity is
equivalent to $K_i(C)=0$ for every $i$ and equality of all $L_i(C)$. For
$C^1$ coefficient families with uniformly positive $C_i$, this compatibility
also makes the last term of \eqref{p3a:eq:v2216-shared-variance}
$O((R-C)^2)$. Hence the linear terminal slope is
$K_\Sigma'(C)=\sum_iK_i'(C)$.
\end{theorem}
\begin{proof}
Complete every square:
\[
 S_i^{(2)}=\tfrac12\langle u,K_i u\rangle+
       \tfrac12\|C_i^{1/2}(z+L_i u)\|^2.
\]
The sum has its unique stationary interior at
$z_*=-\overline L u$. Substitution gives
\eqref{p3a:eq:v2216-shared-variance}. Its extra term is a sum of positive
quadratic forms; it vanishes for every $u$ precisely when
$(L_i-\overline L)u=0$ for every $i,u$. This requires no commutation of the
$C_i$. Independent interiors instead let every completed square vanish
separately. Under the additional $K_i(C)\succeq0$ hypothesis, total nullity
forces every nonnegative summand to vanish, and the converse is immediate.
The uniform inverse bound on the $C_i$ makes $L_i$ and $\overline L$ $C^1$;
their endpoint differences are $O(R-C)$. Squaring proves the order statement
and its derivative consequence.
\end{proof}

\begin{remark}[Sign, gap and physical-incidence scope]
\label{p3a:rem:v2216-multiaction-scope}
The identity concerns positive stationary interior pivots. It is not a
positivity theorem for a Lorentzian constraint block or a Gaussian with a
negative or zero pivot. Closed operator forms require a common form domain
and the existence of the displayed inverse products. A gap closing at the
endpoint lies outside the uniform-gap conclusion and requires its actual
source and boundary scaling. For example
$S_\pm=c(R-C)(z\mp au)^2/2$ has separate zero reductions but shared
$K_\Sigma=2ca^2(R-C)$ for $R>C$; its endpoint inverse does not exist.
This proves a mathematical mechanism, not the generation of that pair in the
framework. Distinct action owners can therefore contribute consistently,
but co-realization alone supplies neither the physical sum, a required
coefficient, nor terminal compatibility. Paper~III-D's complete-jet test,
Theorem~\ref{p3d:thm:v2216-complete-action-jets}, applies after the common
physical incidence has been supplied.
\end{remark}


\begin{proposition}[Complete two-jet equality on one finite action chart]
\label{p3a:prop:v245-common-action-two-jet-equality}
Let $\mathcal U$ be a connected open subset of a finite real
configuration/source coefficient chart, and let $S_\star,S_{\rm comp}$
be $C^2$ real-valued functionals \emph{already constituted on the same
chart}. If $D^2S_\star(q)=D^2S_{\rm comp}(q)$ for every $q\in\mathcal U$
and their values and first differentials agree at one $q_0\in\mathcal U$,
then $S_\star=S_{\rm comp}$ throughout $\mathcal U$.
\end{proposition}
\begin{proof}
The difference has vanishing Hessian, so its first differential is
locally constant, hence constant by connectedness. Its differential
and value vanish at $q_0$ and therefore everywhere.
\end{proof}
\begin{remark}[What a complete action two-jet must contain]
\label{p3a:rem:v245-two-jet-owner-firewall}
This is a uniqueness \emph{criterion}, not an origin theorem for
$S_\star$ or $S_{\rm comp}$. For a graded fermionic action the full
coefficient word, all mixed chiral/source/connection derivatives, and
its Grassmann ordering must be retained; ordinary bosonic two-jets
alone are not a substitute. A fixed source, nontrivial topological
sector, different boundary condition, or distinct nonlocal action
owner cannot be erased by comparing one local chart.
The exact finite-heat identity is
Proposition~\ref{p3b:prop:v245-exact-finite-heat-two-jet}.
\end{remark}

\subsection{Constructive local parent reconstruction and its exact obstruction}
\label{p3a:sec:v246-constructive-parent}

\begin{theorem}[Constructive finite-chart parent from a complete Hessian field]
\label{p3a:thm:v246-integrable-parent-reconstruction}
Let $\mathcal U\subset\mathbb R^n$ be open and star-shaped about $q_0$.
Let $H:\mathcal U\to\operatorname{Sym}^2(\mathbb R^n)^*$ be $C^1$,
with every configuration, source and boundary coefficient consumed by the
claim included among the $q^i$.  There is a $C^3$ function $S$ with
$D^2S=H$, $S(q_0)=s_0$ and $dS(q_0)=\ell_0$ if and only if
\begin{equation}
 \partial_kH_{ij}=\partial_iH_{kj}\qquad\hbox{for all }i,j,k.
 \label{p3a:eq:v246-hessian-integrability}
\end{equation}
When it exists it is unique and has the constructive formula
\begin{equation}
 \boxed{S(q_0+x)=s_0+\ell_0(x)
       +\int_0^1(1-t)H(q_0+tx)[x,x]\,dt.}
 \label{p3a:eq:v246-radial-parent}
\end{equation}
This is an existence and uniqueness theorem on the stated affine chart;
its constitutional use additionally requires independent provenance of
$H,s_0,\ell_0$ and their common coefficient incidence.
\end{theorem}
\begin{proof}
Necessity follows by commuting the third partial derivatives of $S$.
For sufficiency put
$F_i(q_0+x)=(\ell_0)_i+\int_0^1H_{ij}(q_0+tx)x^j\,dt$.
Symmetry and \eqref{p3a:eq:v246-hessian-integrability} make all three
indices of $\partial_kH_{ij}$ symmetric. Consequently
\[
 \partial_kF_i=\int_0^1\left(H_{ik}
       +t x^j\partial_jH_{ik}\right)(q_0+tx)dt
   =[tH_{ik}(q_0+tx)]_{t=0}^{t=1}=H_{ik}(q_0+x).
\]
Thus $F$ is a closed one-form. Its radial primitive is
$s_0+\int_0^1F(q_0+tx)[x]dt$; exchanging the two integrals gives
\eqref{p3a:eq:v246-radial-parent}.  Its differential is $F$, hence its
Hessian is $H$. The stipulated regularity gives $S\in C^3$.
Uniqueness is Proposition~\ref{p3a:prop:v245-common-action-two-jet-equality}.
\end{proof}

\begin{corollary}[Mixed incidence is necessary input to parent reconstruction]
\label{p3a:cor:v246-sector-restriction-kernel}
Separate sector restrictions, even together with all their finite-order
ambient jets, do not determine a common analytic parent. More precisely,
let two nontrivial scalar coefficient directions be $x,y$ and fix a finite
jet order $r$. For any even integer $m>r$ the analytic family
\[
 S_\lambda(x,y)=S_0(x,y)+\lambda x^m y^m
\]
has the same ambient derivatives through order $r$ as $S_0$ at every
point of both axes $x=0$ and $y=0$, while distinct $\lambda$ give
distinct parents on any open neighborhood of the origin. The perturbation
is also even in each direction. Therefore no reconstruction based only
on such separated data can fix its mixed coefficient.
\end{corollary}
\begin{proof}
Every partial derivative of total order at most $r$ leaves a positive
power of both $x$ and $y$ in $x^my^m$ and therefore vanishes on either
axis. At points with $xy\ne0$ the perturbation is nonzero when
$\lambda\ne0$. This is an analytic information-loss example, not an
assertion that any nonzero $\lambda$ is a licensed physical action word.
\end{proof}

\begin{remark}[Graded coefficient scope and pre-action chronology]
\label{p3a:rem:v246-parent-reconstruction-scope}
For finitely many odd generators, first expand the even action in a fixed
ordered Grassmann basis. Apply the theorem separately to every retained
ordinary coefficient, with its own anchored value and first differential;
or use the complete graded derivatives with their ordering signs. A bosonic
coefficient Hessian alone cannot recover an unreported fermion bilinear or
higher Grassmann coefficient. On non-star-shaped charts a closed derivative
field additionally needs vanishing periods for a global primitive; overlapping
local primitives must agree in their affine transition data. The theorem
supplies neither a quantum measure nor a global boundary condition.
Most importantly, deriving $H$ from a proposed $S$ and then reconstructing
that same $S$ cannot provide its independent origin, by
Corollary~\ref{p3a:cor:preaction-no-self-witness}. The constitutive input is
the independently generated, complete common source word, not merely a list
of successful sector outputs.
\end{remark}


\paragraph{Rigor frontier of the one-root action.}
\label{p3a:par:v245-common-action-rigor-frontier}
\emph{Algebraic level:} the born-scalarized heat two-jet, Gram/polar
identities, source-visible finite Feshbach criterion, and source-conditioned
stationary Schur identities are exact on their displayed finite domains.
\emph{Geometric level:} the rooted negative-Hopf triple, harmonic
$163\oplus5\oplus1$ boundary owner, and rooted nonlinear scalar moments are
generated mathematical providers, not automatically physical action slots.
\emph{Action level:} the occurring local joint Dirac operator and certified
finite-fermion handoff supply some real jets, whereas the charged hierarchy
origin test, physical current relative incidence, neutral mass-word
incidence, and boundary-source/normal-port identification remain separate.
\emph{Physical-response level:} mass-only two-pole criteria and exact finite
resolvent algorithms are proved; the complete interacting tagged roots and
residues must still be evaluated on the actual state/constraint/source word.
\emph{Global/continuum level:} the local full-heat expansion does not by
itself identify the restored $f_{2\mid7}$, the unchanged reference cap, the
full nonlocal tensor/boundary action or a continuum QFT measure/limit.
No level is promoted to another by co-realization or numerical agreement.

\subsubsection{Complete cap consumption as a conditional constructor}
\label{p3a:sec:v223-cap-consumption}
\begin{corollary}[Boundary and source presentations consume one action]
\label{p3a:cor:v223-cap-consumption}
If the pre-action signature supplies the extension, conjugate-flux
incidence, curvature and density maps, and complete incoming source of
Theorem~\ref{p3e:thm:v223-exact-cap-reduction}, then that theorem is a
generated-constructor calculation on the supplied signature.  The
positive source comparison of
Theorem~\ref{p1b:thm:v223-positive-source-transport} preserves the full
source word, its mixed jets and finite induced measure; it changes no
eliminated scalar action.  Boundary and Schur presentations of its
gravitational coefficient count the same contribution once.
\end{corollary}
\begin{proof}
Apply the source congruence to the complete word before elimination and
use its pointwise identity and Jacobian.  Apply generated action
constructor closure (Theorem~\ref{p3a:thm:generated-action-constructor-closure})
to the specified maps.  Equality of two representations of one word
does not produce a second independent summand.
\end{proof}
This corollary does not manufacture its pre-action signature from the
desired answer.  In particular an available positive metric isometry
does not identify the native incoming section or its incidence.  The
shared-interior test of Proposition~\ref{p3e:prop:v223-port-incidence}
must be performed on actual common variables, rather than eliminating
two copies of a shared variable independently.  The cap is a conditional
candidate for a full native evaluation, not a new fundamental principle
and not yet a theorem replacing all other action branches.


\subsection{Reference-source transport and shared-variable chronology}
\label{p3a:sec:v224-source-chronology}
The cap consumption result of Corollary~\ref{p3a:cor:v223-cap-consumption}
now also admits the reference-boundary presentation of
Theorem~\ref{p3f:thm:v224-reference-cap-parent}. Its translated source is
the required descendant $j_r=(\mathcal A_r-\mathcal A_0)Q_0$ of
the supplied normal operator and fixed datum. The complete reference
energy, source centre, normalization, endpoint and positive orthogonal
complement must be transported before taking a terminal scalar value.

The form-space map of Proposition~\ref{p3f:prop:v224-collar-finite-transport}
proves an off-shell parent equivalence under its hypotheses, rather than
only equality of two minima. Its four-port kernel consumes the geometric
$u_{\rm op}(r)$. The equality of that coefficient to a stationary
normalized collar flux does not authorize replacing it by a separately
varied off-shell flux: the latter changes the joint Euler equation.
The shared-interior incidence requirement of
Proposition~\ref{p3e:prop:v223-port-incidence} remains in force.

The primitive-owner retraction and retained-Weyl coefficient comparison
of Proposition~\ref{p1b:prop:v224-retained-primitive} likewise act on
the complete contracted word. Their mathematical availability does not
replace the ordinary native action by that readout or erase companions.
I/II co-realization closes all required descendants after the actual
constructor and inputs are licensed; no new physicality premise is
needed for those outputs. The native cap comparison still needs that
constructor's normal band, boundary datum, complete source and energy
allocation. Equal counts or final scalar functions do not supply a
missing action-owner map.



\section{Forward common-action assembly and response-jet coherence}
\label{p3a:v251:sec:coherence}

This construction begins with action contributions and field/source maps
that are already licensed on a common finite chart.  It does not infer a
physical interaction from co-realization alone.  Boundary coefficients,
source marks, density factors and the ordered odd-coordinate presentation
are included before taking any derivative.  A single chart is used below;
global gluing additionally requires the actual transition maps and their
period/cocycle conditions.

\subsection{The complete forward pullback, coefficient by coefficient}
Let $Q$ be an open finite real configuration/source/boundary chart.
For each actually consumed contribution let $\iota_i:Q\to Q_i$ be
$C^3$, let $S_i$ be $C^3$ on its declared chart, and let $S_{\rm int}$
be the generated interaction term, with every overlap allocated once.
Then the actual declared assembly is
\[
 S(q)=\sum_i S_i(\iota_i(q))+S_{\rm int}(q).
\]
At every admitted point its full Hessian is
\begin{equation}
\label{p3a:v251:eq:full-pullback-hessian}
 H(q)[v,w]=\sum_i\left\{
 D^2S_i[D\iota_i v,D\iota_i w]
       +DS_i[D^2\iota_i(v,w)]\right\}
       +D^2S_{\rm int}[v,w].
\end{equation}
The contact $DS_i[D^2\iota_i]$ survives even at a stationary point of
the complete sum unless its separate vanishing is proved: a sum of
Euler covectors can vanish without its summands vanishing.

For a finite ordered Grassmann presentation write
$S=\sum_\alpha s_\alpha(q)\vartheta^\alpha$.  Apply the displayed
formula to every retained ordinary coefficient $s_\alpha$, after the
actual odd-field transition has been pulled back and re-expanded in that
fixed ordering.  Thus the calculation retains fermionic bilinears and
higher odd coefficients rather than attempting to recover them from
the zero-fermion bosonic Hessian.  Existing source-frame or boundary
maps are differentiated in the same pullback; they are not reset between
the sector calculations.

On a star-shaped affine subchart, the independently computed coefficient
Hessian fields $H_\alpha$ form one action precisely when they have the
actual value/gradient anchors and
\[
 \partial_c(H_\alpha)_{ab}=\partial_a(H_\alpha)_{cb}.
\]
The forward assembly above satisfies these identities by commuting
ordinary third derivatives.  Conversely the existing native radial
primitive reconstructs each coefficient uniquely:
\[
 s_\alpha(q_0+x)=s_\alpha(q_0)+ds_\alpha(q_0)[x]
    +\int_0^1(1-t)H_\alpha(q_0+tx)[x,x]\,dt.
\]
This is an exact consistency test for a supplied complete common word.
It neither generates a missing incidence nor licenses a target action
by reconstructing the action from its own derivatives.

For two overlapping charts with actual transition $g$, the same
formula reads
\[
 H_\alpha^{\rm old}=Dg^*H_\alpha^{\rm new}Dg
                  +ds_\alpha^{\rm new}\,D^2g,
\]
with any odd-coordinate mixing handled before coefficient extraction.
The second derivative of a composition is
\[
 D^2(f\circ g)[v,w]
 =D^2f[Dg v,Dg w]+Df[D^2g(v,w)].
\]
Consequently complete two-jet pullbacks associate exactly on triple
overlaps whenever the underlying transition maps do.  Dropping the
second term generally destroys the \emph{Hessian transformation law}
for the actual scalar two-jet: the truncated tensor-congruence rule
\(H\mapsto Dg^*HDg\) still associates algebraically, but in general
fails to be the Hessian of the pulled-back scalar when its first
derivative is nonzero.

\subsection{Second background jets and shared-interior compatibility}
Now let $q\in Q$ denote background/control coefficients, while $u$ and
$z$ are retained and shared interior perturbations.  After all complete
pullbacks above, suppose the actual quadratic constituents are
\[
 S_i^{(2)}(q;u,z)=\tfrac12 u^*A_i(q)u
       +\operatorname{Re}\,u^*B_i(q)z+\tfrac12z^*C_i(q)z,
 \qquad C_i(q)=C_i(q)^*>0.
\]
Assume the matrix coefficients are $C^2$ on a neighborhood of $q_0$;
the interior dimensions, domains and source pairings are fixed there.
The same result holds for real matrices with transpose instead of dagger.
Set
\[
 K_i=A_i-B_iC_i^{-1}B_i^*,\quad L_i=C_i^{-1}B_i^*,\quad
 C_\Sigma=\sum_iC_i,
\]
\[
 \bar L=C_\Sigma^{-1}\sum_iC_iL_i,\qquad
 \Delta_i=L_i-\bar L.
\]
The native complete-square identity gives the exact shared-interior
reduction
\[
 K_\Sigma=\sum_iK_i+V,\qquad
 V=\sum_i\Delta_i^*C_i\Delta_i.
\]

\begin{theorem}[Second background jets of the common Schur response]
\label{p3a:v251:thm:shared-second-jet}
Suppose the actual sector stationary slopes agree at $q_0$:
$L_i(q_0)=L_0$ for every $i$.  Then $V(q_0)=0$ and $DV(q_0)=0$.
For background directions $v,w$ define
\[
 M_v=C_\Sigma(q_0)^{-1}\sum_i C_i(q_0)D_vL_i(q_0),
 \qquad T_{i,v}=D_vL_i(q_0)-M_v.
\]
The exact next response jet is
\begin{equation}
\label{p3a:v251:eq:shared-second-jet}
 \boxed{D^2V(q_0)[v,w]
 =\sum_i\left(T_{i,v}^*C_i(q_0)T_{i,w}
                  +T_{i,w}^*C_i(q_0)T_{i,v}\right).}
\end{equation}
Thus the common reduced Hessian and the sum of separately reduced
Hessians have identical background jets through order two at $q_0$
if and only if
\[
 D_vL_i(q_0)=D_vL_j(q_0)
 \quad\hbox{for all }i,j,v.
\]
For a coordinate basis $e_a$ of the finite background chart, this is
equivalent to the single nonnegative scalar test
\begin{equation}
\label{p3a:v251:eq:finite-variance-certificate}
 \boxed{\sum_{a,i}\|C_i(q_0)^{1/2}T_{i,e_a}\|_{\rm HS}^2=0.}
\end{equation}
\end{theorem}
\begin{proof}
Differentiate $C_\Sigma\bar L=\sum_iC_iL_i$.  At compatibility the
terms containing $DC_i$ cancel, leaving
$D_v\bar L=M_v$ and $D_v\Delta_i=T_{i,v}$.  In the two derivatives
of $V=\sum\Delta_i^*C_i\Delta_i$, all terms except those with one
derivative on each $\Delta_i$ retain an undifferentiated $\Delta_i$
and vanish at $q_0$.  This proves the displayed second-jet formula.

For $v=w$ it is $2\sum_iT_{i,v}^*C_iT_{i,v}\succeq0$.
Strict positivity of the $C_i$ shows that its vanishing is equivalent
to every $T_{i,v}$ vanishing, which says that all $D_vL_i$ equal
their weighted common value.  Checking a coordinate basis is sufficient
by linearity.  The Hilbert--Schmidt sum is zero exactly when each of
these finitely many matrices vanishes.  Polarization then gives all
mixed second jets.
\end{proof}

\paragraph{Independent first-order variational verification.}
Complete each sector quadratic at fixed retained vector $x$ as
$Q_i(x,y;q)=x^*K_i(q)x+(y+L_i(q)x)^*C_i(q)(y+L_i(q)x)$.
Minimization over one \emph{shared} interior variable gives
$y_*(q)=-\bar L(q)x$, whereas independent interiors remove each
positive square separately.  At the compatible point $q_0$ all slopes
are equal.  Along $q(t)=q_0+tv$, the shared minimizer's \emph{first}
variation is $-M_vx$ and the $i$th residual is
$t(D_vL_i-M_v)x+O(t^2)=tT_{i,v}x+O(t^2)$.
Thus its excess over separately minimized sectors is
\[
 x^*V(q(t))x=t^2\sum_i x^*T_{i,v}^*C_i(q_0)T_{i,v}x+o(t^2).
\]
Polarization yields exactly the displayed mixed second jet, and
positivity yields the stated zero test.  Variations of $C_i$ do not
contribute at this order because each zeroth-order residual vanishes.
This linearized stationary-equation proof avoids twice differentiating
the variance itself but retains every positive-pivot, regularity and
compatibility hypothesis; it is not an additional physical occurrence
or mass-source argument.

The theorem concerns background jets of the retained quadratic Hessian
$K(q)$.  When $q$ and the retained perturbation are different variables,
these can be third and fourth variations of the original action.  It
does not assert that fourth derivatives are necessary to compare an
ordinary action two-jet at a single fixed point.

\begin{corollary}[Frame covariance of the second-jet defect]
\label{paper3a:corollary:atlas-frame-covariance-of-the-second-jet-defect-220dfda}

If the common retained and interior frames change through smooth
unitaries $U(q)$ and $W(q)$, then
$\Delta_i\mapsto W^*\Delta_iU$.  At the compatible point
$D_v\Delta_i\mapsto W(q_0)^*D_v\Delta_iU(q_0)$.
The criterion above is therefore independent of these frame choices.
\end{corollary}
\begin{proof}
The actual quadratic pullback gives
$C_i\mapsto W^*C_iW$, $B_i\mapsto U^*B_iW$, and hence
$L_i\mapsto W^*L_iU$, with the same law for $\bar L$.
Derivatives of $U,W$ multiply $\Delta_i(q_0)=0$ and vanish at that
point.  The squared norms and the transformed Hessian are consequently
unchanged up to the retained unitary conjugation.
\end{proof}

\begin{proposition}[First nonzero higher-order response mismatch]
\label{paper3a:proposition:atlas-first-nonzero-higher-order-response-mismatch-4fe1b0f}

Along a one-parameter regular background $q(t)$ suppose
$\Delta_i^{(j)}(0)=0$ for $0\le j<r$ and the coefficients are
sufficiently smooth, with $C_i(0)>0$.  Put $A_i=\Delta_i^{(r)}(0)$.
Then
\[
 V(t)=\frac{t^{2r}}{(r!)^2}\sum_iA_i^*C_i(0)A_i+o(t^{2r}),
\]
\[
 V^{(2r)}(0)=\binom{2r}{r}\sum_iA_i^*C_i(0)A_i.
\]
\end{proposition}
\begin{proof}
Taylor expand $\Delta_i$ to its first potentially nonzero order and
insert in the exact positive variance.  The first coefficient is
nonnegative and cannot cancel between sectors.
\end{proof}

These results require regular positive pivots.  A closing gap, divergent
Legendre pivot, indefinite Lorentzian constraint block or changing
operator domain needs its own weighted/operator analysis.  In particular,
the theorem is not silently extended to the critical terminal scalaron
endpoint simply by taking a finite-matrix formula there.

\subsection{Exact example: individually vanishing reductions do not glue}
For scalar $u,z$ and background $q$ take
\[
 S_\pm^{(2)}(q;u,z)=\tfrac12(z\pm qu)^2.
\]
Each separately reduced $K_\pm$ is identically zero.  Their common
interior has $\bar L=0$, $\Delta_\pm=\pm q$, so
\[
 K_\Sigma(q)=2q^2,\qquad
 K_\Sigma(0)=K_\Sigma'(0)=0,\qquad K_\Sigma''(0)=4.
\]
Thus matching the endpoint response and its first slope does not
establish the next background response jet.  This is an exact diagnostic
of the proposed integration, not a claim that these two particular
actions have native physical occurrence.

\subsection{Normalization is a branched comparison, not a universal patch}
The native common-heat transport is
$U_J^{\rm raw}=96U$ and $G_{J,\rm raw}=G_{\rm scal}/96$.
III-E explicitly admits the branch in which the actual physical action
is the raw joint action, so $G_{\rm phys}=G_{J,\rm raw}$.  The condition
$z=1/96$ belongs to the stronger claim $G_{\rm phys}=G_{\rm scal}$
under a certified uniform whole-action transport.

Therefore a forward integration ledger should retain two conditional
branches: (a) actual raw joint coefficient and its coupling dictionary;
(b) a separately generated complete normalization map if the equality
of the Newton coordinates is claimed.  This does not remove that stronger
claim by redefining $G$, and it does not identify a scalar readout with
a fermionic first-order pairing.  The corresponding Paper~III-E proposition states this qualification explicitly.

\paragraph{Native anchors.}
The starting shared-interior identity is
Theorem~\ref{p3a:thm:v2216-common-action-consumption}, and the coefficient
reconstruction is Theorem~\ref{p3a:thm:v246-integrable-parent-reconstruction}.
The full physical matching criterion is Paper~III-D,
Theorem~\ref{p3d:thm:v2216-complete-action-jets}, with its following remark.
The branched normalization comparison is Paper~III-E,
Proposition~\ref{p3e:prop:v2211-whole-action-newton-certificate}.
The second-jet theorem extends the existing
positive shared-interior identity, not a completed physical parent.

\section{Action information, observable comparison, and metrology}
\label{paper3a:sec:organized-readouts}
\subsection{Claim-relative action information and native observable determination}
\label{p3a:sec:v188-claim-relative-observable}

\begin{definition}[Native claim data and compatible readout family]
\label{p3a:def:v188-native-claim-data}
Fix a claim $C$, its occurring frozen parent $P$, and the consumed domain
of Definition~\ref{p3a:def:claim-complete-action-owner}.  Let
$\mathcal A_C(P)$ be a nonempty complete surviving compatible family
obtained by the already licensed constructors and their result-independent
generated congruence.  Every action owner in this family is complete for
$C$, with the sources, boundary data, state or support, normalization,
reduction and analytic hypotheses actually consumed by $C$ retained in
their declared sorts.  Write
\[
 f_C:\mathcal A_C(P)\longrightarrow D_C:=f_C(\mathcal A_C(P))
\]
for a licensed result-independent record of the data available before the
requested readout.  Let
\[
 O_C:\mathcal A_C(P)\longrightarrow Y_C
\]
be the requested invariant readout of the same licensed native word.
The family, congruence, data record, target sort and readout word are frozen
before any member-specific result is evaluated.  A data fibre used as a
physical compatible family is formed by its licensed exact-data restriction;
a mathematically available comparison is not added to that family by hosting
alone.  A token or analytic domain condition enters this record only through
its existing occurrence or theorem-domain certificate.
\end{definition}

\begin{theorem}[Native observable determination on compatible data fibres]
\label{p3a:thm:v188-native-observable-fibre}
On the domain of Definition~\ref{p3a:def:v188-native-claim-data}, a unique
set readout $\bar O_C:D_C\to Y_C$ satisfying
\[
 \boxed{O_C=\bar O_C\circ f_C}
\]
exists if and only if
\[
 \boxed{f_C(a)=f_C(b)\ \Longrightarrow\ O_C(a)=O_C(b)
 \qquad(a,b\in\mathcal A_C(P)).}
\]
For each already licensed compatible data fibre, its common generated
readout value is native descendant content even when that fibre contains
several inequivalent action or Hamiltonian realizations.  A singleton
complete realization family is sufficient for determination, whereas
uniqueness of the entire action or Hamiltonian is not required for this
particular output.

A claim about $\bar O_C$ as a structured or computable map, or as a map in
an additional retained output schema, uses the factorization and computation
certificates actually required by its declared category and licensed schema.
Set fibre constancy alone supplies none of those certificates and licenses
no new constructor, output sort, occurrence or selected member.
\end{theorem}
\begin{proof}
If $O_C=\bar O_Cf_C$, equal recorded data give equal readout values.
Conversely define $\bar O_C(d)=O_C(a)$ for any $a$ with $f_C(a)=d$.
Fibre constancy makes this independent of the lift, and the definition of
$D_C$ gives existence and uniqueness.  On a licensed compatible fibre,
Paper~II-A Corollary~\ref{p2a:thm:compatible-family-observable-closure}
retains a generated readout with singleton image without selecting a member.
The retained action family may consequently remain non-singleton.
The structured factorization requirement is the one already separated from
set descent in Theorem~\ref{p3a:thm:generated-process-projectability}; a
computation certificate must likewise produce the claimed readout from its
recorded input in the admitted output schema.  The set construction is not
an origin certificate for such a schema.
\end{proof}

\begin{corollary}[Claim-relative same-action completeness]
\label{p3a:cor:v188-claim-relative-completeness}
For $C$, retain its complete consumed action jets, boundary responses,
source and state data, and every coupling or transformation path through
which those data can contribute.  Action contributions collectively
invisible to all of these consumed data may be omitted under
Definition~\ref{p3a:def:claim-complete-action-owner}.  The resulting
licensed readout $O_C$ is unchanged, so such omissions do not require
uniqueness of unconsumed action sectors.  Action-value and partition-function
claims retain their consumed zero-order data.

When the complete parent $Q$ is retained, a licensed $O_C(Q)$ is formed
before a noninjective compression $\pi(Q)$ that discards its required data.
Using only $\pi(Q)$ instead consumes an independently licensed factorization
$O_C=\widetilde O_C\circ\pi$ on the actual domain.  Failure of this proxy
factorization establishes information loss in that proxy.  Native
nondetermination from the recorded data requires two members
$a,b\in\mathcal A_C(P)$ with
\[
 f_C(a)=f_C(b),\qquad O_C(a)\ne O_C(b),
\]
with both members and the same-input comparison admitted by the actual
constructor grammar.
\end{corollary}
\begin{proof}
Claim-complete action descent
(Theorem~\ref{p3a:thm:claim-complete-action-descent}) preserves precisely
the consumed jets and boundary contributions.  The no-re-entry clause of
the owner definition preserves them along the retained coupling and
transformation paths.  Applying the same licensed readout word to those
unchanged inputs therefore gives the same output.  Complete-parent reuse
(Corollary~\ref{p3a:cor:complete-parent-descendant-reuse}) supplies the
uncompressed route and states the additional factorization needed for the
compressed route.  The last assertion is the negation of fibre constancy
in Theorem~\ref{p3a:thm:v188-native-observable-fibre} on its actual compatible
family.  A failed comparison outside that family does not supply its witness.
\end{proof}

\begin{definition}[Claim-relative completion certificate]
\label{p3a:def:v188-native-completion-certificate}
For a requested native output, record the following finite certificate:
\begin{enumerate}
\item the exact target $C$, retained data $f_C$, compatible family and claimed
branch/domain, including each actually consumed state, boundary and limit;
\item the licensed parent-to-output word and its required input certificates,
with consumed action information retained before any lossy readout;
\item the executed output or proved common-fibre value, together with every
regularity, computation or physical-interpretation certificate consumed by
that target;
\item if that result is not established, the earliest unconstructed or unexecuted native map
or unproved consumed predicate; a claim of nondetermination instead records
an admissible same-input, different-output witness as above.
\end{enumerate}
This certificate distinguishes a missing execution or proof from a proved
non-singleton observable image.  Additional action sectors, continuum or
flat limits, and Standard--Model dictionaries belong to the certificate
precisely when the declared target consumes them.  It is a record of the
existing G0/GP1--GP3 calculus and introduces no constitutive principle.
\end{definition}




\subsection{Projective cone protection and constants readout typing}
\label{p3a:sec:v185-cone-rg-metrology}

\begin{corollary}[Common kinetic renormalization does not renormalize the vacuum front cone]
\label{p3a:cor:v185-projective-cone-protection}
Let a later vacuum field be an undeformed descendant of the same principal
owner as the operational quadratic form
\[
 q_{\rm op}(\omega,p)=\omega^2-c_{\rm op}^2|p|_h^2.
\]
If its exact resolution-boundary flow changes the principal Hessian only by a
positive scalar factor,
\[
 \sigma_2(P_{Q,k})=Z_Q(k)\,q_{\rm op}\,A_Q
\]
with $A_Q$ invertible at principal order, then its characteristic set is
independent of $k$ and
\[
 \boxed{c_{{\rm front},Q}=c_{\rm op}.}
\]
More generally, for independent temporal/spatial normalizations
$Z_t,Z_s$,
\[
 \frac{\partial_{\log k}c_Q}{c_Q}
 =\frac12(\eta_t-\eta_s).
\]
Thus a nonzero running of the front speed requires an actual relative
principal deformation; it is not implied by the running of one common kinetic
normalization.
\end{corollary}

\begin{remark}[Operational speed versus constants-paper dimensionless readout]
\label{p3a:rem:v185-constants-speed-firewall}
The physical metric rate $c_{\rm op}=2\pi R_x/T_t$ and a dimensionless
constants-paper number $\widehat c$ are differently typed until an explicit
owner/metrology bridge is supplied.  The equality of a vacuum photon front
speed with $c_{\rm op}$ is a principal-symbol statement.  It neither derives
an SI numeral from scale-free geometry nor licenses an adjustment of the
primitive spatial metric to repair a separate constants-paper residual.
\end{remark}



\subsection{Local descent of the symmetric quadratic branch}
\label{p3a:sec:v189-local-symmetric-quadratic-descent}
This forward interface is conditional on the rank-four relative solder
supplied by Paper~III-B and its stated first-symbol/coframe descent. It is
not needed for the native reciprocal, phase or continuation chain. On that
generated domain, the relative solder reaches the local causal real form
through the reciprocal first-symbol map.  Let $e^A$
be the resulting local coframe and
\[
 g=\eta_{AB}e^A\otimes e^B
\]
the corresponding Lorentzian metric on that local domain.  For an
infinitesimal coframe variation
\[
 \delta e^A=M^A{}_B e^B,
 \qquad
 M_{AB}=\lambda\eta_{AB}+S_{AB}+\Omega_{AB},
\]
with $S_{AB}=S_{BA}$, $S^A{}_A=0$, and
$\Omega_{AB}=-\Omega_{BA}$, Paper~I-D,
Proposition~\ref{p1b:prop:local-coframe}, applied after this first-symbol
descent, gives
\[
 \boxed{\delta g=2\lambda g+2S,\qquad \delta_\Omega g=0.}
\]
Thus the trace and symmetric-traceless pieces of the quadratic parent become
homothety and metric shear only after this local causal/coframe descent; the
antisymmetric infinitesimal is a frame-rotation direction at this stage.
Raw compact-parent $S^2_0V$ content is therefore not identified with physical
Lorentzian gravity before the present typed map.

\begin{corollary}[Spatial shear decomposition]
\label{p3a:cor:v189-spatial-shear-decomposition}
For the generated $3+1$ split $V=\mathbb Re^0\oplus W$,
\[
 \boxed{S^2_0(V)\downarrow SO(3)
 =S^2_0(W)\oplus W\oplus\mathbb R,}
 \qquad 9=5+3+1.
\]
The five-dimensional $S^2_0(W)$ component is the kinematic spatial shear
role; lapse/shift/constraint reduction remains the job of the same-action ADM
and Noether descendants.
The representation decomposition is Paper~I-D,
Proposition~\ref{p1b:prop:local-coframe}, for this generated coframe.
\end{corollary}


\subsection{Complete visible response as a dynamical comparison quotient}
\label{p3a:sec:v227-visible-dynamics}
The source-visible quotient can close a dynamical readout before selection
of an unused complete environment. For two licensed finite Hermitian
parents, Theorem~\ref{p1b:thm:v227-finite-response} gives a finite exact
certificate of equality of the whole retained response. Its source-jet
corollary supplies all differentiated words actually consumed at the fixed
source point. The existing minimal visible realization then identifies the
coupled environment up to unitary equivalence and leaves only reducing,
uncoupled dark data unconstrained.

This is a claim-relative reduction of the comparison obligation. The same
source ports, energy convention, density and readout must be retained.
A complete action/heat trace, another source, or an energy allocation may
see the previously dark data. In particular neutral Liouvillian dynamics
can ignore actual block-central offsets on the branch of
Theorem~\ref{p1b:thm:v227-central-readouts}, while odd source frequencies
and absolute action remain subject to its boundary proposition.




\section{Petersen support and the common-action incidence boundary}
\label{p3a:sec:v319-petersen-action}
This section imports only the \emph{derived} finite packet of Paper~I-B,
Theorems~\ref{p1:thm:petersen-fifteen-resolution} and
\ref{p1:thm:petersen-mirror-envelope}, on an already occurring, embedded
$\mathrm{Supp}_{2I}$ branch.  The resulting ten Hopf axes and fifteen
incidence channels are not the unrelated fifteen-dimensional momentum
first-level carrier $\mathcal H_1$.  No unique physical $2I$ marking is
created on an unmarked $S^3$ or $M_8$ by this import.

\begin{proposition}[Five-core cancellation and marked action jets]
\label{p3a:prop:v319-petersen-marked-jet}
Let $J_m:\mathbf3\to\mathbf3'$ be the five real orthogonal
core intertwiners of Paper~I-B.  Then
\[
  \sum_{m=1}^5J_m=0,\qquad
  J_{S^c}=-J_S,\quad J_S=\sum_{m\in S}J_m.
\]
Consequently every linear functional $\ell$ on the real nine-dimensional
space $\operatorname{Hom}(\mathbf3,\mathbf3')$ obeys
$\sum_m\ell(J_m)=0$.  For a marked two-core subset the quadratic
polar source $J_S^{\mathsf T}J_S$ is complement-even, whereas any
linear contribution in $J_S$ is complement-odd.  Neither statement
implies a nonzero term in the physical action.
\end{proposition}
\begin{proof}
The identities are the exact finite-resolution and simplex identities
of Paper~I-B, Theorems~\ref{p1:thm:petersen-bisector-simplex} and
\ref{p1:thm:petersen-subset-polar}.  Apply linearity and
$( -J_S)^{\mathsf T}(-J_S)=J_S^{\mathsf T}J_S$.
\end{proof}

\begin{proposition}[What the inherited geometry does not generate]
\label{p3a:prop:v319-petersen-action-firewall}
On the unmarked full five-core packet, the $A_5$ action is a finite
rotation action, not an independently generated continuous Noether
symmetry.  The family of local intertwiners and cycle operators is
mathematically available but supplies neither a physical field space
nor a new action coefficient.  If a licensed common action actually
contains a Petersen operator, its complete jets and source markings
must be co-consumed by
Theorem~\ref{p3a:thm:v2216-common-action-consumption}; conversely that
theorem does not insert an otherwise absent operator.
\end{proposition}

\section{Results, limits, and downstream interfaces}
\label{paper3a:sec:organized-conclusion}
\subsection{Conclusion}
\label{p3a:sec:conclusion}
Paper~III-A imports the current Paper-II parent-role constitution and the
already derived physical momentum seed
\(\mathcal H_p\cong\mathcal H_1,\qquad m_p=1\), G0 actualization semantics,
and Paper~II-C's premise-minimal intrinsic-evaluation reciprocal closure.  The
Born/GP2-generated reciprocal role, unquotiented relation-free product, and Born
exchange involution therefore require no metric-curry premise.  Under the
explicit optional G0 certificate $\mathsf P_{\rm mc}$, Paper~III-A promotes
the full canonical profile family to a physical reciprocal test doctrine with
the variance-opposite pullback action and re-proves the BR/IND statuses on
that stronger domain.  The associated function algebras and spectra remain
mathematical envelopes; the scalar physical B-role is supplied instead by the
typed loop-character evaluation constructor.  The paper
constructs the continuous vector reciprocal product $S^3_x\times S^3_p$ and
the scalar Born torus, deriving the primitive Poincar\'e phase class of the
latter.  Paper~I-A's positive primitive Reeb lattice and Paper~II-A's actual
Hopf-to-succession provenance occurrence fix the scalar succession sign; the
unique full-dual primitive fixes the reciprocal energy sign.  No later free
orientation interface remains, although units, Hamiltonian selection, actual
evolution, and any thermodynamic arrow remain open.

The vector product is not globally symplectic, while the exact finite
full-product-invariant commutant is abelian and contains no Heisenberg pair.
These results show that the required scalar canonical pair is not hidden in the
vector product; the separately prior-generated scalar role supplies the scalar
input to the Born completion.  The complete compact
$S^3\times S^3\times T^2$ geometry is likewise not globally symplectic, so
only local/universal-cover Darboux variables and the scalar torus phase class
are used.

After the elapsed scalar cover is constituted and the explicit calibration
\((\mathrm{CAL}_{t})\) assigns one primitive turn the duration \(T_t\), the
retained Hopf provenance, typed parent
inclusion, and premise-minimal reciprocal action derive the typed least
physical clocked operation doctrine.  The closed convex envelope of its
velocity germs supplies the future cone without promoting every relaxed
history to a physical token.  Then
Theorem~\ref{p3:thm:gp1-operation-generator-polar-dual} derives its exact
local generator polar
\[
   E\ge c_{\rm op}\|p\|_{h^{-1}},
   \qquad ((\mathcal C_x^+)^\vee)^\vee=\mathcal C_x^+.
\]
This crossed local operation/generator duality is derived without identifying
the global compact reciprocal carrier with a cotangent bundle.  Its
identity-connected \(O(3)\)-equivariant cone automorphisms are exactly
\(H^{1/4}I_4\) at determinant \(H>0\).  Thus the rank-four root operation is
canonical once a determinant character is supplied, while no response owner
is identified.  Combining the
II-C vector first-symbol dual with the scalar succession/energy dual gives the
canonical generator-germ co-realization of
Theorem~\ref{p3a:thm:canonical-generator-germ-localization}; no separate G0
\(\mathrm{GL}_{xp}\) localization witness remains.  The scalar-harmonic bridge and the finite unit-only map are already
generated upstream. Stronger observable-to-operator and continuum/local--global
comparisons and dimensionful calibrations retain their stated domains.

Theorem~\ref{p3a:thm:complete-joint-zero-defect-locus} proves simultaneous,
not coordinatewise, realizability across the reciprocal products, scalar
orientation and phase class, finite stages, clock winding/doctrine, and active
calibrations.  It retains an unresolved position-stage family until a frontier
occurs, but removes graph supports, wrong provenance signs, higher windings,
and incompatible calibration tuples by their explicit joint defects.  A later
neutral field or photon speed is likewise not a free coordinate: after its
undeformed same-owner principal ancestry is proved, Paper~II-A forces its
front speed to equal \(c_{\rm op}\); support ancestry alone is insufficient.

Theorems~\ref{p3a:thm:universal-generated-exact-continuation} and
\ref{p3a:thm:finite-generator-one-ray-rigidity} supply the framework-wide
dynamics architecture without adding a universal physical generator.
Each sector supplies its own frozen constructor certificate for the complete
family image and licensed quotient; a comparison specification alone is only
mathematical classification.  The
identity-survival corollary blocks motion from a static defect alone; a
generated oriented response comparison is required to exclude the stationary
member.  The law--flow--history theorem now makes the chronology explicit: a
generated law or infinitesimal generator is not yet a certified process flow,
and a certified flow is not yet an actual history token.  Native same-owner
action automorphisms inherit action invariance from Paper~II-A's GP2 theorem,
so their Noether identities and constraints enter \(\Delta_D\) by internal
theorem closure; analytic existence, uniqueness, hyperbolicity, compactness,
or continuum-generator hypotheses remain separately certified.  A fixed
finite-dimensional homogeneous linear branch has an exponential generator,
and the unitary specialization gives \(U(t)=e^{-itG}\) with \(G=G^*\)
without a separate Stone-type physical premise.  A commutant calculation can reduce the exact generator set to one positive
direction. The surviving rate set still obeys every endpoint and strict
comparison, and its actual isomorphism/calibration quotient must be a
singleton for one-ray uniqueness.
Independent process roles remain products, while shared laws or rates require
a generated common owner relation.  Corollary~\ref{p3a:cor:common-process-invariant}
nevertheless retains any licensed invariant common to every unresolved
process member, and Theorem~\ref{p3a:thm:generated-process-projectability}
gives the exact set intertwining test; descent with smooth or other arrow
types additionally uses their category-specific factorization certificate.
Theorem~\ref{p3a:thm:threshold-event-descent} applies that test to ordered
threshold readouts.  It derives the discrete successor law from the parent
history while proving that a strictly crossing integer staircase normally
loses the physical-time phase needed for a homogeneous semigroup.  Retaining
the generated distance to the next threshold reconstructs the closed parent
scalar state and gives an exact set/event description, with cellwise resets
rather than automatic smoothness of the discrete-product coordinates.

Finite reciprocity is asymmetric at this stage: the momentum entry is already
fixed to $1$, while the generic position-side comparison family
$\mathsf K_{\rm res}(m,1)$ is theorem-level and an occurring frontier derives $m_x=m_x^{\rm res}(F_x)$ and therefore
determines the stagewise member $\mathsf K_{\rm res}(m_x,1)$.  Its primitive frontier
already has $m_x=1$.  Later action, state, regularity,
reconstruction, polarization, and dynamical choices retain their own displayed
domains.  Occurrence, preparation, process, trajectory, test, bridge, and
calibration domains share the G0 actualization rule, while analytic regularity
and reconstruction assumptions remain distinct.  The auxiliary momentum-selector countermodels are collected in
Paper~III-G, Appendix~\ref{p3a:app:legacy-selector}.

\printbibliography
\end{document}
